\documentclass[11pt]{ctexart}
\usepackage[a4paper,margin=1in]{geometry}
\usepackage{graphicx}
\usepackage{xcolor}
\usepackage{hyperref}
\definecolor{refgray}{gray}{0.45}
\title{\textbf{市场最新观点汇总}\\[0.4em]{\large 市场主线：AI硬件与供给瓶颈主导定价，中国内需磨底与结构性迁移并行}}
\author{KC桌面 自动生成}
\date{260622}
\begin{document}
\maketitle
\tableofcontents
\newpage
\section{一页摘要}
\begin{itemize}
  \item AI算力投资正从'军备竞赛'转向'资本效率竞赛'，云巨头资本开支增速持续超过现金流，外部融资需求上升。
  \item 内存（DRAM/HBM）供给侧定价权转移，HBM价格预计跳升2-2.5倍，可能迫使云厂商资本开支上调30\%。
  \item 中国A股硬科技盈利增速（CSI300 +9\%）显著领先离岸市场（MSCI中国 -8\%），AI利润池向硬件回摆。
  \item 中国家庭资产正从房产向金融资产结构性迁移，房产占比52\%且持续收缩，股权和保险占比将扩大。
  \item 欧洲化工面临冲突结束后暴露的结构性过剩，中国产能外溢与配方重构双重冲击。
  \item 日本零售业接近加速整合拐点，人力成本压力与消费行为分裂正在淘汰中间业态。
  \item 全球游戏业2026年增速放缓至0.7\%，行业进入整合驱动阶段，GTA VI将成为情绪转折点。
  \item 代币化RWA市值突破510亿美元，竞争焦点从资产发行转向交易基础设施的合规与规模化。
  \item 数据中心供电瓶颈不在技术，而在电网并网等待时间（4年+），运营商首选电网连接但容忍度极低。
  \item 人民币中间价模型显示逆周期因子正在发挥比表面数字更重要的校准作用，市场对单向波动的定价可能过于激进。
\end{itemize}
\section{宏观与利率：中国内需磨底，家庭资产负债表再平衡}
\textbf{中国出口强劲与内需疲软的裂口已固化，家庭部门主动去杠杆是理解未来资产价格走向的主线。}

\begin{itemize}
  \item 五月零售销售同比下降0.6\%，为除疫情封控期外首次负增长，固定资产投资同步收缩，而出口同比增长19.4\%。
  \item GS指出，消费收缩并非单纯收入下降，而是家庭部门在房价持续下跌背景下主动减少消费、增加储蓄以修复资产负债表。
  \item 中国家庭资产中房产占比高达52\%，现金/存款占25\%，GS预计未来几年股权和保险占比将显著扩大，房产占比持续收缩。
  \item 央行在陆家嘴论坛宣布将隔夜回购利率走廊从70bp缩窄至50bp，进一步向价格型货币政策框架靠拢。
  \item 央行同时推出境外机构人民币回购工具（FIMA）、香港5年期国债期货等，人民币国际化明显提速。
  \item NOM模型预测美元兑人民币中间价为6.7787，较前次下移343个基点，但计入逆周期因子后为6.8008，与官方收盘价偏离仅122个基点，显示政策微调意图。
  \item NOM报告提示，2026年下半年关键时间节点包括7月底政治局会议、10月国庆、11月APEC峰会及12月中央经济工作会议，这些事件可能影响汇率预期。
\end{itemize}
\begin{figure}[htbp]
\centering
\includegraphics[width=0.88\linewidth]{figures/fig_001.jpg}
\caption{Exhibit 1 - Exhibit 1: According to Our Nowcast, Global EV Car Sales Penetration Has Increased By 3.4pp Since the Beginning of the Hormuz Shock}
\end{figure}
\begin{figure}[htbp]
\centering
\includegraphics[width=0.88\linewidth]{figures/fig_002.jpg}
\caption{Exhibit 1 - Exhibit 3: The Increase in Global EV Car Sales Penetration Since February Has Been Driven 61\% By China, 21\% By OECD, and 19\% By Non OECD ex China}
\end{figure}
\begin{figure}[htbp]
\centering
\includegraphics[width=0.88\linewidth]{figures/fig_003.jpg}
\caption{Exhibit 1 - Exhibit 3: The Increase in Global EV Car Sales Penetration Since February Has Been Driven 61\% By China, 21\% By OECD, and 19\% By Non OECD ex China Exhibit 4: 12 Out of the 15 Largest EV Markets Have Experienced an Increase in EV C}
\end{figure}
\begin{figure}[htbp]
\centering
\includegraphics[width=0.88\linewidth]{figures/fig_004.jpg}
\caption{Exhibit 2 - Exhibit 4: 12 Out of the 15 Largest EV Markets Have Experienced an Increase in EV Car Sales Penetration Since February For 15 largest EV markets by 2025 sales volume. Solid bars indicate that May share is calculated using realize}
\end{figure}
\begin{figure}[htbp]
\centering
\includegraphics[width=0.88\linewidth]{figures/fig_005.jpg}
\caption{Exhibit 3 - Exhibit 6: EV Car Sales Have Accelerated Meaningfully Since February Across the World ex US}
\end{figure}
\begin{figure}[htbp]
\centering
\includegraphics[width=0.88\linewidth]{figures/fig_006.jpg}
\caption{Exhibit 6 - Exhibit 6: EV Car Sales Have Accelerated Meaningfully Since February Across the World ex US Red diamonds indicate data points that are nowcasted. The vertical line on each chart indicates the start of the Hormuz shock. Exhibit}
\end{figure}
\begin{figure}[htbp]
\centering
\includegraphics[width=0.88\linewidth]{figures/fig_007.jpg}
\caption{Exhibit 6 - Exhibit 6: EV Car Sales Have Accelerated Meaningfully Since February Across the World ex US Red diamonds indicate data points that are nowcasted. The vertical line on each chart indicates the start of the Hormuz shock. Exhibit}
\end{figure}
{\small\color{refgray}\textbf{References:} [R001] GS：市场低估了霍尔木兹冲击对油需的结构性重塑; [R015] GS：中国家庭资产正在经历从房产到金融资产的结构性迁移; [R018] NOM：人民币中间价模型正在传递一个比市场预期更克制的信号}

\section{AI/算力：云巨头资本效率竞赛与内存成本冲击}
\textbf{云市场正从'建设能力'竞争转向'运营效率'与'资本结构'的复合博弈，内存成本压力可能迫使AI投资回报模型重估。}

\begin{itemize}
  \item Bernstein报告显示，四大云巨头（AWS、Azure、GCP、OCI）2026年Q1现金资本支出合计约1300亿美元，全年预计达6230亿美元（含Meta约7600亿），但运营现金流增速明显滞后。
  \item 资本支出增速持续超过运营现金流增长，亚马逊、谷歌和甲骨文已陆续进入资本市场融资，外部融资能力成为新护城河。
  \item 谷歌云成为最大黑马，Q1利润率从17.8\%飙升至约33\%，积压订单接近翻倍至约4620亿美元，AI云收入增量已与微软不相上下。
  \item 微软AI ARR达370亿美元，Office 365 Copilot和GitHub AI功能使用量明显增长，但整体收入增速仍需加速以匹配资本开支节奏。
  \item Bernstein报告揭示，常规DRAM价格已上涨约4.5倍，而HBM因年度合同定价几乎未动，同样晶圆产能做常规DRAM收入是HBM的2倍以上。
  \item HBM价格预计跳升2-2.5倍，经GPU供应商加价传导后，云厂商数据中心资本开支可能被迫上调约30\%。
  \item 内存供给侧定价权转移正在扭曲产能分配，三星、美光等供应商有动机将产能从HBM转向常规DRAM，加剧AI算力紧缺。
  \item GS提出软件行业'好粘性'与'坏粘性'之分：AI正在系统性地瓦解依赖历史迁移成本的'坏粘性'，依赖持续创新的'好粘性'公司护城河更深。
\end{itemize}
\begin{figure}[htbp]
\centering
\includegraphics[width=0.88\linewidth]{figures/fig_047.jpg}
\caption{EXHIBIT 1 - EXHIBIT 2: While cloud revenue growth has accelerated in the recent quarters, it needs to continue accelerating to keep up with the pace of capex growth Quarterly Cloud Revenue: AWS+Azure+GCP+OCI Based on calendar quarters (Calen}
\end{figure}
\begin{figure}[htbp]
\centering
\includegraphics[width=0.88\linewidth]{figures/fig_048.jpg}
\caption{EXHIBIT 2 - EXHIBIT 3: On the other hand, ignoring contract duration, total backlog (RPO) growth has accelerated significantly across all hypersclaers in recent quarters, totaling \textbackslash{}\$2T as of 1Q26}
\end{figure}
\begin{figure}[htbp]
\centering
\includegraphics[width=0.88\linewidth]{figures/fig_049.jpg}
\caption{EXHIBIT 3 - EXHIBIT 5: Meanwhile, as hyperscalers are running out of internal funding, to increase Capex further they likely need to (and have been) raise capital, through either debt or equity}
\end{figure}
\begin{figure}[htbp]
\centering
\includegraphics[width=0.88\linewidth]{figures/fig_050.jpg}
\caption{EXHIBIT 3 - EXHIBIT 5: Meanwhile, as hyperscalers are running out of internal funding, to increase Capex further they likely need to (and have been) raise capital, through either debt or equity Hyperscaler Cash Flow vs. Capex, CY 26E Based o}
\end{figure}
\begin{figure}[htbp]
\centering
\includegraphics[width=0.88\linewidth]{figures/fig_051.jpg}
\caption{EXHIBIT 4 - EXHIBIT 4: So, which tale would you rather believe? Y/Y Growth\%: Revenue vs. Capex vs. CFO vs. Backlog EXHIBIT 5: Meanwhile, as hyperscalers are running out of internal funding, to increase Capex further they likely need to (and}
\end{figure}
\begin{figure}[htbp]
\centering
\includegraphics[width=0.88\linewidth]{figures/fig_052.jpg}
\caption{EXHIBIT 6 - EXHIBIT 6: The five hyperscale providers report Cloud revenue, which includes not only IaaS/PaaS but also other revenues Google includes Workspace. Microsoft Azure revenue is based on our estimates. EXHIBIT 7: In Q1 CY26, Microso}
\end{figure}
\begin{figure}[htbp]
\centering
\includegraphics[width=0.88\linewidth]{figures/fig_053.jpg}
\caption{EXHIBIT 6 - EXHIBIT 8: On a Q/Q basis, AWS added \textbackslash{}\$2.0B vs Azure at \textbackslash{}\$2.5B vs Google Cloud at \textbackslash{}\$2.4B vs Oracle at \textbackslash{}\$0.9B EXHIBIT 9: Of the incremental dollars added Q/Q by the hyperscale Cloud providers (excluding BABA), Google and Oracle ha}
\end{figure}
\begin{figure}[htbp]
\centering
\includegraphics[width=0.88\linewidth]{figures/fig_054.jpg}
\caption{EXHIBIT 7 - EXHIBIT 10: Q over Q, incremental share gained by each hyperscaler (excluding Alibaba)}
\end{figure}
\begin{figure}[htbp]
\centering
\includegraphics[width=0.88\linewidth]{figures/fig_055.jpg}
\caption{EXHIBIT 8 - EXHIBIT 10: Q over Q, incremental share gained by each hyperscaler (excluding Alibaba) All quarters represent calendar quarters, Google's cloud revenues include both GCP and Workspace. Microsoft Azure numbers are based on estimate}
\end{figure}
\begin{figure}[htbp]
\centering
\includegraphics[width=0.88\linewidth]{figures/fig_056.jpg}
\caption{EXHIBIT 9 - EXHIBIT 9: Of the incremental dollars added Q/Q by the hyperscale Cloud providers (excluding BABA), Google and Oracle have gained incremental shares this quarter. Microsoft Azure numbers are based on estimates EXHIBIT 10: Q over}
\end{figure}
\begin{figure}[htbp]
\centering
\includegraphics[width=0.88\linewidth]{figures/fig_057.jpg}
\caption{EXHIBIT 11 - EXHIBIT 11: AWS Backlog grew to \textbackslash{}\$364B up \textbackslash{}\textasciitilde{}50\% from the \textbackslash{}\$244B announced in 4Q25 EXHIBIT 12: AWS average contract life rose Q/Q to 5.5 years in 1Q26 AWS Weighted average contract life (Yrs.) \#\# Microsoft Azure}
\end{figure}
\begin{figure}[htbp]
\centering
\includegraphics[width=0.88\linewidth]{figures/fig_058.jpg}
\caption{EXHIBIT 11 - EXHIBIT 11: AWS Backlog grew to \textbackslash{}\$364B up \textbackslash{}\textasciitilde{}50\% from the \textbackslash{}\$244B announced in 4Q25 EXHIBIT 12: AWS average contract life rose Q/Q to 5.5 years in 1Q26 AWS Weighted average contract life (Yrs.) \#\# Microsoft Azure \textbackslash{}- Microsoft Azur}
\end{figure}
\begin{figure}[htbp]
\centering
\includegraphics[width=0.88\linewidth]{figures/fig_059.jpg}
\caption{EXHIBIT 13 - EXHIBIT 13: Microsoft Azure constant currency growth Azure Revenue growth (CC) EXHIBIT 14: CAPEX has gone up as Microsoft continues to invest in AI to meet demand Microsoft Total Capex (Cash + Financial Lease), in \textbackslash{}\$B}
\end{figure}
\begin{figure}[htbp]
\centering
\includegraphics[width=0.88\linewidth]{figures/fig_060.jpg}
\caption{EXHIBIT 14 - EXHIBIT 14: CAPEX has gone up as Microsoft continues to invest in AI to meet demand Microsoft Total Capex (Cash + Financial Lease), in \textbackslash{}\$B \#\# Google Cloud \textbackslash{}- Google Cloud is gaining solid ground, with revenue surging 63\% Y/Y to ex}
\end{figure}
\begin{figure}[htbp]
\centering
\includegraphics[width=0.88\linewidth]{figures/fig_061.jpg}
\caption{EXHIBIT 15 - EXHIBIT 15: Google cloud's backlog of \textbackslash{}\$460B nearly doubled Q/Q EXHIBIT 16: Revenue/backlog ratio significantly increased sequentially for Google Cloud \#\# Alibaba}
\end{figure}
\begin{figure}[htbp]
\centering
\includegraphics[width=0.88\linewidth]{figures/fig_062.jpg}
\caption{EXHIBIT 15 - EXHIBIT 15: Google cloud's backlog of \textbackslash{}\$460B nearly doubled Q/Q EXHIBIT 16: Revenue/backlog ratio significantly increased sequentially for Google Cloud \#\# Alibaba \textbackslash{}- Alibaba's cloud revenue is reported on a gross basis inclusive}
\end{figure}
\begin{figure}[htbp]
\centering
\includegraphics[width=0.88\linewidth]{figures/fig_063.jpg}
\caption{EXHIBIT 17 - EXHIBIT 17: OCI revenue guidance through FY30 (CY29-30) EXHIBIT 18: Oracle RPO and CRPO OCI 5Y Revenue Guidance (FY26-30) \#\# WHAT ABOUT MARGINS? AI, we believe, is going to be a drag on gross margins at least during the build-ou}
\end{figure}
\begin{figure}[htbp]
\centering
\includegraphics[width=0.88\linewidth]{figures/fig_064.jpg}
\caption{EXHIBIT 17 - EXHIBIT 17: OCI revenue guidance through FY30 (CY29-30) EXHIBIT 18: Oracle RPO and CRPO OCI 5Y Revenue Guidance (FY26-30) \#\# WHAT ABOUT MARGINS? AI, we believe, is going to be a drag on gross margins at least during the build-ou}
\end{figure}
\begin{figure}[htbp]
\centering
\includegraphics[width=0.88\linewidth]{figures/fig_065.jpg}
\caption{EXHIBIT 20 - EXHIBIT 20: Software Moat Overview EXHIBIT 21: GPU Cloud Landscape (Scale v. Sophistication) Sophistication / (SemiAnalysis ClusterMAX 2.0) EXHIBIT 22: CRWV Revenue Growth CRWV Revenue Growth}
\end{figure}
\begin{figure}[htbp]
\centering
\includegraphics[width=0.88\linewidth]{figures/fig_066.jpg}
\caption{EXHIBIT 22 - EXHIBIT 22: CRWV Revenue Growth CRWV Revenue Growth EXHIBIT 23: CRWV Debt Overview (Times, \%) CRWV Debt Overview Leverage calculated as LTM Adj. EBITDA over net debt}
\end{figure}
\begin{figure}[htbp]
\centering
\includegraphics[width=0.88\linewidth]{figures/fig_067.jpg}
\caption{EXHIBIT 23 - EXHIBIT 23: CRWV Debt Overview (Times, \%) CRWV Debt Overview Leverage calculated as LTM Adj. EBITDA over net debt EXHIBIT 24: CRWV CAPEX Expense (\textbackslash{}\$B, \%) CRWV CAPEX Expense EXHIBIT 25: CRWV Capital Intensity (Times) CRWV Capital}
\end{figure}
\begin{figure}[htbp]
\centering
\includegraphics[width=0.88\linewidth]{figures/fig_068.jpg}
\caption{EXHIBIT 24 - EXHIBIT 24: CRWV CAPEX Expense (\textbackslash{}\$B, \%) CRWV CAPEX Expense EXHIBIT 25: CRWV Capital Intensity (Times) CRWV Capital Intensity \textbackslash{}- CoreWeave's margin profile is currently depressed by a structural OpEx–revenue lag. The company begi}
\end{figure}
\begin{figure}[htbp]
\centering
\includegraphics[width=0.88\linewidth]{figures/fig_069.jpg}
\caption{EXHIBIT 24 - EXHIBIT 24: CRWV CAPEX Expense (\textbackslash{}\$B, \%) CRWV CAPEX Expense EXHIBIT 25: CRWV Capital Intensity (Times) CRWV Capital Intensity \textbackslash{}- CoreWeave's margin profile is currently depressed by a structural OpEx–revenue lag. The company begi}
\end{figure}
\begin{figure}[htbp]
\centering
\includegraphics[width=0.88\linewidth]{figures/fig_070.jpg}
\caption{EXHIBIT 26 - EXHIBIT 26: CRWV Adj. Operating Income (\textbackslash{}\$M) CRWV Adj. Operating Income Adj. Operating Income — Adj. Operating Income Margin EXHIBIT 27: CRWV Adj. EBITDA (\textbackslash{}\$M) CRWV Adj. EBITDA}
\end{figure}
\begin{figure}[htbp]
\centering
\includegraphics[width=0.88\linewidth]{figures/fig_071.jpg}
\caption{EXHIBIT 27 - EXHIBIT 27: CRWV Adj. EBITDA (\textbackslash{}\$M) CRWV Adj. EBITDA EXHIBIT 28: CRWV Diluted EPS (\textbackslash{}\$) CRWV Diluted EPS \#\# I. REQUIRED DISCLOSURES}
\end{figure}
\begin{figure}[htbp]
\centering
\includegraphics[width=0.88\linewidth]{figures/fig_072.jpg}
\caption{EXHIBIT 27 - EXHIBIT 27: CRWV Adj. EBITDA (\textbackslash{}\$M) CRWV Adj. EBITDA EXHIBIT 28: CRWV Diluted EPS (\textbackslash{}\$) CRWV Diluted EPS \#\# I. REQUIRED DISCLOSURES Bernstein is part of a joint venture between SG (SG) and AllianceBernstein, L.P. (AB). Unless spec}
\end{figure}
\begin{figure}[htbp]
\centering
\includegraphics[width=0.88\linewidth]{figures/fig_104.jpg}
\caption{Exhibit 3 - EXHIBIT 1: We forecast 2-2.5x increase in HBM prices YoY next year so that the profitability of HBM \& conventional DRAM can narrow. EXHIBIT 2: We expect the price increase to take place at all generations of HBM. EXHIBIT 3: We}
\end{figure}
\begin{figure}[htbp]
\centering
\includegraphics[width=0.88\linewidth]{figures/fig_105.jpg}
\caption{EXHIBIT 1 - EXHIBIT 1: We forecast 2-2.5x increase in HBM prices YoY next year so that the profitability of HBM \& conventional DRAM can narrow. EXHIBIT 2: We expect the price increase to take place at all generations of HBM. EXHIBIT 3: We}
\end{figure}
\begin{figure}[htbp]
\centering
\includegraphics[width=0.88\linewidth]{figures/fig_106.jpg}
\caption{EXHIBIT 2 - Exhibit 1-Exhibit 2). HBM profitability hence will remain below that of conventional DRAM next year, but the gap should be much smaller than this year.}
\end{figure}
\begin{figure}[htbp]
\centering
\includegraphics[width=0.88\linewidth]{figures/fig_107.jpg}
\caption{Exhibit 4 - EXHIBIT 6: Conventional DRAM now generates considerably more revenue per wafer capacity than HBM, and we project an HBM price hike next year to narrow that gap.}
\end{figure}
\begin{figure}[htbp]
\centering
\includegraphics[width=0.88\linewidth]{figures/fig_108.jpg}
\caption{Exhibit 8 - EXHIBIT 6: Conventional DRAM now generates considerably more revenue per wafer capacity than HBM, and we project an HBM price hike next year to narrow that gap. EXHIBIT 7: We believe the margin gap between HBM and conventional DR}
\end{figure}
\begin{figure}[htbp]
\centering
\includegraphics[width=0.88\linewidth]{figures/fig_109.jpg}
\caption{EXHIBIT 5 - EXHIBIT 7: We believe the margin gap between HBM and conventional DRAM to narrow next year. S. Chungcheong Export N. Chungcheong Export + Icheon Export Other Provinces}
\end{figure}
\begin{figure}[htbp]
\centering
\includegraphics[width=0.88\linewidth]{figures/fig_110.jpg}
\caption{EXHIBIT 6 - EXHIBIT 8: Value per weight remained largely in the same range though the export value per weight for Samsung rose notably \& suggested HBM4 shipment in May. Korea Multichip Memory Export Value per Weight to TW+MY}
\end{figure}
\begin{figure}[htbp]
\centering
\includegraphics[width=0.88\linewidth]{figures/fig_111.jpg}
\caption{EXHIBIT 7 - EXHIBIT 10: With the expected HBM price increase, we forecast HBM's revenue mix in DRAM to rebound next year. HBM Contribution to DRAM Revenue}
\end{figure}
\begin{figure}[htbp]
\centering
\includegraphics[width=0.88\linewidth]{figures/fig_112.jpg}
\caption{EXHIBIT 8 - EXHIBIT 10: With the expected HBM price increase, we forecast HBM's revenue mix in DRAM to rebound next year. HBM Contribution to DRAM Revenue \#\# DRAM stocks will benefit from a rapid upward earnings revision in the near term.}
\end{figure}
\begin{figure}[htbp]
\centering
\includegraphics[width=0.88\linewidth]{figures/fig_113.jpg}
\caption{EXHIBIT 9 - EXHIBIT 10: With the expected HBM price increase, we forecast HBM's revenue mix in DRAM to rebound next year. HBM Contribution to DRAM Revenue \#\# DRAM stocks will benefit from a rapid upward earnings revision in the near term. con}
\end{figure}
\begin{figure}[htbp]
\centering
\includegraphics[width=0.88\linewidth]{figures/fig_114.jpg}
\caption{EXHIBIT 13 - EXHIBIT 13: Our Micron FY27 EPS forecast is 38\% above consensus too. EXHIBIT 14: We expect EPS to grow exponentially going forward. EXHIBIT 15: We model earnings to reach peak in 2HCY27... EXHIBIT 16: ... and then gradually decl}
\end{figure}
\begin{figure}[htbp]
\centering
\includegraphics[width=0.88\linewidth]{figures/fig_115.jpg}
\caption{EXHIBIT 14 - EXHIBIT 14: We expect EPS to grow exponentially going forward. EXHIBIT 15: We model earnings to reach peak in 2HCY27... EXHIBIT 16: ... and then gradually decline in CY28 as prices normalize. EXHIBIT 17: We forecast DRAM price}
\end{figure}
\begin{figure}[htbp]
\centering
\includegraphics[width=0.88\linewidth]{figures/fig_116.jpg}
\caption{EXHIBIT 15 - EXHIBIT 18: ROE will rise to unprecedented levels.}
\end{figure}
\begin{figure}[htbp]
\centering
\includegraphics[width=0.88\linewidth]{figures/fig_117.jpg}
\caption{EXHIBIT 16 - EXHIBIT 18: ROE will rise to unprecedented levels. EXHIBIT 19: Samsung BVPS can grow by nearly 2x in both 2026 and 2027.}
\end{figure}
\begin{figure}[htbp]
\centering
\includegraphics[width=0.88\linewidth]{figures/fig_118.jpg}
\caption{EXHIBIT 17 - EXHIBIT 20: SK hynix can see even steeper accretion in BVPS.}
\end{figure}
\begin{figure}[htbp]
\centering
\includegraphics[width=0.88\linewidth]{figures/fig_119.jpg}
\caption{EXHIBIT 18 - EXHIBIT 18: ROE will rise to unprecedented levels. EXHIBIT 19: Samsung BVPS can grow by nearly 2x in both 2026 and 2027. EXHIBIT 20: SK hynix can see even steeper accretion in BVPS. EXHIBIT 21: We see 4x increase in Micron BVP}
\end{figure}
\begin{figure}[htbp]
\centering
\includegraphics[width=0.88\linewidth]{figures/fig_120.jpg}
\caption{EXHIBIT 19 - EXHIBIT 19: Samsung BVPS can grow by nearly 2x in both 2026 and 2027. EXHIBIT 20: SK hynix can see even steeper accretion in BVPS. EXHIBIT 21: We see 4x increase in Micron BVPS from now to end of FY27. EXHIBIT 22: Cash will ac}
\end{figure}
\begin{figure}[htbp]
\centering
\includegraphics[width=0.88\linewidth]{figures/fig_121.jpg}
\caption{EXHIBIT 20 - EXHIBIT 20: SK hynix can see even steeper accretion in BVPS. EXHIBIT 21: We see 4x increase in Micron BVPS from now to end of FY27. EXHIBIT 22: Cash will accumulate at unprecedented pace as well. EXHIBIT 23: We forecast cash t}
\end{figure}
\begin{figure}[htbp]
\centering
\includegraphics[width=0.88\linewidth]{figures/fig_122.jpg}
\caption{EXHIBIT 21 - EXHIBIT 21: We see 4x increase in Micron BVPS from now to end of FY27. EXHIBIT 22: Cash will accumulate at unprecedented pace as well. EXHIBIT 23: We forecast cash to reach 70-80\% of book value... EXHIBIT 24: ... unless compan}
\end{figure}
\begin{figure}[htbp]
\centering
\includegraphics[width=0.88\linewidth]{figures/fig_123.jpg}
\caption{EXHIBIT 22 - EXHIBIT 22: Cash will accumulate at unprecedented pace as well. EXHIBIT 23: We forecast cash to reach 70-80\% of book value... EXHIBIT 24: ... unless company meaningful increase payout or use it for M\&A. EXHIBIT 25: We now targ}
\end{figure}
\begin{figure}[htbp]
\centering
\includegraphics[width=0.88\linewidth]{figures/fig_124.jpg}
\caption{EXHIBIT 23 - EXHIBIT 23: We forecast cash to reach 70-80\% of book value... EXHIBIT 24: ... unless company meaningful increase payout or use it for M\&A. EXHIBIT 25: We now target 6.2x 1 year P/E to value Samsung. EXHIBIT 26: We also use 6.2}
\end{figure}
\begin{figure}[htbp]
\centering
\includegraphics[width=0.88\linewidth]{figures/fig_125.jpg}
\caption{EXHIBIT 24 - EXHIBIT 24: ... unless company meaningful increase payout or use it for M\&A. EXHIBIT 25: We now target 6.2x 1 year P/E to value Samsung. EXHIBIT 26: We also use 6.2x P/E for hynix. EXHIBIT 27: Micron target P/E is 7.7x, still}
\end{figure}
\begin{figure}[htbp]
\centering
\includegraphics[width=0.88\linewidth]{figures/fig_126.jpg}
\caption{EXHIBIT 25 - EXHIBIT 25: We now target 6.2x 1 year P/E to value Samsung. EXHIBIT 26: We also use 6.2x P/E for hynix. EXHIBIT 27: Micron target P/E is 7.7x, still close to past trough levels. \#\# APPENDIX - FINANCIAL FORECASTS}
\end{figure}
\begin{figure}[htbp]
\centering
\includegraphics[width=0.88\linewidth]{figures/fig_127.jpg}
\caption{EXHIBIT 26 - EXHIBIT 26: We also use 6.2x P/E for hynix. EXHIBIT 27: Micron target P/E is 7.7x, still close to past trough levels. \#\# APPENDIX - FINANCIAL FORECASTS \#\# EXHIBIT 28: Samsung Income Statement Figures in KRW B Unless Otherwise St}
\end{figure}
\begin{figure}[htbp]
\centering
\includegraphics[width=0.88\linewidth]{figures/fig_227.jpg}
\caption{Exhibit 1 - Exhibit 1: Databricks core growth has accelerated over the last 5 quarters as a function of its product market fit with AI; Snowflake is also starting to show acceleration, albeit to a lesser extent \textbackslash{}*Databricks expects ARR to grow}
\end{figure}
{\small\color{refgray}\textbf{References:} [R003] Bernstein：云巨头正在从“算力军备竞赛”转向“资本效率竞赛”; [R008] Bernstein：AI正被内存成本压到喘不过气，一场供应链“再校准”不可避免; [R020] GS：软件行业的真正分水岭不是AI本身，而是“好粘性”与“坏粘性”的重新定价}

\section{半导体设备：DRAM与NAND的结构性分化}
\textbf{半导体设备支出的真正变量不是总量，而是DRAM提前支出与NAND等待拐点的结构性切换。}

\begin{itemize}
  \item MS将2026年全球DRAM设备支出从446亿美元上调至486亿美元，2027年从571亿美元上调至625亿美元，主要驱动力来自现有产线升级（brownfield），而非新建厂。
  \item 客户正在'不惜一切代价'将DRAM设备采购前置，brownfield部分的上修远超greenfield，反映的是供给侧的主动加速而非需求拉动。
  \item NAND设备支出预测几乎原地踏步：2026年从159亿美元微调至157亿美元，2027年从240亿美元调整至241亿美元。
  \item 但2027年NAND设备支出增速将高达53\%，远超DRAM的29\%，市场对设备股的定价可能忽略这一即将到来的结构性切换。
  \item NAND bit供给增速从21\%/32\%上调至24\%/31\%，关键变量不是产能扩张，而是良率提升——同样设备能产出更多bit。
  \item 设备商受益顺序：2026年DRAM支出加速，Applied Materials受益最直接；2027年NAND增速反超，Lam Research可能更受益。
  \item HBM的'trade ratio'导致约27\%的2028年bits被'损耗'，进一步加剧DRAM供给紧张。
\end{itemize}
\begin{figure}[htbp]
\centering
\includegraphics[width=0.88\linewidth]{figures/fig_086.jpg}
\caption{Exhibit 1 - Exhibit 1: We now forecast DRAM WFE to average \textbackslash{}\$58bn during 2026-28 and average bit output of 31\% Exhibit 2: As we exit 2026/27/28 at 2,251/2,661/3,081kwpm Exhibit 3: Wafer adds will be driven by all 4 major players}
\end{figure}
\begin{figure}[htbp]
\centering
\includegraphics[width=0.88\linewidth]{figures/fig_087.jpg}
\caption{Exhibit 1 - Exhibit 1: We now forecast DRAM WFE to average \textbackslash{}\$58bn during 2026-28 and average bit output of 31\% Exhibit 2: As we exit 2026/27/28 at 2,251/2,661/3,081kwpm Exhibit 3: Wafer adds will be driven by all 4 major players Exhibit}
\end{figure}
\begin{figure}[htbp]
\centering
\includegraphics[width=0.88\linewidth]{figures/fig_088.jpg}
\caption{Exhibit 2 - Exhibit 2: As we exit 2026/27/28 at 2,251/2,661/3,081kwpm Exhibit 3: Wafer adds will be driven by all 4 major players Exhibit 5: HBM's trade ratio is driving approx. 27\% of 2028 bits to be "lost bits" Exhibit 4: Non-HBM wafer}
\end{figure}
\begin{figure}[htbp]
\centering
\includegraphics[width=0.88\linewidth]{figures/fig_089.jpg}
\caption{Exhibit 3 - Exhibit 3: Wafer adds will be driven by all 4 major players Exhibit 5: HBM's trade ratio is driving approx. 27\% of 2028 bits to be "lost bits" Exhibit 4: Non-HBM wafer adds will surpass HBM wafer adds from 2027 Exhibit 6: As}
\end{figure}
\begin{figure}[htbp]
\centering
\includegraphics[width=0.88\linewidth]{figures/fig_090.jpg}
\caption{Exhibit 3 - Exhibit 7: Our NAND WFE forecast doesn't materially change, but we continue to see capacity adds as a contributor to bit supply in 2027/28}
\end{figure}
\begin{figure}[htbp]
\centering
\includegraphics[width=0.88\linewidth]{figures/fig_091.jpg}
\caption{Exhibit 4 - Exhibit 8: We model eSSD bits to grow at a 3-year CAGR of 66\% from 2025-28}
\end{figure}
\begin{figure}[htbp]
\centering
\includegraphics[width=0.88\linewidth]{figures/fig_092.jpg}
\caption{Exhibit 6 - Exhibit 9: We model eSSD mix to increase from 29\% in 2025 to 67\% in 2028}
\end{figure}
\begin{figure}[htbp]
\centering
\includegraphics[width=0.88\linewidth]{figures/fig_093.jpg}
\caption{Exhibit 7 - Exhibit 10: Our 2026-28 NAND WFE forecast translates to bit supply of 29\% during 2026-28}
\end{figure}
\begin{figure}[htbp]
\centering
\includegraphics[width=0.88\linewidth]{figures/fig_094.jpg}
\caption{Exhibit 8 - Exhibit 11: We expect non-China players to be a meaningful contributor to NAND Wafer adds in 2027-28}
\end{figure}
\begin{figure}[htbp]
\centering
\includegraphics[width=0.88\linewidth]{figures/fig_095.jpg}
\caption{Exhibit 9 - Exhibit 12: We don't model NAND wafer starts to surpass prior peak of 1,762kwpm (Q3 2022) in our forecast horizon}
\end{figure}
\begin{figure}[htbp]
\centering
\includegraphics[width=0.88\linewidth]{figures/fig_096.jpg}
\caption{Exhibit 10 - Exhibit 10: Our 2026-28 NAND WFE forecast translates to bit supply of 29\% during 2026-28 Exhibit 11: We expect non-China players to be a meaningful contributor to NAND Wafer adds in 2027-28 Exhibit 12: We don't model NAND wafer}
\end{figure}
\begin{figure}[htbp]
\centering
\includegraphics[width=0.88\linewidth]{figures/fig_097.jpg}
\caption{Exhibit 11 - Exhibit 11: We expect non-China players to be a meaningful contributor to NAND Wafer adds in 2027-28 Exhibit 12: We don't model NAND wafer starts to surpass prior peak of 1,762kwpm (Q3 2022) in our forecast horizon \#\# Valuation}
\end{figure}
{\small\color{refgray}\textbf{References:} [R006] 摩根斯坦利：半导体设备市场的真正变量不是总量，而是DRAM与NAND的结构性分化; [R017] 摩根斯坦利：市场低估的不是AI需求，而是供给端结构性约束的定价权}

\section{能源与大宗：霍尔木兹冲击重塑油需，AI压缩深水油田周期}
\textbf{霍尔木兹冲击正在加速电动汽车渗透，结构性压缩石油需求；AI同时在上游压缩深水油田开发周期，降低盈亏平衡点。}

\begin{itemize}
  \item GS报告显示，自2月霍尔木兹冲击以来，全球电动汽车渗透率上升3.4个百分点至26.1\%的历史新高，15大市场中有12个出现上升。
  \item 中国是绝对主力，渗透率同期上升11.4个百分点，贡献全球增量的61\%；OECD国家贡献21\%。
  \item GS估算，到2027年12月，全球石油需求可能因此面临每日13万至32万桶的下行风险，其中存量替代效应（中国汽油销售同比下降超20\%）被市场低估。
  \item AI和数字化可将典型深水油田从发现到产油周期从12年压缩至7年（降幅38\%），90\%的时间节省发生在最终投资决策（FID）之前。
  \item GS估算，AI可使深水项目内部收益率提升3.5个百分点（从15.5\%升至19.0\%），盈亏平衡价格下降约15\%。
  \item AI杠杆中，已证明的包括钻井效率提升、地震数据处理加速（速度提升5-10倍）、产量优化；预测性维护和无人平台仍处试点阶段。
  \item 欧洲化工行业面临'双重挤压'：短期需求疲软与物流紊乱，中期中国产能外溢（3-4月中国对亚洲化工品出口增长70\%）和配方重构。
  \item GS圆桌会议共识：冲突结束可能暴露潜在过剩产能，欧洲本地采购比例已从3年前的80\%降至约60\%。
\end{itemize}
\begin{figure}[htbp]
\centering
\includegraphics[width=0.88\linewidth]{figures/fig_001.jpg}
\caption{Exhibit 1 - Exhibit 1: According to Our Nowcast, Global EV Car Sales Penetration Has Increased By 3.4pp Since the Beginning of the Hormuz Shock}
\end{figure}
\begin{figure}[htbp]
\centering
\includegraphics[width=0.88\linewidth]{figures/fig_002.jpg}
\caption{Exhibit 1 - Exhibit 3: The Increase in Global EV Car Sales Penetration Since February Has Been Driven 61\% By China, 21\% By OECD, and 19\% By Non OECD ex China}
\end{figure}
\begin{figure}[htbp]
\centering
\includegraphics[width=0.88\linewidth]{figures/fig_003.jpg}
\caption{Exhibit 1 - Exhibit 3: The Increase in Global EV Car Sales Penetration Since February Has Been Driven 61\% By China, 21\% By OECD, and 19\% By Non OECD ex China Exhibit 4: 12 Out of the 15 Largest EV Markets Have Experienced an Increase in EV C}
\end{figure}
\begin{figure}[htbp]
\centering
\includegraphics[width=0.88\linewidth]{figures/fig_004.jpg}
\caption{Exhibit 2 - Exhibit 4: 12 Out of the 15 Largest EV Markets Have Experienced an Increase in EV Car Sales Penetration Since February For 15 largest EV markets by 2025 sales volume. Solid bars indicate that May share is calculated using realize}
\end{figure}
\begin{figure}[htbp]
\centering
\includegraphics[width=0.88\linewidth]{figures/fig_005.jpg}
\caption{Exhibit 3 - Exhibit 6: EV Car Sales Have Accelerated Meaningfully Since February Across the World ex US}
\end{figure}
\begin{figure}[htbp]
\centering
\includegraphics[width=0.88\linewidth]{figures/fig_006.jpg}
\caption{Exhibit 6 - Exhibit 6: EV Car Sales Have Accelerated Meaningfully Since February Across the World ex US Red diamonds indicate data points that are nowcasted. The vertical line on each chart indicates the start of the Hormuz shock. Exhibit}
\end{figure}
\begin{figure}[htbp]
\centering
\includegraphics[width=0.88\linewidth]{figures/fig_007.jpg}
\caption{Exhibit 6 - Exhibit 6: EV Car Sales Have Accelerated Meaningfully Since February Across the World ex US Red diamonds indicate data points that are nowcasted. The vertical line on each chart indicates the start of the Hormuz shock. Exhibit}
\end{figure}
\begin{figure}[htbp]
\centering
\includegraphics[width=0.88\linewidth]{figures/fig_073.jpg}
\caption{Exhibit 1 - Exhibit 2: At 11-15\% commercial WACC, the 75th percentile of the oil cost curve would move from \textbackslash{}\$66/bbl to \textbackslash{}\$57/bbl}
\end{figure}
\begin{figure}[htbp]
\centering
\includegraphics[width=0.88\linewidth]{figures/fig_074.jpg}
\caption{Exhibit 1 - Exhibit 2: At 11-15\% commercial WACC, the 75th percentile of the oil cost curve would move from \textbackslash{}\$66/bbl to \textbackslash{}\$57/bbl}
\end{figure}
\begin{figure}[htbp]
\centering
\includegraphics[width=0.88\linewidth]{figures/fig_075.jpg}
\caption{Exhibit 2 - Exhibit 4: Drivers of the improvement of oil field economics through AI}
\end{figure}
\begin{figure}[htbp]
\centering
\includegraphics[width=0.88\linewidth]{figures/fig_076.jpg}
\caption{Exhibit 2 - Exhibit 4: Drivers of the improvement of oil field economics through AI \#\# AI levers on a producing offshore field Productivity and cost levers — sized for a typical deepwater hub. Sourced from operator and vendor calls.}
\end{figure}
\begin{figure}[htbp]
\centering
\includegraphics[width=0.88\linewidth]{figures/fig_077.jpg}
\caption{Exhibit 6 - Exhibit 6: We estimate AI and advanced computing could compress the full greenfield deepwater cycle — from licence award to first oil — from c.12 years to c.7 years, a savings of c.4 years (-c.38\%). Block durations match the indi}
\end{figure}
\begin{figure}[htbp]
\centering
\includegraphics[width=0.88\linewidth]{figures/fig_078.jpg}
\caption{Exhibit 7 - Exhibit 8: Around 80\% of the potential calendar savings sits in three phases on the critical path: interpretation, well design and seismic processing. Total cycle saving: 28 months (51 → 23 mo, -56\%). Marginal contribution per phas}
\end{figure}
\begin{figure}[htbp]
\centering
\includegraphics[width=0.88\linewidth]{figures/fig_079.jpg}
\caption{Exhibit 8 - Exhibit 8: Around 80\% of the potential calendar savings sits in three phases on the critical path: interpretation, well design and seismic processing. Total cycle saving: 28 months (51 → 23 mo, -56\%). Marginal contribution per phas}
\end{figure}
\begin{figure}[htbp]
\centering
\includegraphics[width=0.88\linewidth]{figures/fig_080.jpg}
\caption{Exhibit 9 - Exhibit 9: Appraisal: potential \textbackslash{}\textasciitilde{}40-50\% cycle compression, almost entirely from the subsurface workflow \#\# 3) Pre-FID/Concept: FEED — c.40-50\% time compression driven by three roughly equal levers Average FEED shortens from c.12}
\end{figure}
\begin{figure}[htbp]
\centering
\includegraphics[width=0.88\linewidth]{figures/fig_081.jpg}
\caption{Exhibit 10 - Exhibit 10: The phase compresses by only c.10\% with AI / digital tools — from c.4.5 years to c.4 years on a typical deepwater project — making it the least AI-impacted block in the lifecycle Development / Construction phase — stairc}
\end{figure}
\begin{figure}[htbp]
\centering
\includegraphics[width=0.88\linewidth]{figures/fig_082.jpg}
\caption{Exhibit 10 - Exhibit 10: The phase compresses by only c.10\% with AI / digital tools — from c.4.5 years to c.4 years on a typical deepwater project — making it the least AI-impacted block in the lifecycle Development / Construction phase — stairc}
\end{figure}
\begin{figure}[htbp]
\centering
\includegraphics[width=0.88\linewidth]{figures/fig_083.jpg}
\caption{Exhibit 11 - Exhibit 13: We estimate AI could lift the IRR of an average greenfield field by 3.5pp, with capex and time to market driving the majority IRR decomposition of our AI scenario on an average greenfield field}
\end{figure}
\begin{figure}[htbp]
\centering
\includegraphics[width=0.88\linewidth]{figures/fig_084.jpg}
\caption{Exhibit 11 - Exhibit 13: We estimate AI could lift the IRR of an average greenfield field by 3.5pp, with capex and time to market driving the majority IRR decomposition of our AI scenario on an average greenfield field This report draws on com}
\end{figure}
\begin{figure}[htbp]
\centering
\includegraphics[width=0.88\linewidth]{figures/fig_085.jpg}
\caption{Exhibit 12 - Exhibit 12: Capex and OCF profile of an average greenfield field, base vs AI scenario (\textbackslash{}\$bn) Exhibit 13: We estimate AI could lift the IRR of an average greenfield field by 3.5pp, with capex and time to market driving the majority}
\end{figure}
\begin{figure}[htbp]
\centering
\includegraphics[width=0.88\linewidth]{figures/fig_205.jpg}
\caption{Exhibit 1 - Exhibit 1: Under weak domestic demand conditions, better than anticipated feedstock supply and though inventory draw down, China's exports have ramped significantly in April thousand tonnes per month \#\# Valuation \& Risks}
\end{figure}
\begin{figure}[htbp]
\centering
\includegraphics[width=0.88\linewidth]{figures/fig_206.jpg}
\caption{Exhibit 1 - Exhibit 1: Under weak domestic demand conditions, better than anticipated feedstock supply and though inventory draw down, China's exports have ramped significantly in April thousand tonnes per month \#\# Valuation \& Risks \#\# BAS}
\end{figure}
{\small\color{refgray}\textbf{References:} [R001] GS：市场低估了霍尔木兹冲击对油需的结构性重塑; [R004] GS：AI正在重塑全球深水油田的经济模型，但市场低估了“时间压缩”的定价权; [R016] GS：欧洲化工的真正危机不是红海冲突，而是冲突结束后暴露的结构性过剩}

\section{地产与消费：中国内需磨底，日本零售整合拐点}
\textbf{中国消费端动力流失，家庭去杠杆结构性地压制内需；日本零售业在人力成本压力下接近加速整合拐点。}

\begin{itemize}
  \item 五月家电内需数据延续疲软：零售额同比下降16\%，空调国内出厂量同比下降15\%，低于市场预期的-9\%。
  \item GS指出，去年6月'以旧换新'政策效果消退后，从今年6月下旬起同比基数压力将明显缓解，但读数改善不等于需求复苏。
  \item 家电出口出现超预期改善：五月整体出口增速回正，空调出口从-12\%收窄至-4\%，美欧消费者信心回升提供支撑。
  \item 中国家庭资产从房产向金融资产的结构性迁移正在加速，GS预计股权和保险占比将显著扩大。
  \item 日本零售业超市经营利润率仅1.5\%，但人工成本率高达14\%，最低工资自2015年累计上涨超30\%，远超CPI约10\%的涨幅。
  \item Bernstein报告认为，日本消费者行为分裂（工作日药妆店快速购物 vs 周末折扣店'寻宝'）正在结构性淘汰综合超市（GMS）等中间业态。
  \item 日本零售业集中度提升空间巨大：前五大 grocer 市场份额仅约30\%，远低于西欧和美国的水平，整合窗口已打开。
  \item 欧盟暂缓对中国割草机器人征收临时反倾销税，为中国企业争取了6-12个月缓冲期，GS认为九号公司受益最为显著。
\end{itemize}
\begin{figure}[htbp]
\centering
\includegraphics[width=0.88\linewidth]{figures/fig_142.jpg}
\caption{Exhibit 2 - Exhibit 2: AC domestic ex-factory shipment growth weakened sequentially in May Domestic ex-factory shipment vol. (10,000 units) Exhibit 3: AC domestic shipment moderated in May (-15\% yoy in May vs -7\% yoy in Apr)... AC domestic e}
\end{figure}
\begin{figure}[htbp]
\centering
\includegraphics[width=0.88\linewidth]{figures/fig_143.jpg}
\caption{Exhibit 3 - Exhibit 7: Domestic refrigerator ex-factory shipments improved to -4\% yoy growth in Apr Refrigerator domestic ex-factory shipment vol. (LHS, 10,000 units); yoy growth (RHS, \%)}
\end{figure}
\begin{figure}[htbp]
\centering
\includegraphics[width=0.88\linewidth]{figures/fig_144.jpg}
\caption{Exhibit 3 - Exhibit 7: Domestic refrigerator ex-factory shipments improved to -4\% yoy growth in Apr Refrigerator domestic ex-factory shipment vol. (LHS, 10,000 units); yoy growth (RHS, \%) Exhibit 4: ...with Midea/Gree/Haier/Hisense's domestic}
\end{figure}
\begin{figure}[htbp]
\centering
\includegraphics[width=0.88\linewidth]{figures/fig_145.jpg}
\caption{Exhibit 7 - Exhibit 5: Washing machine domestic shipment growth stayed at -6\% yoy in Apr Exhibit 6: Haier/Midea grew above industry at -1\%/-1\% yoy in Apr Washing machine (WM) domestic ex-factory shipment yoy growth by company}
\end{figure}
\begin{figure}[htbp]
\centering
\includegraphics[width=0.88\linewidth]{figures/fig_146.jpg}
\caption{Exhibit 7 - Exhibit 5: Washing machine domestic shipment growth stayed at -6\% yoy in Apr Exhibit 6: Haier/Midea grew above industry at -1\%/-1\% yoy in Apr Washing machine (WM) domestic ex-factory shipment yoy growth by company Exhibit 8: Mi}
\end{figure}
\begin{figure}[htbp]
\centering
\includegraphics[width=0.88\linewidth]{figures/fig_147.jpg}
\caption{Exhibit 5 - Exhibit 5: Washing machine domestic shipment growth stayed at -6\% yoy in Apr Exhibit 6: Haier/Midea grew above industry at -1\%/-1\% yoy in Apr Washing machine (WM) domestic ex-factory shipment yoy growth by company Exhibit 8: Mi}
\end{figure}
\begin{figure}[htbp]
\centering
\includegraphics[width=0.88\linewidth]{figures/fig_148.jpg}
\caption{Exhibit 8 - Exhibit 9: Domestic VRF ex-factory sales improved to -6\% yoy in Apr VRF domestic sales (LHS, RMB mn); yoy growth (RHS, \%)}
\end{figure}
\begin{figure}[htbp]
\centering
\includegraphics[width=0.88\linewidth]{figures/fig_149.jpg}
\caption{Exhibit 9 - Exhibit 11: China white goods exports improved in Apr, with better EU/US consumer sentiment in May/Jun}
\end{figure}
\begin{figure}[htbp]
\centering
\includegraphics[width=0.88\linewidth]{figures/fig_150.jpg}
\caption{Exhibit 9 - Exhibit 11: China white goods exports improved in Apr, with better EU/US consumer sentiment in May/Jun Exhibit 12: US pending home sales, a leading indicator of existing home sales, showed improvements in May}
\end{figure}
\begin{figure}[htbp]
\centering
\includegraphics[width=0.88\linewidth]{figures/fig_151.jpg}
\caption{Exhibit 10 - Exhibit 12: US pending home sales, a leading indicator of existing home sales, showed improvements in May Exhibit 13: AC ex-factory export shipments decreased by 4\% yoy in May, up from -12\% yoy in Apr AC export ex-factory shipment}
\end{figure}
\begin{figure}[htbp]
\centering
\includegraphics[width=0.88\linewidth]{figures/fig_152.jpg}
\caption{Exhibit 11 - Exhibit 14: Washing machine export growth was +14\% yoy in Apr, vs. +7\% yoy in Mar}
\end{figure}
\begin{figure}[htbp]
\centering
\includegraphics[width=0.88\linewidth]{figures/fig_153.jpg}
\caption{Exhibit 12 - Exhibit 14: Washing machine export growth was +14\% yoy in Apr, vs. +7\% yoy in Mar Exhibit 15: Refrigerator export shipment growth improved to +13\% yoy in Apr, vs. -3\% yoy in Mar Refrigerator export ex-factory shipment vol. (LHS, 1}
\end{figure}
\begin{figure}[htbp]
\centering
\includegraphics[width=0.88\linewidth]{figures/fig_154.jpg}
\caption{Exhibit 13 - Exhibit 15: Refrigerator export shipment growth improved to +13\% yoy in Apr, vs. -3\% yoy in Mar Refrigerator export ex-factory shipment vol. (LHS, 10,000 units); yoy growth (RHS, \%) Exhibit 17: Shipping rates increased in the past}
\end{figure}
\begin{figure}[htbp]
\centering
\includegraphics[width=0.88\linewidth]{figures/fig_155.jpg}
\caption{Exhibit 14 - Exhibit 17: Shipping rates increased in the past month... Weekly CCFI \#\# Exhibit 16: VRF export ex-factory sales grew +1\% yoy in Apr}
\end{figure}
\begin{figure}[htbp]
\centering
\includegraphics[width=0.88\linewidth]{figures/fig_156.jpg}
\caption{Exhibit 15 - Exhibit 16: VRF export ex-factory sales grew +1\% yoy in Apr VRF export sales (LHS, RMB mn); yoy growth (RHS, \%)}
\end{figure}
\begin{figure}[htbp]
\centering
\includegraphics[width=0.88\linewidth]{figures/fig_157.jpg}
\caption{Exhibit 17 - Exhibit 17: Shipping rates increased in the past month... Weekly CCFI \#\# Exhibit 16: VRF export ex-factory sales grew +1\% yoy in Apr VRF export sales (LHS, RMB mn); yoy growth (RHS, \%) Exhibit 18: ... and yoy growth increased as}
\end{figure}
\begin{figure}[htbp]
\centering
\includegraphics[width=0.88\linewidth]{figures/fig_158.jpg}
\caption{Exhibit 16 - Exhibit 16: VRF export ex-factory sales grew +1\% yoy in Apr VRF export sales (LHS, RMB mn); yoy growth (RHS, \%) Exhibit 18: ... and yoy growth increased as of mid-June Weekly CCFI yoy growth \% \#\# Retail sales remained weak on a}
\end{figure}
\begin{figure}[htbp]
\centering
\includegraphics[width=0.88\linewidth]{figures/fig_159.jpg}
\caption{Exhibit 22 - Exhibit 24: We note select brands like Midea and Xiaomi lowered prices during 618 vs 2025 Singles' Day, but higher than last 618}
\end{figure}
\begin{figure}[htbp]
\centering
\includegraphics[width=0.88\linewidth]{figures/fig_160.jpg}
\caption{Exhibit 22 - Exhibit 24: We note select brands like Midea and Xiaomi lowered prices during 618 vs 2025 Singles' Day, but higher than last 618 Pricing of entry-level AC}
\end{figure}
\begin{figure}[htbp]
\centering
\includegraphics[width=0.88\linewidth]{figures/fig_161.jpg}
\caption{Exhibit 23 - Exhibit 23: Midea, Xiaomi and Haier showed mom market share increase in Apr, while Gree lost share mom Online market share for AC on Taobao, Tmall and JD (\%) Exhibit 24: We note select brands like Midea and Xiaomi lowered prices d}
\end{figure}
\begin{figure}[htbp]
\centering
\includegraphics[width=0.88\linewidth]{figures/fig_162.jpg}
\caption{Exhibit 23 - Exhibit 23: Midea, Xiaomi and Haier showed mom market share increase in Apr, while Gree lost share mom Online market share for AC on Taobao, Tmall and JD (\%) Exhibit 24: We note select brands like Midea and Xiaomi lowered prices d}
\end{figure}
\begin{figure}[htbp]
\centering
\includegraphics[width=0.88\linewidth]{figures/fig_163.jpg}
\caption{Exhibit 25 - Exhibit 27: Growth of RVC moderated mom on Amazon US in May}
\end{figure}
\begin{figure}[htbp]
\centering
\includegraphics[width=0.88\linewidth]{figures/fig_164.jpg}
\caption{Exhibit 25 - Exhibit 27: Growth of RVC moderated mom on Amazon US in May Sales/Volume/ASP yoy growth by category \& Market share/Sales/Volume/ASP yoy growth by company on EC channel (Amazon US) Monthly sales data showed volatility so the absolute}
\end{figure}
\begin{figure}[htbp]
\centering
\includegraphics[width=0.88\linewidth]{figures/fig_165.jpg}
\caption{Exhibit 28 - Exhibit 28: Robotic lawn mower demand remained solid with DD\% growth and Ninebot's continued market share gains with \$100\textbackslash{}\%+\$ growth in May Download growth of leading robotic lawn mower apps (yoy, \%) Exhibit 29: Residential prop}
\end{figure}
\begin{figure}[htbp]
\centering
\includegraphics[width=0.88\linewidth]{figures/fig_166.jpg}
\caption{Exhibit 28 - Exhibit 28: Robotic lawn mower demand remained solid with DD\% growth and Ninebot's continued market share gains with \$100\textbackslash{}\%+\$ growth in May Download growth of leading robotic lawn mower apps (yoy, \%) Exhibit 29: Residential prop}
\end{figure}
\begin{figure}[htbp]
\centering
\includegraphics[width=0.88\linewidth]{figures/fig_167.jpg}
\caption{Exhibit 29 - Exhibit 29: Residential property completions decline narrowed slightly to -20\% in May (vs. -22\% in Apr), 45\% below pre-Covid levels Residential property completions (LHS, mn sqm); yoy growth (RHS, \%); growth vs. 2019 (RHS, \%) Exhi}
\end{figure}
\begin{figure}[htbp]
\centering
\includegraphics[width=0.88\linewidth]{figures/fig_168.jpg}
\caption{Exhibit 29 - Exhibit 29: Residential property completions decline narrowed slightly to -20\% in May (vs. -22\% in Apr), 45\% below pre-Covid levels Residential property completions (LHS, mn sqm); yoy growth (RHS, \%); growth vs. 2019 (RHS, \%) Exhi}
\end{figure}
\begin{figure}[htbp]
\centering
\includegraphics[width=0.88\linewidth]{figures/fig_169.jpg}
\caption{Exhibit 31 - Exhibit 32: SHFE copper price increased mildly in the past month, +2\%/+34\% mom/yoy as of mid-Jun SHFE copper price (LHS, RMB/ton); yoy growth (RHS, \%) Exhibit 33: SHFE aluminum price grew by -2\%/+16\% mom/yoy as of mid-Jun}
\end{figure}
\begin{figure}[htbp]
\centering
\includegraphics[width=0.88\linewidth]{figures/fig_170.jpg}
\caption{Exhibit 31 - Exhibit 32: SHFE copper price increased mildly in the past month, +2\%/+34\% mom/yoy as of mid-Jun SHFE copper price (LHS, RMB/ton); yoy growth (RHS, \%) Exhibit 33: SHFE aluminum price grew by -2\%/+16\% mom/yoy as of mid-Jun SHFE alu}
\end{figure}
\begin{figure}[htbp]
\centering
\includegraphics[width=0.88\linewidth]{figures/fig_171.jpg}
\caption{Exhibit 32 - Exhibit 34: Our covered appliances companies' share prices changed by -25\% to +14\% in the past month, yielding YTD returns of -38\% to +12\% Prices as of May 19, 2026.}
\end{figure}
\begin{figure}[htbp]
\centering
\includegraphics[width=0.88\linewidth]{figures/fig_172.jpg}
\caption{Exhibit 34 - Exhibit 34: Our covered appliances companies' share prices changed by -25\% to +14\% in the past month, yielding YTD returns of -38\% to +12\% Prices as of May 19, 2026. Exhibit 35: Valuations of most covered appliances companies show}
\end{figure}
\begin{figure}[htbp]
\centering
\includegraphics[width=0.88\linewidth]{figures/fig_173.jpg}
\caption{Exhibit 34 - Exhibit 34: Our covered appliances companies' share prices changed by -25\% to +14\% in the past month, yielding YTD returns of -38\% to +12\% Prices as of May 19, 2026. Exhibit 35: Valuations of most covered appliances companies show}
\end{figure}
\begin{figure}[htbp]
\centering
\includegraphics[width=0.88\linewidth]{figures/fig_216.jpg}
\caption{Exhibit 4 - EXHIBIT 2: Japan has the highest number of retail stores per capita among the developed nations Retail density per 1,000 inhabitants Year: 2025 EXHIBIT 3: Per-kilometer density makes the over-supply even starker-Japan's 2.61 stor}
\end{figure}
\begin{figure}[htbp]
\centering
\includegraphics[width=0.88\linewidth]{figures/fig_217.jpg}
\caption{EXHIBIT 3 - EXHIBIT 3: Per-kilometer density makes the over-supply even starker-Japan's 2.61 stores per sq km is 2x the next-closest peer Retail density per square kilometer Year: 2025 EXHIBIT 4: The math of the shakeout—the survivors inheri}
\end{figure}
\begin{figure}[htbp]
\centering
\includegraphics[width=0.88\linewidth]{figures/fig_218.jpg}
\caption{EXHIBIT 4 - EXHIBIT 4: The math of the shakeout—the survivors inherit a larger slice of a stagnant pie \#\# THIS OVER-STORED MARKET IS NOW TOO EXPENSIVE TO REMAIN FRAGMENTED Global grocery retail has largely consolidated into a game of scale,}
\end{figure}
\begin{figure}[htbp]
\centering
\includegraphics[width=0.88\linewidth]{figures/fig_219.jpg}
\caption{Exhibit 6 - EXHIBIT 5: The consolidation runway made quantitative—Japan has 40+ percentage points of concentration to catch up Top 5 grocery retailers' combined market share by country Year: 2025 Before Why fragmentation survived EXHIBIT 6:}
\end{figure}
\begin{figure}[htbp]
\centering
\includegraphics[width=0.88\linewidth]{figures/fig_220.jpg}
\caption{EXHIBIT 6 - EXHIBIT 6: Three structural shields, all down simultaneously—the old fragmentation no longer has a defense After Why fragmentation is collapsing \#\# Pillar 1 Freshness culture Japan CR5\textbackslash{}\textasciitilde{}30\% (1990s-2010s) Store-level processing fo}
\end{figure}
\begin{figure}[htbp]
\centering
\includegraphics[width=0.88\linewidth]{figures/fig_221.jpg}
\caption{EXHIBIT 7 - EXHIBIT 7: Grocery is Japan's largest retail category by a wide margin 2024: Japan retail sales breakdown by industry}
\end{figure}
\begin{figure}[htbp]
\centering
\includegraphics[width=0.88\linewidth]{figures/fig_222.jpg}
\caption{EXHIBIT 7 - EXHIBIT 9: Why "over-stored" becomes "impossibly overstored"—the denominator itself is shrinking}
\end{figure}
\begin{figure}[htbp]
\centering
\includegraphics[width=0.88\linewidth]{figures/fig_223.jpg}
\caption{EXHIBIT 8 - EXHIBIT 8: A long tail waiting to be consolidated—the top two names hold barely one-fifth of the market FY2024 Japan grocery retail sales breakdown by company EXHIBIT 9: Why "over-stored" becomes "impossibly overstored"—the denom}
\end{figure}
\begin{figure}[htbp]
\centering
\includegraphics[width=0.88\linewidth]{figures/fig_224.jpg}
\caption{EXHIBIT 9 - EXHIBIT 9: Why "over-stored" becomes "impossibly overstored"—the denominator itself is shrinking EXHIBIT 10: The inflation that finally matters—minimum wage has broken out above CPI, and the weak supermarkets cannot follow EXHI}
\end{figure}
\begin{figure}[htbp]
\centering
\includegraphics[width=0.88\linewidth]{figures/fig_225.jpg}
\caption{EXHIBIT 10 - EXHIBIT 10: The inflation that finally matters—minimum wage has broken out above CPI, and the weak supermarkets cannot follow EXHIBIT 11: Rising labor costs are pushing supermarkets toward a breaking point – only scale can solve i}
\end{figure}
\begin{figure}[htbp]
\centering
\includegraphics[width=0.88\linewidth]{figures/fig_226.jpg}
\caption{EXHIBIT 10 - EXHIBIT 10: The inflation that finally matters—minimum wage has broken out above CPI, and the weak supermarkets cannot follow EXHIBIT 11: Rising labor costs are pushing supermarkets toward a breaking point – only scale can solve i}
\end{figure}
{\small\color{refgray}\textbf{References:} [R011] GS：家电内需压力尚未出清，但出口的结构性机会比想象中更值得关注; [R014] GS：欧盟暂缓对中国割草机器人征收临时反倾销税，真正的变量不是关税本身，而是企业能否利用这个窗口期重塑竞争壁垒; [R015] GS：中国家庭资产正在经历从房产到金融资产的结构性迁移; [R019] Bernstein：日本零售业真正被低估的不是消费复苏，而是供给侧的整合拐点}

\section{区域市场：A股硬科技领跑，印度AI自主性困境}
\textbf{A股硬科技正主导中国盈利增长结构性分化，而印度面临在借来的模型上构建AI未来的结构性风险。}

\begin{itemize}
  \item 2026年Q1，CSI300盈利同比增长9\%，创业板+22\%，科创板+198\%，而MSCI中国指数同比下滑8\%，A股与离岸市场盈利鸿沟显著。
  \item GS维持对A股相对于H股的偏好，预测CSI300全年盈利增速20\%，MSCI中国仅8\%，差距源于A股'硬科技'权重远高于离岸市场的'软科技'。
  \item 中国AI供应链内部，硬件利润占比从2024年的41\%反弹至54\%，半导体（存储、IC设计、材料）利润份额提升最快，AI应用利润占比从59\%跌至46\%。
  \item 中国上市公司全球化仍在深化：海外销售占比升至18\%，汽车和医药是全球化领先行业，且海外市场盈利能力普遍优于国内。
  \item Bernstein报告指出，印度无法在借来的模型上构建AI未来，基础模型正从SaaS产品转变为受地缘政治管控的战略资产。
  \item 印度科技生态长期以服务外包为主，缺乏面向消费者的平台生成高质量训练数据，导致未能出现自己的'DeepSeek时刻'。
  \item 印度AI战略资金分散、缺乏重点，政府1040亿卢比AI拨款分配过于分散，未形成任何单一领域的显著优势。
  \item 日本MLCC相关股票的上涨并非单纯由AI需求驱动，JPM量化团队认为其本质是'供给瓶颈溢价'，与2000年代EM热潮中的瓶颈股高度相似。
\end{itemize}
\begin{figure}[htbp]
\centering
\includegraphics[width=0.88\linewidth]{figures/fig_174.jpg}
\caption{Exhibit 1 - Exhibit 1: A-share tech/small-mid caps led year-on-year earnings growth in 1Q26 Exhibit 2: Better earnings revision trend in A-shares/IT compared to MSCI China/Internet Exhibit 3: We now expect 8\%/12\% earnings growth for MSCI C}
\end{figure}
\begin{figure}[htbp]
\centering
\includegraphics[width=0.88\linewidth]{figures/fig_175.jpg}
\caption{Exhibit 1 - Exhibit 4: Profit growth momentum for offshore tech could improve in 2H26 per GS and consensus expectations}
\end{figure}
\begin{figure}[htbp]
\centering
\includegraphics[width=0.88\linewidth]{figures/fig_176.jpg}
\caption{Exhibit 2 - Exhibit 5: The 1H/2Q26 earnings season will kick off in August}
\end{figure}
\begin{figure}[htbp]
\centering
\includegraphics[width=0.88\linewidth]{figures/fig_177.jpg}
\caption{Exhibit 3 - Exhibit 3: We now expect 8\%/12\% earnings growth for MSCI China in 2026/27E, below consensus estimates Exhibit 4: Profit growth momentum for offshore tech could improve in 2H26 per GS and consensus expectations Exhibit 5: The 1H}
\end{figure}
\begin{figure}[htbp]
\centering
\includegraphics[width=0.88\linewidth]{figures/fig_178.jpg}
\caption{Exhibit 4 - Exhibit 4: Profit growth momentum for offshore tech could improve in 2H26 per GS and consensus expectations Exhibit 5: The 1H/2Q26 earnings season will kick off in August \#\# 2) AI profits shifting towards hardware, both globa}
\end{figure}
\begin{figure}[htbp]
\centering
\includegraphics[width=0.88\linewidth]{figures/fig_179.jpg}
\caption{Exhibit 5 - Exhibit 5: The 1H/2Q26 earnings season will kick off in August \#\# 2) AI profits shifting towards hardware, both globally and domestically Global AI profit pool: We laid out our global AI equity universe maps in March, covering}
\end{figure}
\begin{figure}[htbp]
\centering
\includegraphics[width=0.88\linewidth]{figures/fig_180.jpg}
\caption{Exhibit 9 - Exhibit 7: China's share in global AI profit pool declined as other regions gained Exhibit 8: The margin profile for Chinese AI companies still lags global peers Exhibit 9: Both US and Chinese hyperscalers keep scaling up AI Ca}
\end{figure}
\begin{figure}[htbp]
\centering
\includegraphics[width=0.88\linewidth]{figures/fig_181.jpg}
\caption{Exhibit 9 - Exhibit 10: The semiconductors layer grew the fastest in the Chinese AI universe in 1Q26}
\end{figure}
\begin{figure}[htbp]
\centering
\includegraphics[width=0.88\linewidth]{figures/fig_182.jpg}
\caption{Exhibit 8 - Exhibit 11: Within China's AI supply chain, the internal profit pool has undergone a rebound toward hardware}
\end{figure}
\begin{figure}[htbp]
\centering
\includegraphics[width=0.88\linewidth]{figures/fig_183.jpg}
\caption{Exhibit 9 - Exhibit 9: Both US and Chinese hyperscalers keep scaling up AI Capex estimates by GS Internet team (US\textbackslash{}\$bn) Exhibit 10: The semiconductors layer grew the fastest in the Chinese AI universe in 1Q26 Exhibit 11: Within China's AI}
\end{figure}
\begin{figure}[htbp]
\centering
\includegraphics[width=0.88\linewidth]{figures/fig_184.jpg}
\caption{Exhibit 10 - Exhibit 10: The semiconductors layer grew the fastest in the Chinese AI universe in 1Q26 Exhibit 11: Within China's AI supply chain, the internal profit pool has undergone a rebound toward hardware \#\# 3) Recovering aggregate rev}
\end{figure}
\begin{figure}[htbp]
\centering
\includegraphics[width=0.88\linewidth]{figures/fig_185.jpg}
\caption{Exhibit 33 - Exhibit 12: Aggregate revenue growth for the All China universe accelerated to 4\% in 1Q26 vs. +1\% in 2025 Exhibit 13: Mixed margin profile across the “involuted” sectors in 1Q26 \#\# 4) "Going Global" is still ongoing}
\end{figure}
\begin{figure}[htbp]
\centering
\includegraphics[width=0.88\linewidth]{figures/fig_186.jpg}
\caption{Exhibit 33 - Exhibit 12: Aggregate revenue growth for the All China universe accelerated to 4\% in 1Q26 vs. +1\% in 2025 Exhibit 13: Mixed margin profile across the “involuted” sectors in 1Q26 \#\# 4) "Going Global" is still ongoing Sustained mo}
\end{figure}
\begin{figure}[htbp]
\centering
\includegraphics[width=0.88\linewidth]{figures/fig_187.jpg}
\caption{Exhibit 16 - Exhibit 14: Foreign sales exposure of Chinese listed universe rose further to 18\% despite tariff headwinds Exhibit 15: Autos and Pharma are the leaders in the globalization trend Exhibit 16: Most sectors enjoy better profitabili}
\end{figure}
\begin{figure}[htbp]
\centering
\includegraphics[width=0.88\linewidth]{figures/fig_188.jpg}
\caption{Exhibit 14 - Exhibit 17: A slight FX loss was recorded in 2025 due to CNY appreciation}
\end{figure}
\begin{figure}[htbp]
\centering
\includegraphics[width=0.88\linewidth]{figures/fig_189.jpg}
\caption{Exhibit 15 - Exhibit 15: Autos and Pharma are the leaders in the globalization trend Exhibit 16: Most sectors enjoy better profitability in the overseas markets Exhibit 17: A slight FX loss was recorded in 2025 due to CNY appreciation \#\# 5}
\end{figure}
\begin{figure}[htbp]
\centering
\includegraphics[width=0.88\linewidth]{figures/fig_190.jpg}
\caption{Exhibit 16 - Exhibit 16: Most sectors enjoy better profitability in the overseas markets Exhibit 17: A slight FX loss was recorded in 2025 due to CNY appreciation \#\# 5) Record cash returns along with more R\&D and acquisitions Record-high cas}
\end{figure}
\begin{figure}[htbp]
\centering
\includegraphics[width=0.88\linewidth]{figures/fig_191.jpg}
\caption{Exhibit 18 - Exhibit 18: Chinese corporates spent more cash on dividends, acquisitions, and R\&D in 2025 (Rmb tn/\%) All China listed universe - Use of Cash Exhibit 19: Dividend payout ratio edged up to 39\% in 2025 Exhibit 20: Net buybacks rem}
\end{figure}
\begin{figure}[htbp]
\centering
\includegraphics[width=0.88\linewidth]{figures/fig_192.jpg}
\caption{Exhibit 18 - Exhibit 18: Chinese corporates spent more cash on dividends, acquisitions, and R\&D in 2025 (Rmb tn/\%) All China listed universe - Use of Cash Exhibit 19: Dividend payout ratio edged up to 39\% in 2025 Exhibit 20: Net buybacks rem}
\end{figure}
\begin{figure}[htbp]
\centering
\includegraphics[width=0.88\linewidth]{figures/fig_193.jpg}
\caption{Exhibit 19 - Exhibit 19: Dividend payout ratio edged up to 39\% in 2025 Exhibit 20: Net buybacks remain accretive to EPS despite lower absolute volumes Represents the net change in the carrying value of common and preferred stock 6) Key theme}
\end{figure}
\begin{figure}[htbp]
\centering
\includegraphics[width=0.88\linewidth]{figures/fig_194.jpg}
\caption{Exhibit 24 - Exhibit 24: The refreshed portfolio exhibits strong performance and EPS revision trends Exhibit 25: The portfolio trades at 23x fP/E and 1.4x fPEG on an equal-weighted basis}
\end{figure}
\begin{figure}[htbp]
\centering
\includegraphics[width=0.88\linewidth]{figures/fig_195.jpg}
\caption{Exhibit 24 - Exhibit 24: The refreshed portfolio exhibits strong performance and EPS revision trends Exhibit 25: The portfolio trades at 23x fP/E and 1.4x fPEG on an equal-weighted basis}
\end{figure}
\begin{figure}[htbp]
\centering
\includegraphics[width=0.88\linewidth]{figures/fig_196.jpg}
\caption{Exhibit 27 - Exhibit 27: Top-line growth may pick up for MSCI China universe Exhibit 28: Exporters saw resilient top-line and earnings growth in 1Q26 Exhibit 29: All listed developers incurred around Rmb400bn of net losses in 2025}
\end{figure}
\begin{figure}[htbp]
\centering
\includegraphics[width=0.88\linewidth]{figures/fig_197.jpg}
\caption{Exhibit 27 - Exhibit 30: Still around 60\% of companies in property sector are loss-making}
\end{figure}
\begin{figure}[htbp]
\centering
\includegraphics[width=0.88\linewidth]{figures/fig_198.jpg}
\caption{Exhibit 28 - Exhibit 30: Still around 60\% of companies in property sector are loss-making Including both developers and property management companies; \% in recent years are underestimated as some real estate companies stopped/delayed disclosin}
\end{figure}
\begin{figure}[htbp]
\centering
\includegraphics[width=0.88\linewidth]{figures/fig_199.jpg}
\caption{Exhibit 29 - Exhibit 33: Consensus estimates Capex growth to moderate in most sectors while tech hardware notably upscales the Capex}
\end{figure}
\begin{figure}[htbp]
\centering
\includegraphics[width=0.88\linewidth]{figures/fig_128.jpg}
\caption{EXHIBIT 3 - EXHIBIT 3: Too lumpy, too little: India's AI strategy so far has been that of trying to do a lot from less resources Allocation and Actual Spend in India's AI mission (INR bn) \#\# DEFOCUSED, THINLY SPREAD Beyond funding, the broad}
\end{figure}
\begin{figure}[htbp]
\centering
\includegraphics[width=0.88\linewidth]{figures/fig_129.jpg}
\caption{EXHIBIT 4 - EXHIBIT 4: Too defocused, thinly spread: Even with lesser and lumpy allocations, the focus is all over the place, creating no single area with meaningful advantage for India How Government's INR 104 billion AI allocation is split}
\end{figure}
\begin{figure}[htbp]
\centering
\includegraphics[width=0.88\linewidth]{figures/fig_130.jpg}
\caption{Figure 3 - Figure 3: Factor exposure of electrical equipment (TSE 33 sector): As of June 12}
\end{figure}
\begin{figure}[htbp]
\centering
\includegraphics[width=0.88\linewidth]{figures/fig_131.jpg}
\caption{Figure 10 - Figure 3: Factor exposure of electrical equipment (TSE 33 sector): As of June 12 Figure 4: Factor exposure of two MLCC-related companies (Murata Manufacturing and Taiyo Yuden): As of June 12}
\end{figure}
\begin{figure}[htbp]
\centering
\includegraphics[width=0.88\linewidth]{figures/fig_132.jpg}
\caption{Figure 2 - Figure 4: Factor exposure of two MLCC-related companies (Murata Manufacturing and Taiyo Yuden): As of June 12 Figure 5: Estimated net position and profit margin momentum of Cash Equity Trend-Followers for two MLCC-related compan}
\end{figure}
\begin{figure}[htbp]
\centering
\includegraphics[width=0.88\linewidth]{figures/fig_133.jpg}
\caption{Figure 3 - Figure 6: Estimated net position and profit margin momentum of Cash Equity Trend-Followers for two MLCC-related companies: Taiyo Yuden}
\end{figure}
\begin{figure}[htbp]
\centering
\includegraphics[width=0.88\linewidth]{figures/fig_134.jpg}
\caption{Figure 4 - Figure 6: Estimated net position and profit margin momentum of Cash Equity Trend-Followers for two MLCC-related companies: Taiyo Yuden Figure 7: Regional Allocation of global active investment trusts (DM+EM) (mainland China stoc}
\end{figure}
\begin{figure}[htbp]
\centering
\includegraphics[width=0.88\linewidth]{figures/fig_135.jpg}
\caption{Figure 5 - Figure 7: Regional Allocation of global active investment trusts (DM+EM) (mainland China stocks + Hong Kong stocks) and US import ratios from China Figure 8: Relative differences in regional allocation of global active investmen}
\end{figure}
\begin{figure}[htbp]
\centering
\includegraphics[width=0.88\linewidth]{figures/fig_136.jpg}
\caption{Figure 6 - Figure 9: Regional allocation of global active investment trusts (DM+EM): As of end-April}
\end{figure}
\begin{figure}[htbp]
\centering
\includegraphics[width=0.88\linewidth]{figures/fig_137.jpg}
\caption{Figure 7 - Figure 9: Regional allocation of global active investment trusts (DM+EM): As of end-April Figure 10: Regional allocation of global active investment trusts (DM+EM) – Trends for Japan, South Korea, and Taiwan: As of end-April}
\end{figure}
\begin{figure}[htbp]
\centering
\includegraphics[width=0.88\linewidth]{figures/fig_138.jpg}
\caption{Figure 8 - Figure 9: Regional allocation of global active investment trusts (DM+EM): As of end-April Figure 10: Regional allocation of global active investment trusts (DM+EM) – Trends for Japan, South Korea, and Taiwan: As of end-April F}
\end{figure}
\begin{figure}[htbp]
\centering
\includegraphics[width=0.88\linewidth]{figures/fig_139.jpg}
\caption{Figure 9 - Figure 12: Changes in sector weightings of Japanese stocks in global active investment trusts (DM) (Active vs. Passive): As of end-April}
\end{figure}
\begin{figure}[htbp]
\centering
\includegraphics[width=0.88\linewidth]{figures/fig_140.jpg}
\caption{Figure 10 - Figure 12: Changes in sector weightings of Japanese stocks in global active investment trusts (DM) (Active vs. Passive): As of end-April}
\end{figure}
\begin{figure}[htbp]
\centering
\includegraphics[width=0.88\linewidth]{figures/fig_141.jpg}
\caption{Figure 11 - Figure 11: Sector weightings of Japanese stocks in global active investment trusts (DM) (Active vs. Passive): As of end-April Figure 12: Changes in sector weightings of Japanese stocks in global active investment trusts (DM) (Act}
\end{figure}
{\small\color{refgray}\textbf{References:} [R012] GS：A股硬科技正在主导中国盈利增长的结构性分化; [R009] Bernstein：印度真正需要的不是算力租用，而是自己的“DeepSeek 时刻”; [R010] JPM：市场低估了“供给瓶颈”在日本MLCC股中的定价权，而非AI需求本身}

\section{全球游戏业：GTA VI前的洗牌前夜}
\textbf{全球游戏业2026年增速放缓至0.7\%，行业从增长驱动转向整合驱动，GTA VI将成为情绪和资金流向的转折点。}

\begin{itemize}
  \item Bernstein报告显示，2025年全球游戏业收入2180亿美元，同比增长4.8\%，但2026年增速预计骤降至0.7\%，2027年回升至3.3\%。
  \item 行业正处于'资本周期后半段'：上市公司收入增速预计达7-8\%，显著高于行业平均水平，驱动力来自西方工作室关闭潮和头部公司IP/资金优势。
  \item GTA VI成为行业'清场'事件：几乎所有大作挤在9月发售以避免竞争，供给过剩导致非顶级IP游戏容易被淹没。
  \item PC端是最大惊喜：2025年中国PC游戏市场增长15\%（《黑神话：悟空》驱动），普通家用电脑显存已能流畅跑PS4/PS5世代游戏。
  \item 平台抽成下降：2026年3月谷歌和苹果宣布降低平台抽成（安卓从30\%降至20\%，iOS降至25\%），成为移动开发者盈利的顺风。
  \item 高内存成本是主机硬件利润率的明显拖累：DRAM+SSD占物料成本的比重在2026年Q2显著上升。
  \item AI对游戏业的影响有限：AI不解决可见性和IP认知问题，每年真正受欢迎的游戏数量基本恒定（约70款获得超5000条评论）。
\end{itemize}
\begin{figure}[htbp]
\centering
\includegraphics[width=0.88\linewidth]{figures/fig_008.jpg}
\caption{Exhibit 2 - EXHIBIT 2: We estimate global video gaming revenue reached US\textbackslash{}\$218bn in 2025, +4.8\% year on year EXHIBIT 3: Our video gaming model breaks down the global revenue pool by platform and by region EXHIBIT 4: 2025 was a strong year}
\end{figure}
\begin{figure}[htbp]
\centering
\includegraphics[width=0.88\linewidth]{figures/fig_009.jpg}
\caption{EXHIBIT 2 - EXHIBIT 5: North American growth stayed muted in 2025, while other regions rebounded}
\end{figure}
\begin{figure}[htbp]
\centering
\includegraphics[width=0.88\linewidth]{figures/fig_010.jpg}
\caption{EXHIBIT 3 - EXHIBIT 3: Our video gaming model breaks down the global revenue pool by platform and by region EXHIBIT 4: 2025 was a strong year for console, helped by the Switch 2 launch EXHIBIT 5: North American growth stayed muted in 2025,}
\end{figure}
\begin{figure}[htbp]
\centering
\includegraphics[width=0.88\linewidth]{figures/fig_011.jpg}
\caption{EXHIBIT 4 - EXHIBIT 4: 2025 was a strong year for console, helped by the Switch 2 launch EXHIBIT 5: North American growth stayed muted in 2025, while other regions rebounded 2025: Global gaming market revenue by region EXHIBIT 6: Asia-Pa}
\end{figure}
\begin{figure}[htbp]
\centering
\includegraphics[width=0.88\linewidth]{figures/fig_012.jpg}
\caption{EXHIBIT 5 - EXHIBIT 5: North American growth stayed muted in 2025, while other regions rebounded 2025: Global gaming market revenue by region EXHIBIT 6: Asia-Pacific is mobile dominated, while console leads elsewhere EXHIBIT 7: Console is}
\end{figure}
\begin{figure}[htbp]
\centering
\includegraphics[width=0.88\linewidth]{figures/fig_013.jpg}
\caption{EXHIBIT 6 - EXHIBIT 6: Asia-Pacific is mobile dominated, while console leads elsewhere EXHIBIT 7: Console is biggest in Europe and North America, while Asia dominates mobile and PC \#\# WINNERS AND LOSERS IN 2025 Company-level growth in the in}
\end{figure}
\begin{figure}[htbp]
\centering
\includegraphics[width=0.88\linewidth]{figures/fig_014.jpg}
\caption{Exhibit 8 - EXHIBIT 8: Tencent, Nintendo, and Roblox led incremental revenue growth in 2025 versus the prior year 2025 vs. 2024: Global video gaming companies incremental revenue 2025 vs. 2024: Global video gaming companies revenue growth EX}
\end{figure}
\begin{figure}[htbp]
\centering
\includegraphics[width=0.88\linewidth]{figures/fig_015.jpg}
\caption{EXHIBIT 8 - EXHIBIT 8: Tencent, Nintendo, and Roblox led incremental revenue growth in 2025 versus the prior year 2025 vs. 2024: Global video gaming companies incremental revenue 2025 vs. 2024: Global video gaming companies revenue growth EX}
\end{figure}
\begin{figure}[htbp]
\centering
\includegraphics[width=0.88\linewidth]{figures/fig_016.jpg}
\caption{Exhibit 14 - EXHIBIT 10: We’ve updated our tally of R\&D expenses per employee across the global video gaming industry... 2025: Global video game companies R\&D expenses per employee EXHIBIT 11: ...the fact Asian developers have a fraction of th}
\end{figure}
\begin{figure}[htbp]
\centering
\includegraphics[width=0.88\linewidth]{figures/fig_017.jpg}
\caption{EXHIBIT 11 - EXHIBIT 12: The companies we cover own some of the world's top IPs, an important input into video game sales success Top 15 bestselling game franchises on VGChartz EXHIBIT 13: The live service game scene is dominated by a collecti}
\end{figure}
\begin{figure}[htbp]
\centering
\includegraphics[width=0.88\linewidth]{figures/fig_018.jpg}
\caption{EXHIBIT 12 - EXHIBIT 14: The ongoing wave of video gaming studio shutdowns in Western markets was overdue in our view, and should help to move the industry to a more favourable part of the capital cycle}
\end{figure}
\begin{figure}[htbp]
\centering
\includegraphics[width=0.88\linewidth]{figures/fig_019.jpg}
\caption{EXHIBIT 12 - EXHIBIT 14: The ongoing wave of video gaming studio shutdowns in Western markets was overdue in our view, and should help to move the industry to a more favourable part of the capital cycle}
\end{figure}
\begin{figure}[htbp]
\centering
\includegraphics[width=0.88\linewidth]{figures/fig_020.jpg}
\caption{Exhibit 19 - EXHIBIT 15: We model 0.7\% aggregate industry revenue growth in 2026, slower year on year to reflect Playstation hardware declines, and a higher base for Nintendo EXHIBIT 16: We model faster PC growth continuing in 2026, low-single}
\end{figure}
\begin{figure}[htbp]
\centering
\includegraphics[width=0.88\linewidth]{figures/fig_021.jpg}
\caption{EXHIBIT 16 - EXHIBIT 16: We model faster PC growth continuing in 2026, low-single digit growth mobile growth, and lower console revenue growth than in 2025 EXHIBIT 17: Consensus forecasts point to slower growth in the next few quarters, follow}
\end{figure}
\begin{figure}[htbp]
\centering
\includegraphics[width=0.88\linewidth]{figures/fig_022.jpg}
\caption{EXHIBIT 17 - EXHIBIT 17: Consensus forecasts point to slower growth in the next few quarters, followed by faster growth into year-end, partly reflecting the GTA VI launch EXHIBIT 18: In hindsight, mid-2024 was the trough of the post-Covid down}
\end{figure}
\begin{figure}[htbp]
\centering
\includegraphics[width=0.88\linewidth]{figures/fig_023.jpg}
\caption{EXHIBIT 18 - EXHIBIT 18: In hindsight, mid-2024 was the trough of the post-Covid down cycle for video gaming industry revenue EXHIBIT 19: The industry's desire to avoid competing with the GTA VI launch has contributed to a logjam of new games}
\end{figure}
\begin{figure}[htbp]
\centering
\includegraphics[width=0.88\linewidth]{figures/fig_024.jpg}
\caption{EXHIBIT 21 - EXHIBIT 21: We think engagement concerns in video gaming are partly backward looking... EXHIBIT 23: PSN engagement has remained resilient into the end of the current console cycle}
\end{figure}
\begin{figure}[htbp]
\centering
\includegraphics[width=0.88\linewidth]{figures/fig_025.jpg}
\caption{EXHIBIT 21 - EXHIBIT 21: We think engagement concerns in video gaming are partly backward looking... EXHIBIT 23: PSN engagement has remained resilient into the end of the current console cycle Sony GNS - PSN MAUs EXHIBIT 24: ...while the s}
\end{figure}
\begin{figure}[htbp]
\centering
\includegraphics[width=0.88\linewidth]{figures/fig_026.jpg}
\caption{EXHIBIT 22 - EXHIBIT 25: The normalisation of Roblox engagement back near the pre-2025 trend-line has been interesting to observe, and puts claims around Roblox disrupting the industry into perspective}
\end{figure}
\begin{figure}[htbp]
\centering
\includegraphics[width=0.88\linewidth]{figures/fig_027.jpg}
\caption{EXHIBIT 23 - EXHIBIT 25: The normalisation of Roblox engagement back near the pre-2025 trend-line has been interesting to observe, and puts claims around Roblox disrupting the industry into perspective}
\end{figure}
\begin{figure}[htbp]
\centering
\includegraphics[width=0.88\linewidth]{figures/fig_028.jpg}
\caption{EXHIBIT 25 - EXHIBIT 25: The normalisation of Roblox engagement back near the pre-2025 trend-line has been interesting to observe, and puts claims around Roblox disrupting the industry into perspective \#\# THE BULL CASE FOR PC GAMING GROWTH Ass}
\end{figure}
\begin{figure}[htbp]
\centering
\includegraphics[width=0.88\linewidth]{figures/fig_029.jpg}
\caption{Exhibit 26 - EXHIBIT 26: Our views on PC growth last year align more with Newzoo's modelling of revenue acceleration EXHIBIT 27: Steam continues to grow faster than the aggregate PC market EXHIBIT 28: A majority of PCs on Steam can now play}
\end{figure}
\begin{figure}[htbp]
\centering
\includegraphics[width=0.88\linewidth]{figures/fig_030.jpg}
\caption{EXHIBIT 26 - EXHIBIT 29: The PC market in China saw a significant acceleration in growth following the Black Myth Wukong success... we think other markets are due to follow in the coming years}
\end{figure}
\begin{figure}[htbp]
\centering
\includegraphics[width=0.88\linewidth]{figures/fig_031.jpg}
\caption{EXHIBIT 27 - EXHIBIT 29: The PC market in China saw a significant acceleration in growth following the Black Myth Wukong success... we think other markets are due to follow in the coming years FY3/21-FY3/25: Capcom unit sales by region}
\end{figure}
\begin{figure}[htbp]
\centering
\includegraphics[width=0.88\linewidth]{figures/fig_032.jpg}
\caption{EXHIBIT 28 - EXHIBIT 30: Capcom's success has partly been driven by expansion on PC, and in emerging markets EXHIBIT 31: eFootball has also seen Steam CCUs ramp considerably in recent years}
\end{figure}
\begin{figure}[htbp]
\centering
\includegraphics[width=0.88\linewidth]{figures/fig_033.jpg}
\caption{EXHIBIT 29 - EXHIBIT 31: eFootball has also seen Steam CCUs ramp considerably in recent years \#\# WE USE SENSOR TOWER TO TRACK MOBILE GAMING BILLINGS}
\end{figure}
\begin{figure}[htbp]
\centering
\includegraphics[width=0.88\linewidth]{figures/fig_034.jpg}
\caption{EXHIBIT 30 - EXHIBIT 30: Capcom's success has partly been driven by expansion on PC, and in emerging markets EXHIBIT 31: eFootball has also seen Steam CCUs ramp considerably in recent years \#\# WE USE SENSOR TOWER TO TRACK MOBILE GAMING BILLI}
\end{figure}
\begin{figure}[htbp]
\centering
\includegraphics[width=0.88\linewidth]{figures/fig_035.jpg}
\caption{EXHIBIT 32 - EXHIBIT 32: We've modelled 1\% global mobile game billings growth in 2026 EXHIBIT 33: Google's announcement of lower platform fees will represent an earnings tailwind for mobile developers EXHIBIT 34: iOS platform fees were recen}
\end{figure}
\begin{figure}[htbp]
\centering
\includegraphics[width=0.88\linewidth]{figures/fig_036.jpg}
\caption{EXHIBIT 32 - EXHIBIT 32: We've modelled 1\% global mobile game billings growth in 2026 EXHIBIT 33: Google's announcement of lower platform fees will represent an earnings tailwind for mobile developers EXHIBIT 34: iOS platform fees were recen}
\end{figure}
\begin{figure}[htbp]
\centering
\includegraphics[width=0.88\linewidth]{figures/fig_037.jpg}
\caption{EXHIBIT 33 - EXHIBIT 33: Google's announcement of lower platform fees will represent an earnings tailwind for mobile developers EXHIBIT 34: iOS platform fees were recently cut in Europe and China, which helps too \#\# WE MODEL CONSOLE BOTTOM-U}
\end{figure}
\begin{figure}[htbp]
\centering
\includegraphics[width=0.88\linewidth]{figures/fig_038.jpg}
\caption{Exhibit 39 - EXHIBIT 37: Software sales revenue outweighs hardware sales for console platforms}
\end{figure}
\begin{figure}[htbp]
\centering
\includegraphics[width=0.88\linewidth]{figures/fig_039.jpg}
\caption{EXHIBIT 36 - EXHIBIT 36: This drives our estimate of the respective console installed bases FY3/08-FY3/28E: Console hardware install base EXHIBIT 37: Software sales revenue outweighs hardware sales for console platforms 2025: Sony and Nintendo}
\end{figure}
\begin{figure}[htbp]
\centering
\includegraphics[width=0.88\linewidth]{figures/fig_040.jpg}
\caption{EXHIBIT 36 - EXHIBIT 38: We model both Sony and Nintendo on the basis of billings, grossing up third-party games 2025: Sony and Nintendo game billing breakdown Hardware billings Software billings Subscription billings}
\end{figure}
\begin{figure}[htbp]
\centering
\includegraphics[width=0.88\linewidth]{figures/fig_041.jpg}
\caption{EXHIBIT 37 - EXHIBIT 39: High memory prices represent an obvious headwind for both in 2026 and onwards Q2 2026 vs. Q2 2025: DRAM+SSD share of bill of materials cost}
\end{figure}
\begin{figure}[htbp]
\centering
\includegraphics[width=0.88\linewidth]{figures/fig_042.jpg}
\caption{EXHIBIT 38 - EXHIBIT 40: Current contract pricing implies further downside to hardware margins \#\# AI ADOPTION UPDATE: JAPAN NOW ON BOARD}
\end{figure}
\begin{figure}[htbp]
\centering
\includegraphics[width=0.88\linewidth]{figures/fig_043.jpg}
\caption{EXHIBIT 39 - EXHIBIT 39: High memory prices represent an obvious headwind for both in 2026 and onwards Q2 2026 vs. Q2 2025: DRAM+SSD share of bill of materials cost EXHIBIT 40: Current contract pricing implies further downside to hardware marg}
\end{figure}
\begin{figure}[htbp]
\centering
\includegraphics[width=0.88\linewidth]{figures/fig_044.jpg}
\caption{Exhibit 44 - EXHIBIT 43: AI doesn't solve for visibility, IP and name recognition do... the number of truly popular games launched each year has remained roughly constant over time}
\end{figure}
\begin{figure}[htbp]
\centering
\includegraphics[width=0.88\linewidth]{figures/fig_045.jpg}
\caption{EXHIBIT 41 - EXHIBIT 43: AI doesn't solve for visibility, IP and name recognition do... the number of truly popular games launched each year has remained roughly constant over time EXHIBIT 44: The video gaming industry has increasingly found w}
\end{figure}
\begin{figure}[htbp]
\centering
\includegraphics[width=0.88\linewidth]{figures/fig_046.jpg}
\caption{EXHIBIT 42 - EXHIBIT 44: The video gaming industry has increasingly found ways to integrate AI into development workflows...}
\end{figure}
{\small\color{refgray}\textbf{References:} [R002] Bernstein：全球游戏业真正的拐点不在2026年，而在GTA VI之后的2027年}

\section{储能与新能源：供应链本地化重构拐点}
\textbf{全球储能产业竞争维度正从电池性能转向供应链的地理韧性，北美合规驱动的本地化制造加速落地。}

\begin{itemize}
  \item Bernstein报告揭示，市场关注的焦点正在从'谁的技术更先进'转向'谁的供应链更能在本土落地'，美国、欧洲、中国同步发生供应链本地化运动。
  \item 北美市场核心驱动力是合规驱动的本地化制造：HyVISION Systems拿下950亿韩元北美ESS设备订单（初始规模扩张近五倍），Sebang与LGES合作在俄亥俄州建设模块工厂。
  \item 电池技术路线之争仍在：GM可能优先推进LMR电池（成本接近LFP但能量密度高33\%），Stellantis开始路测Factorial半固态电池（375Wh/kg，快充18分钟）。
  \item CATL表示固态电池大规模商用化2027年前不太可能，成本仍是核心障碍。
  \item 中国需求结构从'EV单引擎'变成'EV+储能双引擎'，碳酸锂期货反弹反映需求结构变化。
  \item 数据中心供电的真正瓶颈在电网连接：98\%的受访运营商首选电网供电，但并网等待时间普遍超过4年，CAISO区域更久。
  \item 超过95\%的受访者认为自备备用电源是设计关键，天然气、太阳能+储能、燃料电池是主要选择，但自备电源目前仅占数据中心电力市场的约5\%。
\end{itemize}
\begin{figure}[htbp]
\centering
\includegraphics[width=0.88\linewidth]{figures/fig_200.jpg}
\caption{EXHIBIT 2 - EXHIBIT 3: Interconnection queues are \textbackslash{}\textasciitilde{}4+ years for all regions, CAISO is the worst served region... Duration from IR to COD by region Most respondents indicated a preference for grid connected power, supplemented by onsite reso}
\end{figure}
\begin{figure}[htbp]
\centering
\includegraphics[width=0.88\linewidth]{figures/fig_201.jpg}
\caption{EXHIBIT 2 - EXHIBIT 3: Interconnection queues are \textbackslash{}\textasciitilde{}4+ years for all regions, CAISO is the worst served region... Duration from IR to COD by region Most respondents indicated a preference for grid connected power, supplemented by onsite reso}
\end{figure}
\begin{figure}[htbp]
\centering
\includegraphics[width=0.88\linewidth]{figures/fig_202.jpg}
\caption{EXHIBIT 4 - EXHIBIT 5: 45\% of respondents preferred fossil fuel onsite back up; followed by Solar+ BESS, and Fuel Cells Which Type of Onsite Gen is preferred? Finally, We asked a long form question around how sustainability commitments shape}
\end{figure}
\begin{figure}[htbp]
\centering
\includegraphics[width=0.88\linewidth]{figures/fig_203.jpg}
\caption{EXHIBIT 5 - EXHIBIT 6: Survey question: What role, if any, do decarbonization or sustainability commitments play in shaping your organization's preference around power supply sourcing and configuration?}
\end{figure}
\begin{figure}[htbp]
\centering
\includegraphics[width=0.88\linewidth]{figures/fig_204.jpg}
\caption{EXHIBIT 6 - EXHIBIT 6: Survey question: What role, if any, do decarbonization or sustainability commitments play in shaping your organization's preference around power supply sourcing and configuration? Total 47 responses \#\# I. REQUIRED DISC}
\end{figure}
{\small\color{refgray}\textbf{References:} [R005] Bernstein：全球储能市场的真正拐点不在技术突破，而在供应链的本地化重构; [R013] Bernstein：数据中心供电的真正瓶颈不在技术，在电网的“连接组织”}

\section{代币化与数字资产：基础设施竞争决定胜负}
\textbf{代币化RWA正从'发行试验'进入'基础设施竞争'阶段，合规、高效、可规模化的交易和结算层是决胜关键。}

\begin{itemize}
  \item 截至2026年6月，代币化现实世界资产（RWA）市场总值突破510亿美元，年初至今增长40\%，而同期加密市场整体下跌约20\%，走出独立增长曲线。
  \item 私人信贷以47\%的占比主导代币化资产类别，美国国债占30\%，商品占9\%，代币化股票虽仅占不到3\%但年初至今增长130\%（从7亿至16亿美元）。
  \item 当前代币化股票的核心缺陷是所有权与收益权分离：代币持有者不是股东登记簿上的注册所有人，不享有股息和投票权。
  \item 两条路线在博弈：交易基础设施型（快速、低成本但无股东权利）vs 结算与交易所型（区块链作为实际股票的结算层，投资者享有完整所有权）。
  \item 美国监管动作在松绑：SEC提议废除某些规则，批准了纽交所和纳斯达克的代币化证券交易试点，行业还在等'创新豁免'的明确信号。
  \item 以太坊和Provenance链承载了超过70\%的代币化资产，链上地址数已突破90万。
\end{itemize}
\begin{figure}[htbp]
\centering
\includegraphics[width=0.88\linewidth]{figures/fig_098.jpg}
\caption{EXHIBIT 1 - EXHIBIT 1: Tokenized RWA - Market Cap (\textbackslash{}\$Bn) Tokenized RWA - Market Cap (\textbackslash{}\$Bn) Private credit includes Figure tokenized credit EXHIBIT 2: Tokenized RWA market cap split by asset class Private credit includes Figure tokenized cr}
\end{figure}
\begin{figure}[htbp]
\centering
\includegraphics[width=0.88\linewidth]{figures/fig_099.jpg}
\caption{EXHIBIT 1 - EXHIBIT 3: Tokenized RWA split by blockchain network Tokenized RWA - Blockchain network split (\%)}
\end{figure}
\begin{figure}[htbp]
\centering
\includegraphics[width=0.88\linewidth]{figures/fig_100.jpg}
\caption{EXHIBIT 3 - EXHIBIT 3: Tokenized RWA split by blockchain network Tokenized RWA - Blockchain network split (\%) EXHIBIT 4: Top RWA tokenization platforms}
\end{figure}
\begin{figure}[htbp]
\centering
\includegraphics[width=0.88\linewidth]{figures/fig_101.jpg}
\caption{EXHIBIT 4 - EXHIBIT 4: Top RWA tokenization platforms EXHIBIT 5: Tokenized real world asset holders RWA Asset Holders (in '000) EXHIBIT 6: FIGR - Consumer Loan Volumes (\textbackslash{}\$Bn) FIGR - Consumer Loan Volumes (\textbackslash{}\$Bn) EXHIBIT 7: Tokenized Equitie}
\end{figure}
\begin{figure}[htbp]
\centering
\includegraphics[width=0.88\linewidth]{figures/fig_102.jpg}
\caption{EXHIBIT 5 - EXHIBIT 5: Tokenized real world asset holders RWA Asset Holders (in '000) EXHIBIT 6: FIGR - Consumer Loan Volumes (\textbackslash{}\$Bn) FIGR - Consumer Loan Volumes (\textbackslash{}\$Bn) EXHIBIT 7: Tokenized Equities - Monthly Transfer Volumes Tokenized Equ}
\end{figure}
\begin{figure}[htbp]
\centering
\includegraphics[width=0.88\linewidth]{figures/fig_103.jpg}
\caption{EXHIBIT 6 - EXHIBIT 6: FIGR - Consumer Loan Volumes (\textbackslash{}\$Bn) FIGR - Consumer Loan Volumes (\textbackslash{}\$Bn) EXHIBIT 7: Tokenized Equities - Monthly Transfer Volumes Tokenized Equity - Transfer Volumes (\textbackslash{}\$Bn) June'26 numbers based on runrate as of June}
\end{figure}
{\small\color{refgray}\textbf{References:} [R007] Bernstein：代币化的真正战场不在资产发行，而在交易基础设施的重新定价}

\section{风险偏好与资金流：供给瓶颈定价权重估}
\textbf{市场低估的不是AI需求，而是供给端结构性约束的定价权，MLCC、存储等瓶颈环节享有超额利润。}

\begin{itemize}
  \item JPM量化团队将当前MLCC股票上涨与2000年代EM热潮对比，认为超额利润来源不是技术壁垒，而是供应链中的'瓶颈位置'。
  \item 在EM热潮中，住友金属矿山（镍冶炼）、三井金属（超薄铜箔）、信越化学（硅片）等瓶颈股都经历过类似的超额利润阶段。
  \item 当前MLCC股票处于AI热潮的'第六阶段'（类比EM热潮），仍有上涨空间，但风险在于'瓶颈效应'向下游转移。
  \item MS上调美光盈利预期，但强调比强劲数字更重要的是供给约束的可持续性：DRAM和NAND的ASP涨幅均超预期，但新厂建设周期和HBM良率损耗限制供应响应。
  \item 全球主动投资基金对日本科技板块的配置仍处于'中性偏低'，大资金尚未充分参与，动量交易者和Beta追随者是主要买家。
  \item 中国上市公司资本配置逻辑转变：2025年现金分红、并购和研发支出均创纪录，股息支付率升至39\%，净回购对EPS仍具增厚效应。
\end{itemize}
\begin{figure}[htbp]
\centering
\includegraphics[width=0.88\linewidth]{figures/fig_130.jpg}
\caption{Figure 3 - Figure 3: Factor exposure of electrical equipment (TSE 33 sector): As of June 12}
\end{figure}
\begin{figure}[htbp]
\centering
\includegraphics[width=0.88\linewidth]{figures/fig_131.jpg}
\caption{Figure 10 - Figure 3: Factor exposure of electrical equipment (TSE 33 sector): As of June 12 Figure 4: Factor exposure of two MLCC-related companies (Murata Manufacturing and Taiyo Yuden): As of June 12}
\end{figure}
\begin{figure}[htbp]
\centering
\includegraphics[width=0.88\linewidth]{figures/fig_132.jpg}
\caption{Figure 2 - Figure 4: Factor exposure of two MLCC-related companies (Murata Manufacturing and Taiyo Yuden): As of June 12 Figure 5: Estimated net position and profit margin momentum of Cash Equity Trend-Followers for two MLCC-related compan}
\end{figure}
\begin{figure}[htbp]
\centering
\includegraphics[width=0.88\linewidth]{figures/fig_133.jpg}
\caption{Figure 3 - Figure 6: Estimated net position and profit margin momentum of Cash Equity Trend-Followers for two MLCC-related companies: Taiyo Yuden}
\end{figure}
\begin{figure}[htbp]
\centering
\includegraphics[width=0.88\linewidth]{figures/fig_134.jpg}
\caption{Figure 4 - Figure 6: Estimated net position and profit margin momentum of Cash Equity Trend-Followers for two MLCC-related companies: Taiyo Yuden Figure 7: Regional Allocation of global active investment trusts (DM+EM) (mainland China stoc}
\end{figure}
\begin{figure}[htbp]
\centering
\includegraphics[width=0.88\linewidth]{figures/fig_135.jpg}
\caption{Figure 5 - Figure 7: Regional Allocation of global active investment trusts (DM+EM) (mainland China stocks + Hong Kong stocks) and US import ratios from China Figure 8: Relative differences in regional allocation of global active investmen}
\end{figure}
\begin{figure}[htbp]
\centering
\includegraphics[width=0.88\linewidth]{figures/fig_136.jpg}
\caption{Figure 6 - Figure 9: Regional allocation of global active investment trusts (DM+EM): As of end-April}
\end{figure}
\begin{figure}[htbp]
\centering
\includegraphics[width=0.88\linewidth]{figures/fig_137.jpg}
\caption{Figure 7 - Figure 9: Regional allocation of global active investment trusts (DM+EM): As of end-April Figure 10: Regional allocation of global active investment trusts (DM+EM) – Trends for Japan, South Korea, and Taiwan: As of end-April}
\end{figure}
\begin{figure}[htbp]
\centering
\includegraphics[width=0.88\linewidth]{figures/fig_138.jpg}
\caption{Figure 8 - Figure 9: Regional allocation of global active investment trusts (DM+EM): As of end-April Figure 10: Regional allocation of global active investment trusts (DM+EM) – Trends for Japan, South Korea, and Taiwan: As of end-April F}
\end{figure}
\begin{figure}[htbp]
\centering
\includegraphics[width=0.88\linewidth]{figures/fig_139.jpg}
\caption{Figure 9 - Figure 12: Changes in sector weightings of Japanese stocks in global active investment trusts (DM) (Active vs. Passive): As of end-April}
\end{figure}
\begin{figure}[htbp]
\centering
\includegraphics[width=0.88\linewidth]{figures/fig_140.jpg}
\caption{Figure 10 - Figure 12: Changes in sector weightings of Japanese stocks in global active investment trusts (DM) (Active vs. Passive): As of end-April}
\end{figure}
\begin{figure}[htbp]
\centering
\includegraphics[width=0.88\linewidth]{figures/fig_141.jpg}
\caption{Figure 11 - Figure 11: Sector weightings of Japanese stocks in global active investment trusts (DM) (Active vs. Passive): As of end-April Figure 12: Changes in sector weightings of Japanese stocks in global active investment trusts (DM) (Act}
\end{figure}
\begin{figure}[htbp]
\centering
\includegraphics[width=0.88\linewidth]{figures/fig_207.jpg}
\caption{Exhibit 3 - Exhibit 3: Weekly stock performance}
\end{figure}
\begin{figure}[htbp]
\centering
\includegraphics[width=0.88\linewidth]{figures/fig_208.jpg}
\caption{Exhibit 3 - Exhibit 3: Weekly stock performance Exhibit 4: Monthly stock performance Exhibit 5: YTD stock performance}
\end{figure}
\begin{figure}[htbp]
\centering
\includegraphics[width=0.88\linewidth]{figures/fig_209.jpg}
\caption{Exhibit 3 - Exhibit 3: Weekly stock performance Exhibit 4: Monthly stock performance Exhibit 5: YTD stock performance \#\# Semis Week Ahead}
\end{figure}
\begin{figure}[htbp]
\centering
\includegraphics[width=0.88\linewidth]{figures/fig_210.jpg}
\caption{Exhibit 4 - Exhibit 4: Monthly stock performance Exhibit 5: YTD stock performance \#\# Semis Week Ahead Exhibit 6: June 22nd - June 26th, 2026 Monday, June 15th}
\end{figure}
\begin{figure}[htbp]
\centering
\includegraphics[width=0.88\linewidth]{figures/fig_211.jpg}
\caption{Exhibit 8 - Exhibit 8: MS vs. Street Revenue and EPS Growth Exhibit 9: 2025-27e Revenue CAGR Exhibit 10: 2025-27e EPS CAGR \#\# Inventory}
\end{figure}
\begin{figure}[htbp]
\centering
\includegraphics[width=0.88\linewidth]{figures/fig_212.jpg}
\caption{Exhibit 9 - Exhibit 9: 2025-27e Revenue CAGR Exhibit 10: 2025-27e EPS CAGR \#\# Inventory Exhibit 11: Semiconductor Company inventory is at 111 days, up 3 days q/q which is 4 days behind of the seasonal decrease of 1 days. DOI is 22 days abo}
\end{figure}
\begin{figure}[htbp]
\centering
\includegraphics[width=0.88\linewidth]{figures/fig_213.jpg}
\caption{Exhibit 11 - Exhibit 12: Semi Customer DOI decreased 6 days sequentially to 52 days, below the seasonal decrease of 3 days. DOI is below the historical median and trailing 4 quarters increased from last quarter. Semi Customers (DOI)}
\end{figure}
\begin{figure}[htbp]
\centering
\includegraphics[width=0.88\linewidth]{figures/fig_214.jpg}
\caption{Exhibit 12 - Exhibit 12: Semi Customer DOI decreased 6 days sequentially to 52 days, below the seasonal decrease of 3 days. DOI is below the historical median and trailing 4 quarters increased from last quarter. Semi Customers (DOI) Exhibit 13}
\end{figure}
\begin{figure}[htbp]
\centering
\includegraphics[width=0.88\linewidth]{figures/fig_215.jpg}
\caption{Exhibit 12 - Exhibit 14: Short Interest as \% of Float}
\end{figure}
{\small\color{refgray}\textbf{References:} [R010] JPM：市场低估了“供给瓶颈”在日本MLCC股中的定价权，而非AI需求本身; [R017] 摩根斯坦利：市场低估的不是AI需求，而是供给端结构性约束的定价权; [R012] GS：A股硬科技正在主导中国盈利增长的结构性分化}

\section{结语}
综合来看，当前市场主线围绕AI硬件与供给瓶颈的定价权展开，同时中国内需磨底与家庭资产负债表再平衡构成另一条暗线。区域市场分化显著，A股硬科技与日本零售整合值得关注，而欧洲化工与印度AI则面临结构性挑战。代币化与储能产业正处于基础设施竞争的关键拐点。
\newpage
\section{报告覆盖清单}
\begin{itemize}
  \item [R001] 宏观与利率：中国内需磨底，家庭资产负债表再平衡 / 能源与大宗：霍尔木兹冲击重塑油需，AI压缩深水油田周期 - GS：市场低估了霍尔木兹冲击对油需的结构性重塑
  \item [R002] 全球游戏业：GTA VI前的洗牌前夜 - Bernstein：全球游戏业真正的拐点不在2026年，而在GTA VI之后的2027年
  \item [R003] AI/算力：云巨头资本效率竞赛与内存成本冲击 - Bernstein：云巨头正在从“算力军备竞赛”转向“资本效率竞赛”
  \item [R004] 能源与大宗：霍尔木兹冲击重塑油需，AI压缩深水油田周期 - GS：AI正在重塑全球深水油田的经济模型，但市场低估了“时间压缩”的定价权
  \item [R005] 储能与新能源：供应链本地化重构拐点 - Bernstein：全球储能市场的真正拐点不在技术突破，而在供应链的本地化重构
  \item [R006] 半导体设备：DRAM与NAND的结构性分化 - 摩根斯坦利：半导体设备市场的真正变量不是总量，而是DRAM与NAND的结构性分化
  \item [R007] 代币化与数字资产：基础设施竞争决定胜负 - Bernstein：代币化的真正战场不在资产发行，而在交易基础设施的重新定价
  \item [R008] AI/算力：云巨头资本效率竞赛与内存成本冲击 - Bernstein：AI正被内存成本压到喘不过气，一场供应链“再校准”不可避免
  \item [R009] 区域市场：A股硬科技领跑，印度AI自主性困境 - Bernstein：印度真正需要的不是算力租用，而是自己的“DeepSeek 时刻”
  \item [R010] 区域市场：A股硬科技领跑，印度AI自主性困境 / 风险偏好与资金流：供给瓶颈定价权重估 - JPM：市场低估了“供给瓶颈”在日本MLCC股中的定价权，而非AI需求本身
  \item [R011] 地产与消费：中国内需磨底，日本零售整合拐点 - GS：家电内需压力尚未出清，但出口的结构性机会比想象中更值得关注
  \item [R012] 区域市场：A股硬科技领跑，印度AI自主性困境 / 风险偏好与资金流：供给瓶颈定价权重估 - GS：A股硬科技正在主导中国盈利增长的结构性分化
  \item [R013] 储能与新能源：供应链本地化重构拐点 - Bernstein：数据中心供电的真正瓶颈不在技术，在电网的“连接组织”
  \item [R014] 地产与消费：中国内需磨底，日本零售整合拐点 - GS：欧盟暂缓对中国割草机器人征收临时反倾销税，真正的变量不是关税本身，而是企业能否利用这个窗口期重塑竞争壁垒
  \item [R015] 宏观与利率：中国内需磨底，家庭资产负债表再平衡 / 地产与消费：中国内需磨底，日本零售整合拐点 - GS：中国家庭资产正在经历从房产到金融资产的结构性迁移
  \item [R016] 能源与大宗：霍尔木兹冲击重塑油需，AI压缩深水油田周期 - GS：欧洲化工的真正危机不是红海冲突，而是冲突结束后暴露的结构性过剩
  \item [R017] 半导体设备：DRAM与NAND的结构性分化 / 风险偏好与资金流：供给瓶颈定价权重估 - 摩根斯坦利：市场低估的不是AI需求，而是供给端结构性约束的定价权
  \item [R018] 宏观与利率：中国内需磨底，家庭资产负债表再平衡 - NOM：人民币中间价模型正在传递一个比市场预期更克制的信号
  \item [R019] 地产与消费：中国内需磨底，日本零售整合拐点 - Bernstein：日本零售业真正被低估的不是消费复苏，而是供给侧的整合拐点
  \item [R020] AI/算力：云巨头资本效率竞赛与内存成本冲击 - GS：软件行业的真正分水岭不是AI本身，而是“好粘性”与“坏粘性”的重新定价
\end{itemize}
\section{逐篇报告摘录}
\subsection{[R001] GS：市场低估了霍尔木兹冲击对油需的结构性重塑}
{\small\color{refgray}归类：宏观与利率：中国内需磨底，家庭资产负债表再平衡 / 能源与大宗：霍尔木兹冲击重塑油需，AI压缩深水油田周期}

霍尔木兹海峡的供应冲击正在改变全球石油市场的需求曲线，而不仅仅是供给曲线。这是GS最新研报中最值得关注的判断。市场习惯性地将地缘政治风险简化为短期价格脉冲和库存扰动，但这份报告揭示了一个更深层的机制：当油价因供给冲击而上涨时，消费者行为并非被动接受，而是加速了结构性替代。这份报告的独特价值在于，它用实时数据捕捉了这种替代正在发生，并且量化了它对2027年油需的潜在影响。

报告的核心发现是：自今年2月霍尔木兹冲击开始以来，全球电动汽车渗透率已上升3.4个百分点，达到26.1\%的历史新高。这不是一个温和的线性外推，而是一个加速信号。GS据此估算，到2027年12月，全球石油需求可能因此面临每日13万至32万桶的下行风险。这个数字在每日约1亿桶的全球市场中看似微小，但它的意义不在于绝对值，而在于它揭示了需求侧正在发生的结构性变化——而市场对此的定价可能远远不足。

理解这份报告，需要跳出“电动汽车替代石油”的陈旧叙事。真正重要的是三个层次的问题：第一，替代的速度为何在此时加速；第二，这种加速是否具有持续性；第三，当前的市场定价是否已经充分反映了这一变化。

GS使用的“nowcast”模型显示，全球电动汽车渗透率并非均匀上升，而是在霍尔木兹冲击发生后出现了明显的斜率变化。从2月到5月，全球渗透率从约22.7\%跃升至26.1\%，这一增幅超过了此前12个月的总和。更值得注意的是，在15个最大的电动汽车市场中，有12个出现了渗透率上升，这是一个罕见的广泛性信号。

中国是这一变化的绝对主力。中国的电动汽车渗透率在同期内上升了11.4个百分点，贡献了全球增量的61\%。OECD国家贡献了21\%，非OECD非中国地区贡献了19\%。这种分布结构说明，加速并非由单一政策驱动，而是呈现全球性特征。即使考虑到韩国因补贴政策提前而出现的短期脉冲，整体趋势依然清晰。

这里的关键洞察是：供给冲击正在改变消费者的时间偏好。当高油价成为现实，消费者从“未来可能买电动车”转变为“现在就买电动车”的决策加速。这种时间偏好的改变，一旦形成规模，就会产生路径依赖——充电基础设施的扩张、电池成本的下降、产品选择的丰富，都会进一步强化这一趋势。

GS设定了两个场景来量化影响。“暂时性加速”场景假设区域渗透率维持在2026年5月的水平不变，测算出的需求冲击为每日13万桶。“持续性加速”场景假设渗透率按照2月至5月的趋势线性增长，测算出的冲击为每日32万桶。

两个场景的差异不在于幅度，而在于持续性。这正是市场最容易低估的部分。如果渗透率只是脉冲式上升后维持高位，影响是有限的；但如果加速趋势得以延续，那么到2027年底，累计的存量替代将产生可观的边际效应。

GS使用的经验法则值得注意：每100万辆从燃油车转向电动车的销售转移，在美国会降低道路石油需求每日3万桶，在其他地区为每日2万桶。这一规则暗示，中国以外的市场单位替代效率更高，因为中国的电动车平均能效可能与美国存在差异，或者因为中国道路燃料的碳氢化合物组成不同。这背后涉及的细节报告没有完全展开，但对于理解区域差异至关重要。

投资者需要问自己的问题是：当前的市场定价中，隐含的是哪个场景？如果市场已经定价了“暂时性加速”，而实际走势更接近“持续性加速”，那么当前的石油远期曲线可能系统性高估了2027年的需求。反之，如果市场已经在定价“持续性加速”，那么GS的测算可能只是确认了已知信息。

GS明确指出了两个未纳入分析的重要维度，而这恰恰是理解完整影响的关键。

第一，报告只考虑了新增销量对存量的贡献，但没有充分量化存量中的替代行为。GS提到，中国汽油及相关产品销量同比下降超过20\%，同时电动车充电量激增，这证明存量替代正在发生。但报告没有给出存量替代的量化模型。这意味着，GS测算的每日13万至32万桶可

[... middle omitted ...]

被动等待下一份研报的更新。

第一个维度是渗透率的斜率变化。GS的nowcast提供了月度数据，但投资者可以跟踪更高频的指标，例如中国每周的电动汽车上险数据、欧洲主要市场的月度注册数据。斜率的变化比绝对水平更重要。如果渗透率增速从当前水平进一步加快，那么GS的“持续性加速”场景可能仍然偏保守。

第二个维度是政策反馈循环。霍尔木兹冲击导致的高油价，正在改变各国政府的政策权衡。高油价增加了消费者对电动车补贴的接受度，也降低了政府补贴的财政压力（因为燃油税收入增加）。这种正反馈循环可能在未来6至12个月催生更多支持性政策。报告没有深入分析这一点，但这是逻辑的自然延伸。

第三个维度是供给侧的响应。GS在报告中提到，其下行场景假设布伦特原油在2027年底跌至50美元/桶中段，部分原因是供给端的超预期增长。如果需求侧因电动车加速而出现结构性下行，那么OPEC+的减产策略将面临更大挑战。供给端的竞争性增产与需求端的替代加速，可能形成双重压力。这是报告尚未完全展开的“二阶效应”。

\subsection{[R002] Bernstein：全球游戏业真正的拐点不在2026年，而在GTA VI之后的2027年}
{\small\color{refgray}归类：全球游戏业：GTA VI前的洗牌前夜}

Bernstein：全球游戏业真正的拐点不在2026年，而在GTA VI之后的2027年

全球视频游戏行业在2025年实现了2180亿美元的收入，同比增长4.8\%。这个数字本身并不令人意外——行业早已从疫情后的回调中恢复。真正值得关注的信号，是Bernstein这份年度深度报告揭示的一个结构性转折：行业正在从“增长驱动”转向“整合驱动”，而2026年可能成为近年来最被低估的过渡年份。

报告的核心判断并不复杂：全球游戏行业的收入增长将在2026年放缓至0.7\%，然后在2027年回升至3.3\%。但数字背后的叙事远比表面丰富。这份报告真正想告诉投资者的，不是增长率的短期波动，而是行业竞争格局正在经历一次自移动游戏爆发以来最深刻的洗牌。

我们整理了Bernstein这份报告的五个核心洞察，以及一个报告尚未完全回答的关键问题。

Bernstein在报告中反复强调一个概念：行业正处于“资本周期的后半段”。这意味着，过去几年大量资本涌入、新工作室遍地开花的阶段已经结束，取而代之的是市场集中度的提升和落后产能的出清。

报告明确指出，虽然整体行业收入增速放缓，但上市公司的收入增速预计将达到7-8\%，显著高于行业平均水平。这背后有两个驱动力：一是西方市场长期被推迟的工作室关闭潮终于到来，二是头部公司凭借IP和资金优势在竞争中持续扩大份额。

这一判断的意义在于，投资者不应再将注意力集中在行业总规模的增长上，而应关注哪些公司能够在整合中成为赢家。Bernstein用一组数据说明这个逻辑：2025年Steam上发布的新游戏数量从2019年的8100款激增至21400款，但获得超过5000条评论的游戏数量却从73款降到了72款。供应的爆炸式增长并没有带来成功产品的相应增加——这恰恰是资本周期后半段的典型特征。

2. GTA VI的发布将成为一个“清算事件”，而非单纯的利好

Bernstein对2026年最引人注目的判断，是关于《侠盗猎车手VI》（GTA VI）的。报告认为，这款游戏在2026年秋季的发布，可能成为整个行业投资者情绪和仓位调整的转折点。

报告通过夏季游戏展的观察发现，整个行业为了避免与GTA VI正面竞争，已经将大量新游戏的发布窗口挤到了9月。这种拥挤带来了明显的短期风险：对于那些没有强大开发商或IP信誉支撑的游戏，销售不及预期的概率显著上升。

但Bernstein的洞察不止于此。报告指出，目前投资者对GTA VI的预期已经“从看涨到圣经级别”。这种极端预期意味着，无论GTA VI的实际表现如何，它都可能成为其他游戏公司股票重新定价的催化剂。如果GTA VI大获成功，它会带动整个行业的关注度和参与度；但如果它只是符合预期甚至略有不及，那些因避让而拥挤在9月的游戏可能会面临更大的压力。

对于投资者来说，2026年11月前后可能是一个关键的观察窗口——不是看GTA VI卖了多少份，而是看其他游戏公司在那之后能否获得新的增长动能。

3. PC正在取代主机成为增长最快的平台，驱动力来自中国和硬件升级

报告中最具颠覆性的判断之一，是PC正在取代主机成为行业增长最快的细分市场。2025年，PC游戏收入同比增长8.4\%，远超此前预期的3\%，而主机游戏增速为9.8\%，主要受Switch 2发布推动。

更重要的是，Bernstein认为PC的增长有结构性支撑。报告指出，随着普通家用PC的GPU显存（VRAM）达到8-12GB，足以支持PS4/PS5时代的游戏体验，PC游戏的潜在用户群正在急剧扩大。同时，主机和移动端的开发者越来越多地采用多平台发布策略，进一步推动了PC市场的增长。

中国市场的表现尤其突出。2025年中国PC游戏市场增长了15\%，《黑神话：悟空》的成功将大量玩家从移动端吸引到了大屏幕。Bernste

[... middle omitted ...]

以及多司法管辖区的诉讼持续进行，平台费率的长期方向仍然是下行。

这对于内容创作者而言无疑是利好。但报告同时提出了一个重要的限定条件：Bernstein的基本假设是，发行和二次合作关系在可预见的未来不会发生实质性变化。这意味着，平台费率的下调主要惠及的是现有体系内的参与者，而非改变行业权力结构。

换句话说，费率下调是增量利好，但不是结构性变革。真正决定内容创作者命运的，仍然是IP实力和社区可见度——这与报告在其他部分反复强调的观点一致。

Bernstein对AI在游戏行业的影响给出了一个明确的判断：恐惧是被误导的。报告承认，行业确实在寻找利用AI加速开发的方法，但IP实力和社区可见度仍然是新游戏成功的关键驱动因素。

报告用一组数据支撑这个观点：尽管Steam上新游戏的数量从2019年的8100款增长到2025年的21400款，但获得超过5000条评论的游戏数量仅从73款变为72款。这意味着，AI虽然降低了游戏开发的门槛，但并没有改变“成功产品稀缺”这一核心事实。

\subsection{[R003] Bernstein：云巨头正在从“算力军备竞赛”转向“资本效率竞赛”}
{\small\color{refgray}归类：AI/算力：云巨头资本效率竞赛与内存成本冲击}

Bernstein：云巨头正在从“算力军备竞赛”转向“资本效率竞赛”

市场对AI基础设施投资的焦虑从未停止。每一季度云巨头的资本支出数字都在刷新纪录，而关于“何时能看到回报”的质疑声浪也随之升高。但Bernstein最新发布的这份涵盖全球软件、美国和中国互联网以及通信基础设施的季度云报告，给出了一个值得深思的信号：真正决定未来几年赢家的，可能不再是谁能建起最大的GPU集群，而是谁能以最高的资本效率将算力转化为可交付的云收入。

这份由Bernstein全球云团队联合撰写的报告，追踪了亚马逊AWS、微软Azure、谷歌云、甲骨文OCI、阿里云以及新纳入的“Neocloud”代表CoreWeave的季度表现。报告覆盖的周期已延伸至2026年第一季度，并基于长达21个季度的历史数据，构建了一个观察超大规模云市场结构性变化的参照系。其核心判断是：市场正在从单纯的“建设能力”竞争，转向“运营效率”与“资本结构”的复合博弈。

1. 资本支出与运营现金流的剪刀差扩大，融资能力成为新护城河

报告中最具冲击力的数字之一，是四大云巨头（亚马逊、微软、谷歌、甲骨文）在2026年第一季度的现金资本支出总和已接近1300亿美元。Bernstein模型显示，2026年全年这一数字将达到约6230亿美元，若加上Meta，则逼近7600亿美元。与此同时，这四家公司全年的经营活动现金流估计约为6350亿美元。

这组数据揭示了一个关键矛盾：资本支出的增长速度正在持续超过运营现金流的增长。报告中的图表清晰展示了这一趋势——当资本支出曲线陡峭上行时，现金流增长曲线却呈现波动盘整。这意味着，云巨头们仅靠内部造血已难以支撑其扩张速度。事实上，亚马逊、谷歌和甲骨文近期已陆续进入资本市场进行融资。

这一变化的含义是深远的。在上一轮云基础设施周期中，现金充裕的巨头几乎不需要外部融资。但现在，能否以合理成本获取长期资本，将直接影响企业的扩张速度和投资回报率。对于投资者而言，观察重点不应仅仅是资本支出的绝对规模，更应关注其融资结构、债务成本以及自由现金流的覆盖倍数。

对于“回报在哪里”的追问，Bernstein给出了两个维度的回答。第一个维度是实际收入增长。报告显示，AWS、Azure、GCP和OCI的合计云收入增速在最近几个季度确实有所加快，但报告明确指出，这一增速“仍需继续提升”，才能与资本支出的增速相匹配。

第二个维度是合同积压，即剩余履约义务。这一指标在最近几个季度出现了显著跃升，截至2026年第一季度，四大云巨头的合计RPO已突破2万亿美元。谷歌云的RPO更是接近环比翻倍，达到约4620亿美元，其中超过50\%预计将在未来24个月内转化为收入。甲骨文的RPO则高达6380亿美元，其中约3000亿美元来自OpenAI的合同。

RPO的高速增长是一个强烈的乐观信号，因为它代表了未来收入的可见性。但报告也提醒，这一指标存在两个“混杂因素”：合同期限的延长和客户集中度。如果大客户签署的是多年期大额合同，RPO的增长可能夸大了短期收入的爆发力。因此，投资者需要进一步拆解RPO的期限结构、客户集中度以及收入确认节奏，才能判断其真正的含金量。

在各大云厂商中，Bernstein对微软的定位最为清晰。报告指出，微软的AI业务年化收入已超过370亿美元，AI已成为“有意义的收入贡献者，而非负担”。更关键的是，微软拥有一个独特的“双轮驱动”结构：一方面，其自有的第一方应用（如Office 365 Copilot、GitHub Copilot）正在快速消耗算力；另一方面，其第三方客户合同（不仅限于OpenAI，还包括大量企业客户）提供了多元化的收入来源。

这种结构的战略价值在于：如果AI投资真的出现泡沫破裂，微软可以通过调整第一方应用的算力分配来保持利用率，

[... middle omitted ...]

软在新增云收入季度环比增长上“并驾齐驱”。

Bernstein特别强调了谷歌云一个容易被忽视的增长引擎：向客户直接销售TPU芯片。这种“硬件即服务”的模式，允许谷歌将其定制AI芯片部署到客户自己的数据中心，从而开辟了一条新的收入路径。虽然这部分收入预计在2026年后期才开始体现，2027年才会更加显著，但它显著提升了谷歌云的长期收入可见性。

不过，谷歌云的市场评级仅为“市场表现”，目标价390美元，低于亚马逊和微软的“跑赢大盘”评级。报告隐含的判断是：谷歌云的利润率改善（1Q26达到约33\%，同比提升超过15个百分点）虽然令人瞩目，但折旧增加和Wiz收购的拖累将在2026年形成压力，其盈利质量仍需持续验证。

5. 报告尚未完全回答的关键问题：资本效率的“临界点”何时到来？

尽管报告提供了丰富的分析框架，但有一个核心问题并未完全展开：资本效率的“临界点”在哪里？换言之，资本支出增速需要放缓到何种水平，或者收入增速需要加速到何种程度，才能让自由现金流开始显著改善？

\subsection{[R004] GS：AI正在重塑全球深水油田的经济模型，但市场低估了“时间压缩”的定价权}
{\small\color{refgray}归类：能源与大宗：霍尔木兹冲击重塑油需，AI压缩深水油田周期}

GS：AI正在重塑全球深水油田的经济模型，但市场低估了“时间压缩”的定价权

一份来自GS的最新报告，以15家以上国际石油公司、油服企业和数字供应商的访谈为基础，试图回答一个对全球能源供给格局至关重要的问题：AI和数字化到底能多大程度改变一个深水油田从发现到产油的时间表，以及这个时间表的变化会如何改写上游项目的投资回报率。

这份报告的核心判断是：AI对油气行业的冲击不是渐进式的效率改善，而是结构性压缩了项目从勘探到最终投资决策(FID)的周期，进而在成本曲线上制造了一个足以影响长期油价预期的位移。GS估算，一个典型的深水绿地项目，全周期时间可从12年压缩至7年，降幅约38\%。而由此带来的内部收益率(IRR)提升可达3.5个百分点，盈亏平衡价格下降约15\%。

这些数字本身已经足够引人注目。但真正值得产业决策者和投资者关注的，不是AI“能否”带来改变，而是“在哪里”改变、改变“由谁捕获”，以及哪些传统假设正在被颠覆。

GS的报告给出了一个反直觉的结论：AI带来的时间节省几乎全部发生在项目最终投资决策(FID)之前。具体来看，勘探阶段可缩短55\%，评价阶段缩短40-50\%，前端工程设计(FEED)同样缩短40-50\%。这三个阶段合计贡献了全部时间节省的90\%。

这意味着什么？意味着AI正在改写的是油气项目“决策前的确定性”，而不是“建设中的速度”。在传统模式下，从发现油田到决定是否投资，往往需要5-8年的数据收集、地震解释、井位设计和工程验证。AI通过加速地震数据处理、自动解释和工程工作流，将这一周期大幅压缩。GS访谈中，有运营商报告地震解释速度提升了5-10倍，在试点盆地中，预测成功率从约50\%提升至约75\%。

但FID之后，物理世界的刚性开始显现。建设阶段仅能压缩约5-10\%，而FPSO(浮式生产储卸装置)的建造周期仍然固定在约4年，受制于船厂产能、钢材供应和长周期设备。这一阶段的瓶颈不是数字技术，而是物理制造。

这个分布特征对投资者的含义是：AI带来的时间价值，主要被那些能够更快做出FID决策的公司捕获，而不是那些擅长建设执行的公司。对于油服企业而言，如果业务集中在建设阶段，AI的增量收益可能相当有限。

2. IRR提升的四个杠杆中，资本支出和“时间价值”贡献了四分之三

GS将AI对上游经济学的改善分解为四个杠杆：资本支出减少(贡献+1.9个百分点IRR)、时间压缩(+0.8个百分点)、产量提升(+0.5个百分点)和运营支出下降(+0.3个百分点)。合计3.5个百分点的IRR提升，将典型绿地项目的回报率从15.5\%推高至19.0\%。

资本支出减少是最大的单一贡献者。AI加速钻井设计和执行，运营商报告钻井时间减少20-25\%，工程设计周期从数月压缩至数周。这直接降低了单井成本，从而降低了整个项目的资本支出强度。

时间压缩的贡献仅次于资本支出。但这里有一个容易被忽略的细节：GS假设的时间压缩仅为3个月，而非全周期压缩的5年。原因在于，IRR计算中，时间价值的影响取决于现金流的前移程度，而FID之前的阶段虽然大幅缩短，但这一阶段本身并不产生现金流。真正影响IRR的是从FID到首次产油的时间缩短。GS估算，这一阶段仅能压缩约3个月。

这意味着，全周期时间压缩38\%的宏大叙事，在财务模型中的直接体现远比直觉上温和。但反过来看，如果AI能够进一步突破FID后的物理瓶颈——比如通过数字孪生优化FPSO建造排程，或者通过模块化设计缩短集成时间——那么IRR的提升空间可能远超当前估算。

GS将AI驱动的各类杠杆分为三个层级：已验证、新兴和愿景型。这个分类本身就是一个投资框架。

已验证的杠杆包括：更快的钻井、地震加速、生产效率提升、上部设施瓶颈消除、气举优化。这些技术已经在多个运营商中产生可衡量的结果，构

[... middle omitted ...]

率提高2-5个百分点——被GS认为是一个5-10年的故事，目前未纳入分析。但值得指出的是，如果这一杠杆最终被证明可行，它对全球油气供给曲线的影响将是量级性的。

这个分层框架对投资者的含义是：当前市场定价中，可能已经反映了已验证杠杆的价值，但尚未充分反映新兴杠杆的期权价值，也几乎完全没有定价愿景型杠杆的潜在冲击。

4. 成本曲线的位移：75分位盈亏平衡从75美元降至64美元

GS将上述AI场景输入其2026年“Top Projects”成本曲线，发现75分位的盈亏平衡价格将从每桶75美元降至64美元(以13-18\%的加权平均资本成本计算)。这个位移的驱动力主要来自资本支出和时间压缩，而非运营支出。

这个数字的意义在于，它直接挑战了市场对长期油价支撑位的共识。如果AI能够系统性地将深水项目的盈亏平衡价格降低约15\%，那么一批原本处于成本曲线高位的项目将变得可行，而一批原本处于边际的项目将被加速推进。这可能导致全球油气供给曲线的右移，进而对长期油价预期形成下行压力。

\subsection{[R005] Bernstein：全球储能市场的真正拐点不在技术突破，而在供应链的本地化重构}
{\small\color{refgray}归类：储能与新能源：供应链本地化重构拐点}

Bernstein：全球储能市场的真正拐点不在技术突破，而在供应链的本地化重构

全球储能产业正在经历一个被大多数投资者低估的结构性转变。这份Bernstein最新发布的全球储能周报，表面上是一系列分散的行业动态汇总，但当我们把这些信号拼合在一起，一个清晰的判断浮出水面：市场关注的焦点正在从“谁的技术更先进”转向“谁的供应链更能在本土落地”。这不是一个渐进式的变化，而是一个正在加速的拐点。

为什么现在这个判断重要？因为在过去两年，投资者习惯用电池化学的突破、能量密度的提升、或者固态电池的进展来评估储能公司的价值。但Bernstein这份报告揭示了一个不同的叙事：真正驱动行业格局重塑的力量，来自美国、欧洲、中国三个区域同步发生的供应链本地化运动。这不是贸易摩擦的简单延伸，而是储能产业从“全球化采购”向“区域化闭环”的模式切换。

报告中最值得关注的信号不是某个技术突破，而是几个看似分散的供应链事件：韩国HyVISION Systems拿下950亿韩元的北美ESS设备订单，Sebang与LG Energy Solution合作在俄亥俄州建设模块工厂，以及中国政策要求高速公路配套重卡充电设施。这些事件的共同指向是：储能产业的竞争维度正在从电池性能转向供应链的地理韧性。

1. 北美市场的核心驱动力不是需求，而是合规驱动的本地化制造

Bernstein报告中最密集的信号来自北美。HyVISION Systems的订单金额从初始规模扩张近五倍，直接原因是客户扩大了美国本土生产计划至2028年。Yujin Technology正在向美国、加拿大、波兰的ESS生产线扩大出货。Sebang与LGES的合作项目预计到2028年累计订单价值1.8万亿韩元。

这些事件表面上是韩国设备商的业务扩张，但更深层的含义是：北美储能市场的竞争壁垒正在从技术门槛转向合规门槛。美国FEOC（外国关注实体）要求正在重塑整个供应链的流向。SK On成为现代汽车佐治亚州Metaplant的独家电池供应商，其核心优势不是电池性能，而是其非中国供应链能够满足FEOC要求。

这意味着什么？对于投资者而言，评估一家储能公司在北美的前景，首先要问的不是它的电池能量密度有多高，而是它的供应链是否能在合规框架内完成本地化。那些能够在美国建立从电芯到系统集成的完整本地供应能力的公司，将获得超越技术指标的定价权。

2. 中国市场的政策转向正在淘汰纯硬件玩家，筛选出具备系统集成能力的公司

Bernstein报告指出，中国储能行业在2026年下半年将面临三项重大政策转变。第114号文件、第41号令以及新的DL/T 并网标准，正在将市场从政策驱动转向合规与性能驱动。容量电价、更严格的安全并网要求、以及分布式可再生能源的储能强制配套，这些措施的共同效果是：储能不再是一个可以靠低价硬件取胜的市场。

报告明确判断，行业将出现整合，那些在电网集成、合规、EMS/软件优化和资产运营方面具备能力的公司将胜出，而低成本硬件导向的玩家将面临困境。这是中国储能市场从“量的扩张”转向“质的竞争”的关键信号。

这里需要特别注意的是，中国锂盐期货在经历近一个月的下跌后反弹，市场解读为ESS、EV和AI数据中心需求的支撑。但Bernstein报告隐含的提醒是：原材料价格的短期波动不应掩盖下游竞争格局的深层变化。当政策迫使行业从“装得多”转向“用得好”，那些只靠电芯出货量讲故事的公司，其估值逻辑需要重新审视。

3. 固态电池的商业化时间表被再次拉长，短期竞争仍围绕LFP与LMR的路线选择

CATL再次确认，固态电池的大规模商业化在2030年前不太可能实现。CEO曾毓群指出，早期的固态电池将首先出现在高端电动车中，而CATL的目标是在2027年左右开始小规模生产。但报告同时强调，当前固

[... middle omitted ...]

欧洲市场的信号显示，储能正在从能源基础设施升级为AI基础设施的关键组件

Bernstein报告中的几个欧洲事件值得深入解读。Farasis Energy与德国WLF Energy合作开发AI驱动的ESS，Pylontech在斯洛伐克完成13.2MWh项目，BYD储能系统为匈牙利最大电池项目供电。这些项目本身并不惊人，但它们的共同背景是：欧洲正在将储能视为AI数据中心、虚拟电厂和智能电网的核心枢纽。

特别值得关注的是英国首相在G7峰会上宣布的13亿英镑投资承诺，其中包括InfraVia Capital的10亿英镑用于电池储能和能源项目。这一信号表明，欧洲政策制定者已经开始将储能投资与AI生态系统建设挂钩。储能不再只是可再生能源的配套，而是数字基础设施的底层支撑。

对于投资者而言，这意味着评估欧洲储能市场时，需要关注的不只是可再生能源装机目标，还要考虑AI数据中心的电力需求增长曲线。那些能够将储能系统与AI能源管理平台深度整合的公司，可能获得超出行业平均的估值溢价。

\subsection{[R006] 摩根斯坦利：半导体设备市场的真正变量不是总量，而是DRAM与NAND的结构性分化}
{\small\color{refgray}归类：半导体设备：DRAM与NAND的结构性分化}

摩根斯坦利：半导体设备市场的真正变量不是总量，而是DRAM与NAND的结构性分化

当市场还在争论2026年半导体资本支出是否会创下新高时，摩根斯坦利最新发布的半导体资本设备研报给出了一个更精细的答案：总量确实在上修，但真正值得关注的不是数字本身，而是DRAM和NAND两条主线正在走向截然不同的投资节奏。

这份报告将2026年全球DRAM设备支出（WFE）从之前的446亿美元上调至486亿美元，2027年从571亿美元上调至625亿美元。与此同时，NAND的WFE预测几乎原地踏步——2026年从159亿美元微调至157亿美元，2027年从240亿美元调整至241亿美元。乍看之下，DRAM是绝对的王者。但报告真正的判断藏在更深的层次里：2027年NAND的WFE增速将高达53\%，远超DRAM的29\%。这意味着，当前市场对设备股的定价，可能正在忽略一个即将到来的结构性切换。

这不是一份简单的“上调预测”报告。它揭示了半导体设备投资中一个容易被忽视的规律——当技术升级遇到产能扩张，谁在提前花钱，谁在等待拐点，决定了下一阶段设备商的赢家格局。

1. DRAM的上修并非需求拉动，而是客户“不惜一切代价”的提前支出

摩根斯坦利将DRAM WFE上修的核心驱动力归结为“brownfield”——即现有产线的升级改造，而非全新的晶圆厂建设。2026年绿色field（新建厂）的月产能仅从258kwpm微调至268kwpm，2027年维持395kwpm不变。真正拉动上修的是brownfield部分的假设：报告明确写道，客户正在“不惜一切代价”（doing everything in their power）将设备采购前置。

这意味着什么？DRAM设备支出的增长，不是需求端的自然扩张，而是供给端竞争加剧的结果。HBM（高带宽存储器）的爆发式需求，迫使DRAM厂商在现有产线上加速转向HBM产能。但HBM的生产效率远低于标准DRAM——报告中的Exhibit 5显示，到2028年，HBM的trade ratio将导致约27\%的bit成为“丢失的bit”。也就是说，同样数量的晶圆投入，产出的可用bit数量因HBM的复杂工艺而大幅缩水。

对设备商而言，这是个好消息：客户为了弥补HBM带来的效率损失，必须购买更多的设备。对投资者而言，这却是个需要警惕的信号：DRAM WFE的高增长，部分是在“为低效率买单”，而非技术或需求的根本性突破。一旦HBM的需求增速放缓，这种前置支出的可持续性将面临考验。

2. NAND的沉默不是停滞，而是一场良率驱动的结构性爆发前夜

与DRAM的高调上修形成鲜明对比，NAND的WFE预测几乎纹丝不动。但报告在NAND部分埋下了一个更重要的伏笔：bit供给增长的上修，完全来自良率提升，而非设备投入增加。2026年NAND bit供给增长从21\%上调至24\%，2027年从32\%微调至31\%，而WFE几乎不变。这意味着，在现有设备基础上，NAND厂商正在通过工艺优化释放出更多的有效产出。

报告坦言，“NAND WFE的加速不可避免，早期迹象已经出现，但我们不认为这一论点在短期内会兑现。”这句话是整份报告中最值得细读的判断之一。它暗示了一个时间差：当前NAND设备支出处于低位，但一旦良率提升达到瓶颈，或eSSD（企业级固态硬盘）需求真正爆发，NAND厂商将不得不开启新一轮大规模设备采购。报告预测eSSD bit在 年的三年复合增长率高达66\%，eSSD在NAND总产出中的占比将从2025年的29\%跃升至2028年的67\%。这个结构性转变，才是NAND WFE爆发的真正引擎。

对设备商而言，NAND的拐点可能比市场预期的更晚，但一旦启动，弹性远大于DRAM。这正是摩根斯坦利在报告中明确表达“相对偏好LAM（

[... middle omitted ...]

AND。

这个判断对当前的市场定价构成挑战。AMAT作为DRAM设备的最大受益者，在2026年的WFE上修中无疑是赢家。报告甚至承认，“2026年更高的DRAM WFE将不可避免地使AMAT受益最大，我们对AMAT的降级（从Overweight降至Equal-weight）可能时机略早。”但报告的核心逻辑是前瞻性的：如果2027年NAND增速远超DRAM，那么当前AMAT的估值溢价就缺乏持续支撑。而LAM在NAND领域的暴露度更高，将在下一阶段获得更强的业绩弹性。

这不是一个简单的“买谁卖谁”的建议，而是一个关于时间窗口的判断。当前市场可能正在为DRAM的短期繁荣定价，却忽略了NAND的结构性拐点即将到来。对于持有设备股的投资者而言，关键问题不是“DRAM好不好”，而是“你愿意为哪个时间段的增长买单”。

这份报告最令人印象深刻的部分，是它坦诚地揭示了NAND bit增长主要来自良率改善，而非设备投入。但这也引出了一个报告没有充分展开的问题：良率提升的极限在哪里？

\subsection{[R007] Bernstein：代币化的真正战场不在资产发行，而在交易基础设施的重新定价}
{\small\color{refgray}归类：代币化与数字资产：基础设施竞争决定胜负}

Bernstein：代币化的真正战场不在资产发行，而在交易基础设施的重新定价

截至2026年6月，代币化现实世界资产的市场总值已突破510亿美元，年初至今增长40\%。同一时期，更广泛的加密市场下跌约20\%。这个数字本身就包含一个反直觉的判断：代币化并非加密行情的附属品，它正在走出一条独立的增长曲线。

但更值得关注的不是总量，而是结构。私人信贷以47\%的占比主导代币化资产类别，美国国债占30\%，商品占9\%。代币化股票虽然仅占不到3\%，但年初至今增长了130\%，从7亿美元跃升至16亿美元。Bernstein这份报告的核心判断是：代币化正在从“发行试验”进入“基础设施竞争”阶段，而真正决定胜负的，不是谁能发行更多资产，而是谁能构建合规、高效、可规模化的交易和结算层。

这不是一个关于“区块链能否替代传统金融”的宏大叙事。这是一个关于“在现有监管框架下，谁有能力重新定义证券的清算、交易和持有方式”的具体商业问题。

理解这一轮基础设施竞争的关键，在于看清当前主流代币化股票模型的根本局限。

目前，Robinhood等平台向非美国投资者提供的代币化股票，采用的是第三方赞助商模式。赞助商购买底层股票、托管在受监管账户中，然后发行一个区块链代币。投资者可以24x7交易这些代币，实现实时结算，平台赚取交易佣金。

但这些代币只运行在区块链上，与实际的股东登记簿没有任何对账机制。代币持有者不是股东登记簿上的注册所有人，赞助商才是。这意味着，代币持有者不享有股息、投票权等与股票所有权相关的权利。

这个模型的问题不在于技术，而在于法律结构。它提供的是“股票的价格敞口”，而不是“股票的完整所有权”。从投资角度看，这相当于一个合成工具。从监管角度看，它只能在现有豁免框架下运营，且对大多数美国投资者关闭。

Bernstein报告指出，市场真正等待的是一种“创新豁免”，允许代币化美国股票在岸交易。但在这之前，行业玩家已经开始布局更彻底的方案。

与第三方赞助商模型形成对比的是发行方赞助模型。在这个框架下，区块链作为公司实际发行股票的结算层。股东获得与传统交易所上市股票相同的所有权和保护权利，同时享受更快的结算周期和更高的资本效率。

Figure、Bullish和Securitize正在构建这个受监管的基础设施层。它们使用SEC注册的过户代理、受监管的ATS、经纪交易商和托管解决方案。

Figure已经在其OPEN平台上推出了自己的代币化股票，并且还有一只股票正在排队上市。Bullish通过收购Equiniti，获得了运营传统证券和代币化证券统一账本的能力，同时提供代币设计和流动性相关服务。Securitize则与纽交所合作开发后者的代币化证券平台，与Computershare合作推动美国发行人的代币化股票，与Jump Trading合作推出去中心化交易所式的代币化股票链上交易。

这些动作的共同特征是：它们不是在传统金融体系之外另建一套系统，而是在传统金融体系的合规框架内，用区块链重新定义清算和结算的效率边界。

Coinbase是这份报告中值得单独讨论的变量。它正在尝试一种与上述基础设施玩家不同的路径：多资产单一交易所。

过去几周，Coinbase为非美国投资者推出了代币化股票、股票永续合约、Pre-IPO永续合约，以及面向美国投资者的受监管加密衍生品市场。其代币化股票由底层股票1:1支持，具备自动股息分配和链上经济的可编程效用。

更关键的是，Coinbase的全球永续合约产品现在覆盖股票、ETF和加密资产，强化了其“一切交易所”的定位。2026年3月，Coinbase推出了面向美国以外零售和机构投资者的24x7 USDC结算股票永续市场，随后又推出了Pre-IPO永续合约。

这里有一个值得深挖的洞察：Coinbase已

[... middle omitted ...]

经纪商的竞争压力，以及美国监管框架的不确定性。

这份报告最具价值的部分，在于它清晰地划分了两种商业模式：中立基础设施定位和交易模型。但它没有完全回答一个根本性问题——在发行方赞助模型中，代币持有者的法律地位究竟如何界定？

目前，Figure和Securitize的模型虽然声称提供完整的股东权利，但这一权利在破产清算、公司治理争议等极端情况下的实际可执行性，尚未经过充分的法律检验。SEC虽然通过无异议函和交易所规则批准释放了积极信号，但距离一个完整的、具有先例效力的监管框架还有距离。

另一个未被充分展开的问题是：代币化是否会改变上市公司的股东沟通和治理方式。如果代币化股票可以编程，那么股息分配、投票权委托、甚至公司行动的通知和执行，都可能实现自动化。但这同时也意味着，传统的中介角色——过户代理、投票顾问、托管银行——可能面临结构性替代。

Bernstein报告提到了“创新豁免”的可能性，但没有深入讨论这一豁免的具体形态和时间表。这恰恰是投资决策中最不确定的变量。

\subsection{[R008] Bernstein：AI正被内存成本压到喘不过气，一场供应链“再校准”不可避免}
{\small\color{refgray}归类：AI/算力：云巨头资本效率竞赛与内存成本冲击}

Bernstein：AI正被内存成本压到喘不过气，一场供应链“再校准”不可避免

当所有人都在关注GPU算力瓶颈时，内存正在悄然成为AI基础设施中最被低估的成本变量。这份Bernstein研报揭示了一个反直觉的判断：市场真正低估的不是需求，而是供给侧定价权转移带来的成本结构重塑。HBM价格即将迎来2-2.5倍的跳升，而经过NVIDIA等GPU供应商的“加价”传导后，超大规模云厂商的AI资本开支可能被迫上调30\%。这不是一个边际变化，而是一个需要重新审视整个AI投资回报模型的系统性冲击。

报告的核心信号是：内存价格压力正在从消费电子蔓延至AI领域。过去三年，HBM作为AI算力的“伴生品”，其定价一直被长期合同锁定，与暴涨的常规DRAM价格形成了巨大背离。这种背离正在扭曲内存供应商的产能分配决策，也将倒逼整个供应链进行一次痛苦的“再校准”。

Bernstein的分析显示，从2025年第三季度到2026年第二季度，常规DRAM价格已经上涨了约4.5倍。而HBM由于采用年度合同定价，价格几乎没有变动。这种定价机制带来的结果是：同样一片晶圆产能，投入常规DRAM产生的收入是HBM的2倍以上，毛利润更是接近3倍。三星和美光在2026年第一季度财报电话会议上均已明确表示，非HBM的DRAM利润率已经领先于HBM，且随着常规DRAM价格继续攀升，这一差距还在扩大。

这里的关键洞察不在于价格差异本身，而在于这个差异对供应行为的扭曲。当内存供应商发现生产常规DRAM比生产HBM更赚钱时，他们会有强烈的动机将产能从HBM转向常规DRAM。而HBM恰恰是当前AI算力最紧缺的组件。这种供给侧的“理性选择”正在制造一个结构性矛盾：AI需求最旺盛的环节，反而面临最严重的供给抑制。

Bernstein估算，要使HBM每片晶圆收入赶上常规DRAM，HBM价格需要上涨约3倍。但报告也指出，内存供应商可能不会采取如此激进的定价策略，因为他们意识到过高的HBM成本会损害整个AI生态的健康发展。SK海力士在财报中也强调，将努力实现“HBM与通用DRAM之间的最优分配”，而非追求收入最大化。因此，Bernstein假设HBM价格在2027年上涨2-2.5倍，即便如此，HBM的盈利能力仍将低于常规DRAM，只是差距会大幅收窄。

这是整份报告中最容易被忽视、但影响最为深远的洞察。HBM并非由超大规模云厂商直接采购，而是被封装在NVIDIA等GPU供应商的加速器内部，成为其销售成本的一部分。当HBM价格上涨时，GPU供应商面临一个两难选择：要么让利压缩自身毛利率，要么通过加价将成本转嫁给下游。

以NVIDIA的Vera Rubin（VR200）机架为例，Bernstein假设NVIDIA的毛利率为75\%。如果HBM成本上升，NVIDIA若要保持毛利率不变，其对机架的提价幅度将是HBM成本涨幅的4倍。具体而言，假设HBM原本占VR200售价的5\%，HBM价格上涨2-2.5倍后，仅HBM成本本身就会推动VR200售价上涨6\%。但经过NVIDIA的4倍加价后，机架价格将上涨24\%。

这个“加价”机制的意义远超数字本身。它意味着HBM的价格波动被GPU供应商的财务结构放大了。对于超大规模云厂商而言，他们面对的不仅是内存供应商的涨价，还有一个被杠杆放大的系统级涨价。Bernstein测算，如果仅将HBM成本直接传导，数据中心总资本开支增加4\%；但如果包含NVIDIA的加价，增幅将扩大到15\%。再加上常规DRAM和NAND价格的大幅上涨，总资本开支增幅将达到约30\%。

30\%的资本开支增幅意味着什么？对于正在建设AI数据中心的超大规模云厂商而言，这意味着他们基于原有预算做的投资回报分析需要全部重算。投资项（I）比预期高出30\%，在收入预期不变的情况下，RO

[... middle omitted ...]

置，甚至可能成为受益者。

4. 三星HBM4技术领先，但“更多HBM”可能意味着更少利润

报告中的一个重要技术判断是：三星在HBM4技术上已取得领先。Bernstein的韩国内存出口追踪数据显示，5月份三星的多芯片内存出口单位价值出现显著提升，这很可能与HBM4开始出货有关。基于此，Bernstein预计三星的HBM市场份额将在今明两年增长。

然而，这里存在一个反直觉的结论：更高的HBM市场份额可能意味着更低的利润率。如前所述，HBM的盈利能力目前低于常规DRAM，且即使价格上涨后，差距仍将存在。三星在财报中表现出对利润率的强烈关注，因此Bernstein认为，即便三星拥有HBM4的技术优势，它也可能选择将更多产能分配给盈利能力更强的常规DRAM。

这意味着HBM市场的份额竞争可能不会像市场预期的那样激烈。技术领先者未必会全力抢占份额，因为这样做会牺牲短期利润。这个判断对于理解HBM供给格局至关重要：供给不会因为技术突破就自动增加，利润考量才是真正的决定因素。

\subsection{[R009] Bernstein：印度真正需要的不是算力租用，而是自己的“DeepSeek 时刻”}
{\small\color{refgray}归类：区域市场：A股硬科技领跑，印度AI自主性困境}

Bernstein：印度真正需要的不是算力租用，而是自己的“DeepSeek 时刻”

这份来自Bernstein的研报，提出了一个在当下争论中容易被忽视的核心判断：印度无法在借来的模型上构建自己的AI未来。当市场将目光集中于算力基础设施的建设和应用层的快速迭代时，报告指出，一个更深层的结构性风险正在浮现——基础模型正在从开放的商品演变为受地缘政治管控的战略资产。对于印度而言，问题不在于是否要参与AI生态，而在于其在AI价值链上的定位是否可持续。

报告选择在此时重提这个话题，本身就意味着一种紧迫感。过去，关于“大语言模型不重要”的论调在印度技术圈并不少见。但近期美国对非公民访问最新AI模型的限制，将这一风险从理论推向了现实。这份研报的价值在于，它没有停留在对依赖性的泛泛批评，而是通过梳理印度信息技术产业的深层结构，解释了为什么印度没有出现自己的“DeepSeek 时刻”，并为决策者提供了一套权衡“开放”与“自主”的政策框架。

1. 基础模型不再是SaaS产品，而是正在成为受管控的“战斗机”

Bernstein报告的核心洞见在于对AI资产属性的重新定义。它指出，技术访问权在过去一直是被配给和管控的。从关键矿产和半导体设备，到GPU，再到如今对前沿模型本身的限制，这一趋势正在打破AI“全球化”的幻觉。报告以Anthropic最新模型对非美国公民的限制为例，认为这不再是孤立事件，而是新常态的开始。

这意味着，基础模型将从一种可以随意订阅的软件服务，转变为类似于“战斗机”的国家战略资源。其背后的逻辑是，谁控制了最前沿的模型，谁就掌握了定义下一个时代应用能力的钥匙。对于正在说服自己“基于外国LLM构建应用、靠数据中心收租”是明智AI战略的印度而言，这一判断具有深刻的警示意义。Bernstein并未全盘否定这种参与方式，但它明确警告：印度至少需要意识到，将核心AI模型外包出去所蕴含的风险。

2. 印度没有出现“DeepSeek 时刻”是结构性的，而非战略选择

一个看似矛盾的现象是：印度是全球数据的重要来源，却未能利用这一优势打造具有全球竞争力的生成式AI系统。报告认为，这种缺失是结构性的。印度科技生态长期以来以服务为导向，缺乏像搜索、社交媒体或即时通讯这样能产生丰富、结构化数据集的消费级平台。这种数据生态的匮乏，从未催生出构建基础模型所需的人才储备和学术深度。

更深层的原因在于，印度IT服务模式的“路径依赖”。在这种模式下，低成本劳动力本质上是为全球巨头开发的软件进行微调。这进一步固化了印度企业停留在应用层的倾向。报告指出，许多印度顶尖机构的负责人曾主张“印度不需要自己的大语言模型，只需聚焦AI应用”。Bernstein认为，这些观点并非深思熟虑的战略选择，而更多是印度走到今天这条路径的自然反应。这揭示了一个核心问题：印度的技术生态缺乏从底层创造“原生AI”的土壤。

3. 系统性依赖的真相：从互联网到AI，印度几乎在所有关键层都处于下游

报告通过详尽的表格（Exhibit 1和Exhibit 2）展示了印度在技术栈上的系统性依赖。从CPU、GPU等核心计算芯片，到操作系统、企业软件，再到云服务、社交媒体，印度几乎在所有关键层都依赖美国公司。这种依赖不是孤立的，而是贯穿整个技术栈的。过去，印度在互联网时代接受了这种模式——保持开放，让全球平台成为其数字经济运行的轨道。

但AI不同。报告认为，AI不仅仅是又一个应用层，它正在成为所有企业流程和消费者交互的底层“轨道”。更重要的是，现在还为时过早。与以往印度总是晚入场的技术周期不同，AI的轮廓仍在塑造之中。这正是当前关于“主权AI栈”争论的根源。印度并非突然发现了技术自主的重要性，而是因为错过这一波的代价将远高于以往。如果印度像过去一样等待20-30年才介入，

[... middle omitted ...]

的“口袋阵地”。这些能力可以逐步演变为可扩展的解决方案，并最终作为全球产品出口。报告甚至指出了AI与物理世界结合的可能性，例如在训练工业机器人或下一代人形机器人时，特定工业场景的数据集可能比使用最前沿的通用AI模型更为关键。

这一判断意味着，印度AI领域的价值创造可能不会来自与OpenAI或Google的正面竞争，而是来自对特定行业知识的数字化和智能化。这并非追求万亿市值的科技巨头，而是通过将AI与真实世界的资产和工作流结合，在价值链中寻找属于自己的位置。

5. 政策框架的两难：在“访问”与“自主”之间，没有完美答案

报告在政策层面提出了一个坦率的框架，承认了决策的复杂性。它否定了“为国内企业圈定AI领域”这种缺乏现实可行性的想法。Bernstein提出了两个战略杠杆，但明确指出两者都不完美。

第一种选择是限制或放缓对全球AI模型的访问，同时将大量资本和人才引导至建设印度本土的LLM能力。这相当于“闭关修炼”，但可能在全球一体化的技术生态中导致自身孤立和落后。

\subsection{[R010] JPM：市场低估了“供给瓶颈”在日本MLCC股中的定价权，而非AI需求本身}
{\small\color{refgray}归类：区域市场：A股硬科技领跑，印度AI自主性困境 / 风险偏好与资金流：供给瓶颈定价权重估}

JPM：市场低估了“供给瓶颈”在日本MLCC股中的定价权，而非AI需求本身

*主判断：** 当前日本MLCC相关股票的上涨，并非单纯由AI需求驱动，更核心的是其作为“供给瓶颈”环节所享有的超额利润定价权。这一轮上涨的结构与2000年代中期的EM商品热潮高度相似，真正的风险不在于需求见顶，而在于“瓶颈效应”向下游转移。

这份来自JPM量化与衍生品策略团队的研报，提供了一个与市场主流叙事截然不同的分析框架。大多数投资者将村田制作所与太阳诱电的强势表现归结为AI带来的基本面增长，但该报告从技术面和资金结构切入，认为这轮上涨的内在逻辑更接近“供给瓶颈溢价”的重现。报告的核心洞察在于：超额利润的可持续性，更多取决于企业在供应链网络中的“瓶颈位置”，而非其自身的竞争力。

为什么现在这个判断重要？因为市场正在将MLCC股票视为AI周期的纯受益者，却忽略了其背后一个更脆弱的结构性特征——高Beta、低Quality。这意味着一旦“瓶颈”转移，这些股票的调整速度可能远超基本面投资者的预期。报告提供了一个历史参照系，以及一套用于观察“瓶颈转移”的信号框架。

1. 超额利润的来源并非技术壁垒，而是供应链中的“瓶颈位置”

报告最引人注目的论点，是将当前MLCC股票的上涨与2000年代EM热潮中的特定日本股票进行对比。在EM热潮初期（ 年），住友金属矿山（镍冶炼）、三井金属（超薄铜箔）、信越化学（硅片）以及新报国材料（低热膨胀合金）等公司，都曾经历过一段超额利润的爆发期。这些公司的共同点是什么？它们都处于当时全球供应链中难以快速扩产的“瓶颈环节”。

报告通过动态时间规整（DTW）方法，将当前MLCC股票的股价走势与上述EM热潮中的瓶颈股进行相似度测量，发现两者呈现出高度一致的模式。这意味着，当前市场对MLCC股票的定价，并非完全基于其技术领先或产品独特性，而是基于一个更结构性的前提：在AI需求爆发的背景下，MLCC的供给结构（寡头垄断、产能扩张缓慢）使其成为供应链中新的“瓶颈”。

*这意味着什么？** 对于投资者而言，评估这些公司价值的关键变量，不应仅仅是AI需求的增速，更应关注其“瓶颈地位”是否被其他环节侵蚀。一旦上游半导体、设备或材料环节出现新的瓶颈，资本可能会从MLCC板块快速撤离，转向新的“瓶颈资产”。这种转移并非因为MLCC公司本身变差了，而是超额利润的来源发生了位移。

报告对MLCC相关股票及所属电气设备行业的因子结构进行了估算，发现一个值得警惕的特征：动量（Momentum）和贝塔（Beta）因子显著偏高，而质量（Quality）因子则处于低位。这意味着，当前股价的支撑主要来自价格趋势和市场相关性，而非盈利质量的持续改善。

从短期交易的角度看，这种因子组合反而意味着较高的预期回报。报告明确指出，核心买家是趋势跟踪型短期动量交易者，只要盈利边际改善的势头得以维持，这些交易者仍有建仓空间。然而，这种结构的脆弱性同样明显——一旦动量逆转，高Beta特性会放大下行速度，而低Quality意味着缺乏基本面支撑来缓冲下跌。

*这引申出一个关键问题：** 对于长线投资者而言，当前价格是否已经充分反映了“瓶颈溢价”？报告给出的暗示是，全球主动型投资信托对日本科技硬件与设备板块的配置仍处于中性至小幅低配状态，这意味着长期资本尚未大规模入场。但这恰恰可能成为未来价格修正时的买入力量，而非当前追高的理由。

报告的另一重要贡献，是将MLCC股票的走势置于全球资产配置的宏观叙事中。数据显示，全球主动型基金正在持续减少对中国（包括A股和港股）的区域配置，并将资金转向日本、韩国和台湾。这种“减中国、增日本”的再平衡，并非短期战术调整，而是基于地缘政治风险、供应链重构以及增长预期差异的战略性转移。

对于日本电气设备板块而言，这一资

[... middle omitted ...]

的预期。但与此同时，投资者需要警惕的是，区域配置的转向同样具有周期性，一旦中国资产重新获得吸引力，或者日本市场出现系统性风险，这股“顺风”也可能迅速转为逆风。

4. 报告尚未完全回答的关键问题：瓶颈转移的触发条件与时间窗口

报告最精彩的部分，是对EM热潮末期“瓶颈转移”的复盘。2007年下半年，EM相关股票集体崩盘，表面诱因是次贷危机，但JPM认为根本原因在于供给瓶颈从电子元器件转向了大宗商品和海运价格。当投机资本发现新的瓶颈环节出现时，它们迅速从原有的股票组合中撤出，转向新的目标。这种“瓶颈级联效应”导致了旧有板块的同步崩溃。

当前，报告认为MLCC板块尚未出现这种转移的明确信号。上游半导体、制造设备、HBM和CoWoS等环节仍处于扩张趋势，瓶颈位置大概率仍将停留在当前股票组合中。但这里存在一个尚未完全回答的问题：“瓶颈转移”的触发条件是什么？ 是某个关键上游环节的产能突破？还是AI需求从训练侧转向推理侧带来的需求结构变化？抑或是地缘政治事件引发的供应链重塑？

\subsection{[R011] GS：家电内需压力尚未出清，但出口的结构性机会比想象中更值得关注}
{\small\color{refgray}归类：地产与消费：中国内需磨底，日本零售整合拐点}

GS：家电内需压力尚未出清，但出口的结构性机会比想象中更值得关注

五月的数据再次确认了一个判断：中国家电行业正处于“内需承压、外需修复”的错位周期中。GS最新发布的五月家电追踪报告显示，国内零售和出厂数据延续了4月的疲软态势，而出口则出现了超预期的环比改善。这份报告的核心信号并非简单的“内冷外热”，而在于两个更值得推敲的结构性变化：一是国内需求的下行斜率正在放缓，基数效应将在6月下旬开始提供缓冲；二是出口的韧性并非仅靠低基数，中国企业在海外市场的份额持续扩张正在成为更持久的驱动力。

对于投资者而言，当前阶段的关键问题不再是“行业是否见底”，而是“底部的形态和分化路径是什么”。GS在报告中明确指出了对美的集团、九号公司和海信家电的偏好，其逻辑并非押注行业整体反转，而是捕捉那些在周期底部仍能通过海外扩张、产品组合和股东回报构建安全边际的公司。

1. 五月内需数据低于预期，但真正的拐点信号在于基数的“前高后低”

五月国内家电零售额同比下降16\%，与4月-15\%的降幅基本持平。空调国内出厂量更是同比下降15\%，显著低于此前行业生产计划所隐含的-9\%的预期。GS认为，这主要是受到去年同期需求前置以及4月以来部分品牌提价的叠加影响。

但比单月数据更重要的是趋势判断。报告指出，由于去年6月“以旧换新”政策的效果开始逐步消退，从今年6月下旬开始，国内需求的同比基数压力将明显缓解。这意味着，即使终端需求没有显著回暖，7月之后的同比读数也可能出现自然改善。换句话说，当前市场对“内需持续恶化”的线性外推，很可能低估了基数效应带来的读数修复。

这里需要谨慎的是：读数的改善并不等于需求的复苏。真正的拐点需要看到终端动销的实质回升，而非仅仅是统计意义上的低基数。GS也据此将二季度主要白电国内出货增长预期从此前的低个位数下调至低十位数负增长。

2. 出口超预期改善的背后，是中国企业份额扩张与海外需求的同步修复

与内需形成鲜明对比的是，五月家电出口出现了明显的环比改善。空调出口降幅从4月的-12\%收窄至-4\%，显著好于生产计划预期的-6\%的降幅。洗衣机出口增速从+7\%提升至+14\%，冰箱出口更是由负转正至+13\%。

GS将这一改善归因于两个因素的叠加：一是去年同期的低基数，二是中国企业在海外市场的持续份额扩张。报告特别提到，尽管中东局势仍存在不确定性，但中国品牌的出口仍展现出韧性。与此同时，美国和欧洲的消费者信心指数以及美国成屋销售领先指标——待售房屋销售数据，均在5月出现了环比改善。

这意味着，出口的改善并非单纯的中国供给端“以价换量”，而是海外需求端本身也在边际修复。对于投资者而言，这比单纯的低基数更具持续性价值。GS认为，随着美伊谅解备忘录签署后更多细节的披露，市场需求的可视度有望进一步提升，这将为中国头部企业进一步扩大海外份额提供空间。

GS的数据清晰地展示了品牌层面的分化。五月空调国内出货量，海尔同比增长5\%，海信下降5\%，而美的和格力分别下降35\%和6\%。海尔在国内市场的相对稳健，与其在高端品牌卡萨帝和渠道效率上的持续投入有关。而美的国内出货的大幅下滑，部分原因可能是其主动调整渠道库存或产品结构。

但在出口端，美的空调出口同比增长4\%，显著好于格力的-33\%和海信的-17\%。这一反差揭示了不同公司的战略侧重：美的通过海外工厂的本地化布局，在应对贸易摩擦和汇率波动时拥有更大的灵活性。报告也明确指出，中国企业的整体出口量应大于IOL统计的出厂数据，因为头部企业如美的也在通过海外工厂直接出口。

在冰箱领域，美的和海尔的分化同样显著。四月冰箱总出货量，美的同比增长25\%，海尔持平，海信增长15\%。美的在国内和出口市场均实现了两位数增长，显示出其在产品组合和渠道覆盖上的优势。

这些数据合在一起意味着：当前周期中，

[... middle omitted ...]

值得深思。在原材料成本——尤其是铜、钢等关键大宗商品价格——过去一个月保持基本稳定的背景下，部分品牌仍选择降价促销，说明终端需求压力已经传导至定价策略层面。这并非健康的竞争信号。如果成本端未来出现上行，而品牌因需求疲软无法转嫁成本，利润率将面临双重挤压。

GS并未在报告中给出具体的利润率预测，但其隐含的判断是：能够通过产品结构升级和海外市场对冲国内价格压力的公司，将拥有更强的盈利韧性。

报告还提及了一个容易被忽视的变量：航运成本。过去一个月，中国出口集装箱运价指数(CCFI)持续上升，同比增速也在6月中旬进一步扩大。对于家电出口企业而言，航运成本上升会直接侵蚀FOB模式下的利润空间，或迫使企业调整海外定价策略。

GS并未对此展开深入分析，但这是一个值得投资者持续跟踪的变量。如果航运成本持续高位运行，那些在海外拥有本地化产能的企业——如美的——将获得相对优势，因为它们可以通过本地生产、本地销售来规避部分运费压力。而对于依赖国内生产再出口的企业，成本压力将更为直接。

\subsection{[R012] GS：A股硬科技正在主导中国盈利增长的结构性分化}
{\small\color{refgray}归类：区域市场：A股硬科技领跑，印度AI自主性困境 / 风险偏好与资金流：供给瓶颈定价权重估}

2026年第一季度财报季刚刚落幕。GS覆盖了约7000家中国上市公司，并分析了超过4500份业绩会纪要。这份报告传递的核心信号是：中国上市公司的盈利增长正在经历一场深刻的结构性分化，而A股硬科技板块正站在这一分化的主导一侧。

这不是一个简单的“A股优于H股”的判断。真正值得关注的，是支撑这一分化的三个底层力量：AI利润池正在从软件向硬件回摆，企业全球化布局在关税压力下仍在深化，以及上市公司资本配置逻辑正在发生根本性转变。这些力量叠加在一起，正在重塑中国权益资产的定价框架。

GS在报告中明确提出，维持对A股相对于H股的偏好。其预测显示，CSI300在2026年全年盈利增速有望达到20\%，而MSCI中国指数仅为8\%。在即将到来的二季度业绩期，这一差距预计仍将显著存在——CSI300的盈利增速预期为11\%，MSCI中国则接近零。

这一判断的底层逻辑，值得每一位置身中国市场的投资者仔细拆解。

1. A股与H股的盈利鸿沟，根源在于“硬科技”与“软科技”的资产结构差异

2026年一季度，全部中国上市公司的净利润同比增长6\%，延续了2025年7\%的增长节奏。但这一整体数字掩盖了指数层面的巨大分化。CSI300一季度盈利同比增长9\%，而MSCI中国指数则同比下滑8\%。更极端的对比出现在创业板和科创板——两者分别录得22\%和198\%的同比盈利增长。

这一差距并非偶然。A股指数中，科技硬件、半导体、电力设备等“硬科技”板块的权重远高于离岸市场。而MSCI中国指数中，互联网、电商、本地生活等“软科技”和消费服务类资产的占比更高。后者的盈利正受到快速商超补贴竞争和AI资本开支扩张的双重挤压。

GS的分析师团队在业绩会纪要的文本分析中捕捉到了一个关键信号：市场参与者的AI讨论焦点，已经从下游软件应用向上游硬件基础设施转移。这一转变并非情绪驱动，而是有企业资本开支作为实物支撑。无论是美国的超大规模云服务商还是中国的云厂商，都在持续上调AI相关资本开支计划。这意味着，硬科技板块的盈利增长并非昙花一现，而是有可见的需求支撑。

对于投资者而言，这一结构差异意味着，选择A股还是H股，本质上是在选择两种不同的盈利增长驱动力。前者更直接受益于AI基础设施投资周期的上行，后者则需要等待互联网平台的盈利拐点。

2. AI利润池正在经历一次向硬件的显著回摆，中国在这一过程中面临份额压力

GS构建了一个覆盖全球3700多家公司、总市值42万亿美元的AI权益宇宙。数据显示，全球AI利润池的分配正在发生重大变化。2024年，中国在全球AI利润中的占比约为10\%，到2025年降至9\%，2026年一季度进一步压缩至8\%。与此同时，以东北亚其他地区（主要是韩国和台湾）为代表的区域，凭借在先进代工和存储芯片领域的优势，全球份额从2024年的24\%上升至2026年一季度的27\%。

在中国内部，AI供应链的利润分配同样出现了明显的硬件回摆。2024年，硬件在中国AI利润池中的占比为41\%，到2026年一季度已反弹至54\%。半导体板块是这一回摆的核心受益者，其中存储芯片、IC设计、材料环节的利润份额增长最为显著。相比之下，AI下游应用的利润份额从2024年的59\%压缩至一季度的46\%。

这一回摆的背后，是中美两国云厂商持续扩张的AI资本开支。GS的互联网分析师团队估算，无论是美国的四大云巨头还是中国的头部云厂商，2026年的AI资本开支计划仍在加速。这意味着，硬件端的订单可见性在短期内不会减弱。

但这里有一个值得关注的问题：中国AI企业的利润率仍然显著低于全球同行。GS的图表显示，中国AI公司的平均净利润率远低于美国AI公司，也低于其他地区的AI公司。这意味着，尽管中国硬件企业受益于资本开支周期，但它们在价值链中的议价能力可能仍受制于技术差距和竞争格

[... middle omitted ...]


这一分化揭示了一个关键事实：利润率的改善并非来自竞争缓和，而是来自技术壁垒和成本结构差异。半导体行业的利润率提升，源于AI需求拉动的产品结构升级。工业金属的利润率改善，则更多与供给侧约束有关。而消费电子、新能源汽车等竞争激烈的行业，利润率仍在被价格战侵蚀。

这意味着，在当前的宏观环境下，选择行业比选择时机更重要。那些具备技术壁垒、能够将规模转化为议价权的企业，正在享受利润率改善的红利。而那些依赖市场份额扩张但缺乏定价权的企业，营收增长可能只是“虚假繁荣”。

关税不确定性并未阻止中国企业的全球化步伐。GS的数据显示，中国上市公司的海外收入占比从2024年的16\%进一步上升至2025年的18\%。汽车板块是这一趋势的领跑者，海外收入占比从2024年的21\%跃升至29\%，主要得益于电动汽车的海外扩张。

这一全球化趋势正在变得更加广泛。半导体、交通运输、资本品、耐用品和医药生物等板块的海外收入占比均在提升。只有科技硬件（受智能手机拖累）和能源板块的海外收入占比出现下降。

\subsection{[R013] Bernstein：数据中心供电的真正瓶颈不在技术，在电网的“连接组织”}
{\small\color{refgray}归类：储能与新能源：供应链本地化重构拐点}

Bernstein：数据中心供电的真正瓶颈不在技术，在电网的“连接组织”

过去一年，市场对AI和数据中心电力需求的讨论，几乎都围绕一个核心问题：电力从哪里来？天然气、核能、可再生能源、燃料电池，每一种技术路线都有自己的拥护者。但Bernstein刚刚发布的一份基于50多家数据中心运营商调研的报告，给出了一个更底层、也更令人不安的回答——问题从来不是“发多少电”，而是“电能不能送到”。

这份报告最值得关注的判断是：数据中心运营商最想要的，是电网连接。不是离网自建，不是孤岛运行，而是接入大电网。但他们面临的现实是，电网的“连接组织”——从并网审批到输电容量到设备交付——正在成为整个AI基础设施扩张的最大瓶颈。这个瓶颈的持续时间，不是一两年，而是以十年为单位。

1. 电网连接是首选，但并网等待时间已经突破运营商的容忍底线

调研中，98\%的受访数据中心运营商表示，电网连接是未来5-10年新项目的首选供电方式。对于100MW以上的大型数据中心，86\%的运营商明确偏好“电网连接为主、现场发电为辅”的模式。现场发电（behind-the-meter）目前仅占数据中心电力市场的约5\%，Bernstein预计到2030年代，这个比例也只会上升到5-10\%。

这个偏好本身并不令人意外。电网供电在成本、可靠性和可扩展性上，通常优于任何形式的现场自发电。真正值得关注的是，运营商对并网等待时间的容忍度极低。调研显示，绝大多数受访者认为“等待超过三年”是不可接受的。但现实是，目前美国主要区域的并网排队时间普遍在4年以上，CAISO区域甚至更久。

这意味着什么？意味着在电网基础设施没有显著改善之前，大量数据中心项目的实际落地时间表，可能要比市场预期的更长。这不是需求不足的问题，而是供给侧的交付时间被电网瓶颈锁死了。

当被问及选择供电方式的最重要因素时，80\%的受访者将“可靠性”排在第一位，20\%将“速度”排在第一位。成本、碳排放、监管等因素，在优先级上都明显靠后。甚至有10\%的受访者明确表示，成本不是驱动因素，他们愿意为现场发电支付更高费用。

这个排序对投资逻辑的影响是直接的。拥有已并网资产的独立发电商（IPP），如Constellation Energy和Vistra，在Bernstein的框架中获得了跑赢评级（Outperform），因为它们不需要等待新的并网审批。而能够快速交付现场发电方案的供应商，如Bloom Energy（燃料电池，55-90天交付），也在速度维度上获得了市场表现评级（Market-Perform）。

但这里有一个容易被忽略的深层含义：如果速度和可靠性压倒成本，那么传统以“最低成本”为核心的电厂估值模型，可能需要调整。在数据中心供电的语境下，能够提供确定性交付时间的资产，应该享有估值溢价。这个溢价目前是否已经被市场完全定价，报告没有展开，但这是一个值得持续追踪的问题。

调研显示，超过95\%的受访者将现场备份电源列为关键设计要素。更值得关注的是，90\%以上的受访者认为，现场备份不是过渡安排，而是长期存在的架构特征。未来大型数据中心将普遍采用“电网+现场发电”的混合供电架构。

在具体技术路线上，45\%的受访者偏好化石燃料（天然气涡轮机或往复式内燃机），其次是太阳能加储能（Solar+BESS）和燃料电池。这个分布说明，在“可靠性第一”的约束下，运营商对技术路线的选择是务实的：天然气是成熟的、可调度的；太阳能加储能在白天有优势；燃料电池在安静和低排放场景中有价值。

这个趋势对电网设备供应商GE Vernova是明确的利好。Bernstein在报告中特别指出，市场可能低估了与数据中心相关的“连接组织”需求——即变压器、开关设备、输电线路等电网基础设施的升级和扩容。这不是一个短暂的周期，而是长达数十年

[... middle omitted ...]

热公司Fervo（跑赢评级）也已经与超大规模客户签署了超过650MW的长协合同。但可持续性本身，目前还不足以成为驱动供电决策的独立变量。

5. 报告尚未完全回答的关键问题：电网升级的“谁买单”和“多快能”

Bernstein的调研清晰地指出了瓶颈所在——并网等待时间、审批不确定性、输电容量限制。但报告没有完全回答的问题是：谁来为电网升级买单？以及，升级的速度能有多快？

从历史经验看，美国电网基础设施的投资决策和成本回收，涉及联邦、州、公用事业公司、监管机构等多个利益相关方。数据中心运营商愿意为速度和可靠性支付溢价，但这个溢价是否足以覆盖电网升级的全部成本，并转化为公用事业公司和设备供应商的可持续盈利，仍需要更清晰的监管和政策信号。

此外，报告提到“互联队列”（interconnection queues）是4年以上的等待时间，但并没有深入分析：这些队列中的项目有多少是真正会建设的？有多少是投机性的？如果大量项目最终无法落地，电网升级的投资回报假设是否需要调整？

\subsection{[R014] GS：欧盟暂缓对中国割草机器人征收临时反倾销税，真正的变量不是关税本身，而是企业能否利用这个窗口期重塑竞争壁垒}
{\small\color{refgray}归类：地产与消费：中国内需磨底，日本零售整合拐点}

GS：欧盟暂缓对中国割草机器人征收临时反倾销税，真正的变量不是关税本身，而是企业能否利用这个窗口期重塑竞争壁垒

2026年6月19日，欧盟委员会宣布，将继续对中国产割草机器人进行反倾销调查，但现阶段不会征收临时反倾销税。理由是案件技术复杂性。

这份GS研报的核心判断是：投资者应当将此视为积极信号。但真正值得深入推敲的，不是“短期内关税没加”这个事实本身，而是这个时间窗口对行业竞争格局意味着什么。

GS覆盖的三家中国公司——九号公司、石头科技和科沃斯——均可能因此受益，其中九号公司受益最为显著。但这份研报的深层逻辑在于：欧盟的决定为中国企业争取了至少6到12个月的缓冲期。这段时间足够做什么？足够一家企业把产能从中国转移到海外，但不足以让转移后的成本劣势完全消失。所以，真正的分化将从这里开始。

1. 2026年关税风险基本排除，但真正考验的是企业能否把窗口期转化为结构性优势

GS在研报中明确指出，多数中国割草机器人企业原本的计划是在2026年上半年完成对欧洲的大部分出货，并预期从6月起可能面临临时反倾销税的影响。现在欧盟决定暂不征收临时关税，意味着2026年全年适用的关税税率大概率保持不变。

这对企业最直接的影响是利润表。GS判断，中国企业可能不需要通过涨价来对冲关税成本，因此短期内的利润率和市场份额竞争力将得以维持。但这里有一个被市场低估的变量：定价纪律。

GS特别提到，至少到调查结果更加明朗之前，中国割草机器人企业仍将保持定价纪律。这句话的潜台词是，一旦企业认为关税风险解除，可能会重新回到价格战的老路上。而欧盟暂缓加税，反而给了头部企业一个机会——在不牺牲利润率的前提下，继续用产品力而非价格来竞争。

所以，2026年关税风险排除，只是短期利好。真正的变量是中国企业会如何利用这个时间窗口——是继续卷价格，还是加速构建品牌和渠道壁垒。

2. 九号公司是最大受益者，但受益逻辑不只是“关税没加”，更是估值压制因素的解除

GS将九号公司列为此次事件的最大受益者，理由是割草机器人业务占其2025年总收入的9\%，利润贡献更高。但更值得关注的是GS对九号公司的估值判断。

GS指出，欧盟反倾销调查一直是九号公司估值的一个压制因素。这个说法值得展开。市场对九号公司的担忧，不只是关税本身，而是关税不确定性带来的“不可预测性”——你不知道关税什么时候加、加多少、是否追溯。这种不确定性使得投资者难以对九号公司的海外业务进行合理定价。

现在，虽然调查仍在继续，但至少短期内的不确定性大幅降低。GS认为，这至少能在短期内缓解投资者情绪。而情绪的修复，往往先于基本面的改善。

但这里有一个需要验证的假设：九号公司是否真的有能力在海外调整产能。GS在研报中表示，领先品牌（包括九号）由于产品定价更高，且具备全球生产调整能力，比小企业更能应对潜在的关税影响。九号公司相对更高的产品定价，意味着它在面对关税时有更大的缓冲空间——它可以适度吸收关税成本，而不必像低端品牌那样被迫大幅涨价或牺牲利润。

所以，九号公司的受益逻辑是双重的：短期是估值压制因素的解除，中期是定价能力和产能布局带来的结构性优势。

3. 产能转移的成本缺口依然存在，但更长的准备时间可以缩小差距

这是研报中最具洞察力的部分之一。GS指出，如果中国企业选择将制造产能转移到海外，初期阶段的制造成本通常会增加10\%-20\%，原因是劳动力效率较低和供应链不成熟。

这个数字本身并不令人意外。但GS接着做了一个关键判断：更长的准备时间可以帮助提高搬迁后的运营效率，尽管无法完全消除成本差距。

这意味着，欧盟暂缓加税，不仅给了中国企业一个“喘口气”的机会，更重要的是给了它们一个“准备得更充分”的机会。如果企业能够利用这段时间优化海外工厂的选址、供应链布局和生产流程，那么即

[... middle omitted ...]

答的关键问题：反倾销税是否会追溯，以及2027年之后的竞争格局

GS在研报中明确提醒，欧盟仍在继续调查，最终仍可能决定征收反倾销税，并且可能具有追溯效力。

这是一个被市场低估的风险。如果最终反倾销税具有追溯效力，那么即使2026年没有加税，企业也可能在未来需要补缴过去一段时间内产生的关税差额。这意味着，企业现在享受的“免税窗口”可能只是暂时的，未来的实际税负可能比预期更高。

GS对此的判断是，当前公告为中国企业留下了更长的准备时间，但并未完全消除风险。这个表述非常克制，但背后的含义是：投资者不应因为短期利好就忽视长期风险。

另一个未被完全回答的问题是：2027年之后的竞争格局会是什么样？如果反倾销税最终落地，中国企业的成本优势将被削弱，欧洲本土企业（如Husqvarna）是否能够借机收复失地？GS对此没有展开，但这是一个值得持续跟踪的问题。

基于这份研报的逻辑，投资者在评估中国割草机器人企业时，应该从以下三个维度建立观察框架，而不是仅仅关注反倾销调查的进展：

\subsection{[R015] GS：中国家庭资产正在经历从房产到金融资产的结构性迁移}
{\small\color{refgray}归类：宏观与利率：中国内需磨底，家庭资产负债表再平衡 / 地产与消费：中国内需磨底，日本零售整合拐点}

五月的经济数据提供了一个令人不安的对照：出口同比增长19.4\%，但零售销售同比下降0.6\%——这是除疫情封控期外，中国首次出现零售负增长。出口强劲与内需疲软之间的裂口，已经不只是周期性的问题，它正在重塑中国经济的底层结构。

GS在最新发布的研报中，用一份长期专题研究给出了一个关键判断：中国家庭资产负债表的再平衡，可能才是理解未来五年资产价格走向的真正主线。这份报告的核心洞察不在于短期数据波动，而在于一个正在发生的结构性迁移——中国家庭资产正在从房产向金融资产转移，而这一过程的深度和速度，可能被市场严重低估。

1. 五月数据暴露的不只是需求疲软，而是家庭部门正在经历主动去杠杆

五月零售销售同比下降0.6\%，固定资产投资同样录得同比收缩。GS指出，这是“除疫情封控期外首次出现负增长”。出口增长19.4\%与内需收缩并存，意味着中国经济的“双轨”格局已经固化：生产端和出口端依然有韧性，但消费端和投资端的动力正在流失。

这里需要追问一个“所以呢”的问题：消费收缩是因为收入下降，还是因为家庭主动减少了支出？

GS的报告虽然没有直接给出答案，但结合其同期发布的家庭资产负债表研究，逻辑链条是清晰的。中国家庭资产中房产占比高达52\%，而房价持续下跌意味着家庭净资产在缩水。当资产端缩水时，家庭部门为了修复资产负债表，会主动减少消费、增加储蓄。这不是需求暂时被压抑，而是家庭部门正在经历一轮自发的去杠杆。

这意味着，短期刺激政策可能难以撬动消费。如果家庭部门的去杠杆行为是结构性的，那么任何依赖消费反弹的经济预测都需要重新审视其假设。

2. 家庭资产结构正在经历从“房产主导”到“金融资产主导”的转折

GS在最新发布的“中国长期系列”研究中，构建了一套中国家庭资产负债表数据集。截至2026年第一季度，中国家庭总资产中，房产占比52\%，现金和存款占比25\%，两者合计接近80\%。而权益和保险类资产的占比仍然较低。

但GS预计，未来几年权益和保险类资产的占比将显著扩大，而房产占比将继续收缩。这个判断的核心驱动因素有两个：一是房价持续下行导致的财富效应逆转，二是政策层面正在推动的资本市场改革和人民币国际化。

这里的关键问题是：这种迁移的速度有多快？GS的报告给出了一个方向性的判断，但没有给出具体的时间表。从历史经验来看，日本家庭资产从房产向金融资产的转移用了将近十年。中国市场可能因为政策推动而加速，但居民对房产的偏好惯性仍然很强。

对投资者而言，这意味着中国金融资产的长期需求结构正在发生变化。如果家庭部门每年将1\%的资产从房产转移至权益市场，对应的资金规模就在万亿级别。这可能会成为A股和港股市场长期资金流入的重要来源。

3. 陆家嘴论坛释放的信号：货币政策框架转型与人民币国际化正在同步推进

GS在另一份报告中详细解读了今年陆家嘴论坛上释放的关键信号。央行行长潘功胜宣布将隔夜回购利率走廊从70个基点收窄至50个基点。这看似是一个技术性调整，但其含义是深远的：中国正在从数量型货币政策框架向价格型框架转型。

利率走廊收窄意味着央行对短期利率的控制更加精准，市场利率的波动性将降低。这对债券市场是利好，但更重要的是，它为未来的利率市场化改革铺平了道路。

与此同时，论坛上宣布的几项措施值得关注：为外国官方机构设立FIMA人民币回购便利、在香港推出五年期国债期货、在上海建立离岸人民币交易中心。这些措施的共同指向是：人民币国际化正在从“贸易结算”阶段进入“金融资产储备”阶段。

GS的报告指出，这些措施“表明人民币国际化正在加速”。但这里有一个尚未被充分讨论的隐含含义：如果人民币国际化成功，中国债券市场将迎来全球央行和主权基金的持续配置需求。这会改变中国利率的定价逻辑，使其从国内因素驱动转向全球因素驱动。

GS的家庭资产负债表研究

[... middle omitted ...]

eleveraging），但没有给出量化终点。这意味着，投资者需要密切关注两个领先指标：一是70城房价数据中一线城市的走势——GS已经注意到五月一线城市房价环比出现上涨；二是居民中长期贷款的变化，如果居民开始重新增加借贷，可能意味着去杠杆接近尾声。

基于GS的这份研报，投资者可以建立一个更清晰的观察框架，而不是被月度经济数据的波动牵着走。

第一个信号是家庭资产配置的边际变化。如果银行存款增速开始下降，而公募基金和保险产品的销售开始回暖，这可能意味着家庭资产从存款向权益的迁移已经开始加速。

第二个信号是人民币资产的全球吸引力。FIMA回购便利和离岸国债期货的推出，为境外机构提供了更便利的人民币资产配置工具。如果境外机构对中国债券的持有量出现趋势性增长，这可能是一个重要的先行指标。

第三个信号是政策框架的转型速度。利率走廊收窄只是第一步。如果未来央行进一步推进利率市场化，甚至考虑引入类似美联储的“前瞻指引”机制，那将意味着中国金融市场的定价逻辑正在发生根本性变化。

\subsection{[R016] GS：欧洲化工的真正危机不是红海冲突，而是冲突结束后暴露的结构性过剩}
{\small\color{refgray}归类：能源与大宗：霍尔木兹冲击重塑油需，AI压缩深水油田周期}

GS：欧洲化工的真正危机不是红海冲突，而是冲突结束后暴露的结构性过剩

当市场还在关注地缘冲突何时结束、物流何时恢复时，一份来自GS的研报给出了一个反直觉的判断：冲突的终结可能才是欧洲化工行业真正考验的开始。

GS分析师团队在参加ICIS与欧洲化工分销商协会（Fecc）联合举办的CEO圆桌会议后，发布了一份基调偏冷的研究笔记。出席会议的包括Equilex董事会主席Patrick Elie、Integra Petrochemicals CEO Gina Fyffe、IMCD CEO Marcus Jordan以及National Chemical Co. CEO Alan Looney。这些行业高管的共识性判断，指向一个令人不安的前景：当前的市场疲软并非单纯的季节性因素，而是价值毁灭的真实信号。

这份研报的核心洞察可以浓缩为一句话：欧洲化工行业正在经历一场“双重挤压”——短期面临需求疲软与物流紊乱带来的交易混乱，中期则面临中国产能外溢、欧洲本地产能永久性关停，以及结构性配方重构的三重冲击。GS认为，真正值得关注的不是冲突何时结束，而是结束后暴露出来的过剩产能将如何重塑全球化工竞争格局。

市场普遍认为红海冲突的结束将带来供应链正常化，但GS从圆桌会议中捕捉到的信号恰恰相反：冲突结束可能让行业变得更糟。

Integra Petrochemicals的CEO明确指出，即使假设冲突已经结束，正常化所需的时间将远超预期。仅清理积压的船舶就需要6至8周。第三季度预计将非常混乱，因为买家害怕做决策，对公平定价存在不确定性。

更深层的担忧在于，一旦冲突结束，生产恢复常态，供应重新变得充足，市场将直接暴露在潜在过剩产能面前。这种过剩将被释放的中东原料进一步放大。Equilex主席警告，当前市场上已经出现了明确的价值毁灭信号，而不仅仅是季节性需求疲软。

这意味着，对于欧洲化工企业而言，过去两年在供应链紧张中获得的定价权和议价能力，可能在冲突结束后迅速瓦解。那些依赖短缺逻辑维持利润的公司，将面临最直接的冲击。

GS研报中一个被反复提及的关键变量，是中国在这场危机中的战略定位。Equilex主席观察到，中国在中间体和溶剂领域显著利用危机实现了自身定位的转变。

背后的逻辑链条值得拆解：当其他亚洲国家遭遇原料供应问题时，中国国内需求依然低迷。过剩的产能迫使产品大量外溢，对全球市场造成了显著扰动。ICIS的数据印证了这一判断：今年3月和4月，中国对亚洲其他地区的化学品出口增长了70\%。

更值得警惕的是，这种出口扩张并非仅仅是量的增长，而是质的迁移。研报指出，中国不仅在基础大宗化学品领域占据主导，而且正在向半精细化学品领域进军。分销商已经开始积极考察中国和亚洲的领先生产商，以应对未来的供应格局。

GS的分析隐含了一个重要判断：中国正在利用这场地缘危机，完成从“大宗出口国”向“精细化工竞争者”的产业升级。对于欧洲精细化工企业而言，这不再是远期的威胁，而是正在发生的竞争态势。

研报中一个容易被忽视但可能最具长期影响力的判断，是关于配方重构（reformulation）的加速。GS明确将其称为“需求的结构性威胁”。

高原料价格迫使运营商和消费者进行配方重构，因为昂贵的投入成本已经无法被市场接受。Equilex主席虽然认为配方重构大体上是暂时性、可逆的措施，但他也承认，部分运营商可能选择永久性改变配方。

IMCD的CEO确认了配方重构需求的显著上升，并将其视为专业分销领域的积极机会。但从行业整体来看，这意味着对某些高成本原料的需求可能永久性流失。

GS在报告中特别强调了一个数据维度：那些拥有高比例客户定制配方销售的公司（如奇华顿，96\%的销售来自客户定制配方，而行业平均为52\%）和高研发投入的公司（如诺维信，研发占销售比例1

[... middle omitted ...]

洲以外。

然而，悖论在于，客户对欧洲生产、欧洲消费的价值认知正在强化。很多产品在欧洲根本买不到，客户反而在主动要求欧洲生产商重启某些化学品的生产。

这揭示了一个两难困境：欧洲既无法维持自给自足的供应能力，又无法完全依赖进口。这种“既不能完全自产，又不愿完全依赖进口”的尴尬位置，恰恰是欧洲化工行业未来几年最核心的战略挑战。

对于投资者而言，这意味着那些拥有多源供应网络、能在全球范围内灵活调配产能的企业，将获得明显的竞争优势。GS在报告中特别提到，IMCD凭借其多源供应网络，在整个冲突期间成功避免了产品短缺，这得益于其供应商（如巴斯夫）在美国和中国都设有工厂，能够重新调配网络。

Integra Petrochemicals的CEO发出了一个严峻的警告：现在可能已经太晚了。欧洲正在关停的工厂很可能不会再重启。

原因有两个层面。其一，一些工厂面临巨大的资本支出压力，因为它们正在接近碳排放上限，合规成本高企。其二，炼油厂的关闭导致化工企业失去了原料供应基础，被迫停产。

\subsection{[R017] 摩根斯坦利：市场低估的不是AI需求，而是供给端结构性约束的定价权}
{\small\color{refgray}归类：半导体设备：DRAM与NAND的结构性分化 / 风险偏好与资金流：供给瓶颈定价权重估}

摩根斯坦利：市场低估的不是AI需求，而是供给端结构性约束的定价权

这份研报的核心判断并不复杂，却容易被当前高涨的情绪掩盖：半导体周期的驱动力正在从需求侧的单边拉动，转向供给侧的结构性约束与客户长期承诺的博弈。摩根斯坦利在最新周报中同时上调了美光（Micron）的盈利预期，并维持对Cerebras Systems的看好，但真正值得关注的不是数字本身，而是支撑这些数字背后的逻辑正在发生质变。

市场习惯于用“AI需求爆发”来解释一切上涨，但这份报告揭示了一个更微妙的层次——当需求已经充分预期，真正决定下一阶段回报率的，是企业能否将规模转化为议价权，以及供给端的技术与资本约束能否被突破。这正是当前半导体投资中容易被忽视的第二层思考。

这份报告发布于6月18日，覆盖美光（MU）和Cerebras Systems（CBRS）两家公司即将发布的财报。摩根斯坦利对两者均给予“超配”（Overweight）评级，美光目标价1050美元，Cerebras目标价未在摘要中明确给出，但分析师强调其“差异化架构”是核心看点。然而，在看似乐观的表象下，报告埋下了几个值得深挖的关键线索。

摩根斯坦利将美光5月季度（F3Q26）的EPS预期从20.57美元上调至21.31美元，8月季度（F4Q26）从25.00美元上调至26.01美元。驱动因素非常直接：DRAM和NAND的ASP涨幅均超出此前预期。报告将DRAM ASP环比涨幅从40\%上调至45\%，NAND从35\%上调至50\%，而第三方预测甚至更高——DRAM约60\%，NAND约75\%。

但真正有洞察力的部分不是这些数字本身，而是分析师如何看待这些数字的可持续性。报告明确指出：“比强劲业绩更重要的是关于当前周期持久性的证据。”这句话点出了市场的核心焦虑——当所有人都知道AI在拉动需求时，周期何时见顶才是真正的分歧点。

摩根斯坦利的回答是：供给端的约束远比市场想象的更持久。报告提到，尽管美光收购力晶（PSMC）的铜锣厂可能加快设备交付，使得全年资本支出从250亿美元上调至270亿美元，但“DRAM整体供应增长仍远低于设备支出所暗示的水平”。原因在于：建设周期、HBM（高带宽内存）的良率折损、以及制程微缩效率下降，这些结构性因素共同限制了供应响应的速度。结果是，比特出货量的季度环比增长仍维持在中等个位数，而需求增速显然更快。

这意味着什么？意味着当前美光的盈利上调不是一次性事件，而是一个持续过程的开始。只要供给约束不被打破，ASP的上升趋势就有望延续到2027年甚至更远。这正是报告将美光目标价设定在1050美元、基于29.5倍周期中段市盈率的逻辑基础——这一倍数高于历史水平，但摩根斯坦利认为这反映了AI带来的新机会，与更广泛的半导体板块估值一致。

2. 战略客户协议（SCA）的透明化程度，决定了美光的估值能否进一步扩张

报告中最具信息量却最容易被忽视的讨论，是关于战略客户协议（SCA）的。一个季度前，美光披露签署了一份为期五年的SCA，但关于条款和财务承诺的规模信息极为有限。摩根斯坦利认为，美光可能在本次财报中宣布更多交易的完成，但不会期望获得更多条款细节——因为美光正在与多个客户谈判，不愿暴露底牌。

这里的关键洞察是：SCA不仅是需求确认的工具，更是估值扩张的催化剂。报告写道：“我们认为这些交易对于市场情绪至关重要，关乎进一步估值扩张的案例。如果新信息有限，股价可能下跌。但即便如此，我们会寻找加仓机会，因为新披露不会改变我们所知道的事实——客户认为DRAM在多年时间范围内日益紧张——这一点在当前9.3倍2027财年EPS的估值中并未被定价。”

这句话值得反复读。摩根斯坦利实际上在说：市场对美光的定价仍然基于历史周期框架，而没有充分计入“客户主动锁定长期供应”这一结构性变化

[... middle omitted ...]

这种影响预计持续到2027年中。

乍看之下，这份财务预测并不令人兴奋——收入持平、毛利率大幅下滑。但摩根斯坦利的结论却是“维持建设性观点”。为什么？因为报告认为，对于Cerebras而言，短期内的财务波动已被市场充分预期，真正的价值在于其差异化架构能否在AI基础设施这一增长最快的细分市场中确立领导地位。

关键变量是250MW产能的部署节奏。报告指出，“产能部署速度是收入和毛利率上行的关键驱动力”，并预期首批250MW产能将在2027年中上线。这意味着，未来几个季度的财报电话会议上，管理层关于产能建设进度的任何更新，都可能比实际财务数字更能影响股价。

这里存在一个尚未完全回答的问题：Wafer Scale Engine（晶圆级引擎）能否证明其作为持久竞争优势的价值？Cerebras需要向客户证明其技术能带来显著价值，并且能够在经济上实现规模化。目前，摩根斯坦利认为投资辩论的核心是执行能力，但长期来看，技术壁垒的可持续性才是决定其能否从“差异化”走向“主导”的关键。

\subsection{[R018] NOM：人民币中间价模型正在传递一个比市场预期更克制的信号}
{\small\color{refgray}归类：宏观与利率：中国内需磨底，家庭资产负债表再平衡}

这份NOM亚洲外汇策略团队的最新报告，看似只更新了一个每日模型预测数字，但拆解其内部结构和参数调整后，一个清晰的信号浮现出来：当前人民币定价机制中的“逆周期因子”正在发挥比表面数字更重要的校准作用，市场对汇率单向波动的定价可能过于激进。

报告的核心判断并不复杂：NOM模型给出的美元兑人民币中间价预测为6.7787，较前次预测下移了343个基点。但当我们把目光从绝对数字移开，聚焦于模型内部的结构性调整时，会发现一个更具深意的变化——计入逆周期因子后的预测值为6.8008，与前一交易日官方收盘价的偏离仅为122个基点。这个差值，才是真正值得关注的信号。

市场往往习惯于将注意力集中在人民币汇率的绝对水平上，讨论“破7”还是“守7”，讨论升值还是贬值的方向。但这份NOM报告提示了一个更精细的观察维度：政策制定者正在通过逆周期因子，对模型输出的“纯粹”市场定价进行微调，而这种微调的幅度和频率，才是理解当前汇率政策意图的真正钥匙。

1. 模型预测的下移并非线性趋势，而是多重短期因素的一次性校准

NOM模型预测从6.8130下移至6.7787，表面上看是343个基点的显著变化。但报告中的“无逆周期因子模型预测”和“含逆周期因子模型预测”两个数值的差异，揭示了这并非一个简单的方向性信号。

从报告提供的图表看，模型误差和隔夜贡献因素的变化，是推动预测下移的主要驱动力。这些因素包括隔夜美元指数的波动、其他亚洲货币的交叉汇率变动、以及市场情绪的短期扰动。它们具有明显的时效性，而非结构性趋势的改变。

一个关键点在于：模型预测的下移幅度，与官方中间价的实际调整幅度之间存在差距。这种差距本身就是一种政策信号。当市场参与者单纯追逐模型预测的方向时，他们可能会忽略逆周期因子正在发挥的“缓冲”作用。

这意味着，任何基于单一模型预测做出的汇率方向性判断，都可能高估了短期波动的持续性。真正的变量不在于6.7787还是6.8008，而在于政策制定者愿意容忍多大的偏离。

报告中最值得反复推敲的数据点，是“含逆周期因子模型预测”6.8008与前一交易日官方收盘价的差值——122个基点。这个数值代表了政策制定者通过逆周期因子对市场定价进行的主动干预幅度。

逆周期因子的本质，是央行在人民币中间价形成机制中引入的一个“调节阀”。当市场情绪过度乐观或悲观时，这个阀门可以被拧紧或放松，以平滑汇率的短期过度波动。122个基点的偏离，在历史数据中属于中等偏上的水平，表明政策制定者认为当前市场定价存在一定程度的偏离，需要进行温和校正。

但更重要的是，这个偏离值本身也在变化。报告显示，此前的模型预测与官方收盘价的偏离为164个基点，而当前已收窄至122个基点。这种收窄趋势，可能意味着两件事中的一件：要么市场定价正在向政策合意区间靠拢，要么政策制定者正在逐步减少干预力度。无论是哪一种，都指向一个结论：当前汇率定价的“政策锚”比市场想象的更稳定，而非更脆弱。

对于关注人民币资产定价的投资者而言，这意味着不应将短期汇率波动解读为政策转向的信号。逆周期因子的存在，恰恰是为了防止这种误读。

3. 报告未完全展开的关键问题：逆周期因子的校准是否存在一个“隐含目标区间”

这份NOM报告的核心价值在于提供了精确的模型预测和校准数据，但它也留下了一个未完全回答的问题：逆周期因子的校准是否存在一个隐含的目标区间？或者更直白地说，政策制定者心中是否有一个“合意汇率水平”？

从报告提供的图表来看，模型误差和每日中间价变动呈现出明显的波动性。逆周期因子的调整似乎是对这些短期波动的反应，而非对某个长期目标的追逐。但这恰恰是问题所在：如果没有一个隐含的目标区间，那么逆周期因子的校准就只是“头痛医头”的被动应对；如果存在目标区间，那么当前的校准幅度就是政策意图的直接体现。



[... middle omitted ...]

结论。

4. 2026年下半年的事件日历，为汇率定价提供了可预期的压力测试场景

报告附录中列出的2026年下半年事件日历，虽然看似只是常规的日程安排，但将其与汇率模型结合起来看，就构成了一个清晰的“压力测试”框架。

7月底的政治局会议、11月的APEC峰会、年底前可能发生的领导人会晤，以及12月的中央经济工作会议——这些事件每一个都可能成为影响汇率预期的关键节点。尤其是APEC会议在中国深圳举办，以及年底前可能的中美领导人会晤，都意味着汇率稳定在特定时间段内具有更高的政策优先级。

这意味着，逆周期因子的校准幅度可能会在这些关键事件前后发生可预期的变化。在市场关注度最高的时间窗口，政策制定者倾向于维持汇率稳定；而在事件间隔期，模型预测的纯粹市场定价可能会获得更大的自由度。

对于交易者和投资者而言，这意味着不应将当前的逆周期因子校准幅度视为常态。它可能在未来几个月内根据事件日程发生显著调整。理解这种“事件驱动型”的汇率管理逻辑，比猜测一个具体的汇率点位更有价值。

\subsection{[R019] Bernstein：日本零售业真正被低估的不是消费复苏，而是供给侧的整合拐点}
{\small\color{refgray}归类：地产与消费：中国内需磨底，日本零售整合拐点}

Bernstein：日本零售业真正被低估的不是消费复苏，而是供给侧的整合拐点

全球投资者长期将日本国内零售视为一个“已被充分讨论”的市场——人口萎缩、增长有限、整合停滞。这种看法本身并没有错，但它遗漏了一个关键变量：过去三十年“失去的岁月”同时也是一个漫长的通缩过滤器，弱者被淘汰，幸存者的运营效率在压力下持续进化。Bernstein最新发布的日本零售系列报告提出了一个与市场共识截然不同的判断：日本国内零售业正在接近一个加速整合的拐点，而这一变化的驱动力已经从周期性转向结构性。

这份报告的核心洞见在于，日本零售业当前面临的不是需求端的短期波动，而是供给端成本结构的根本性重塑。超市行业约1.5\%的经营利润率，对应着约14\%的人工成本率——这个方程式在最低工资自2015年以来累计上涨超过30\%、远超CPI约10\%涨幅的背景下，已经变得不可持续。支撑行业碎片化的传统支柱——店内生鲜加工、地理分散化、对非正规劳动力的依赖——正在同时瓦解。规模较小的企业，缺乏投资自动化和中央加工的能力，正在被系统性边缘化。

这与西方零售业完成整合的历史路径不同。西欧和美国在几十年前就完成了这一过程，而日本现在才进入“中场局”。对于产业决策者和投资者而言，这意味着一个清晰的、未来十年的集中度提升路径正在形成，而非市场普遍预期的“一成不变”。

日本现代消费者呈现出一种鲜明的二元模式：工作日的“效用型”购物——快速、功能性支出集中在药妆店——与周末的“寻宝型”消费在折扣零售商处完成。这种分裂几乎没有给中间路线业态留下生存空间。

综合超市（GMS）曾经是“一站式解决方案”，但现在正被专业零售商和更聚焦的业态持续蚕食品类份额。结果是，GMS逐渐变成了一个难以定义自身价值主张的业态——在运营上仍然存在，但在战略上已经过时。Bernstein的报告用了一个精准的比喻：GMS正在变成“僵尸”——运营上存在，但战略上已经死亡。

这个观察的意义在于，它解释了为什么日本零售业的整合不会是均匀的。整合将首先发生在那些“夹心层”业态，而折扣店和药妆店将受益于消费者行为的结构性转移。对于持有GMS相关资产的企业和投资者来说，这不仅仅是周期性问题，而是需要重新评估资产长期价值的信号。

这是整份报告最关键的论据之一。日本超市行业1.5\%的OPM与14\%的人工成本率之间的差距，在过去十年中因最低工资大幅上涨而急剧扩大。传统的应对方式——依赖非正规劳动力、分散化运营、店内现做——正在失去效力。

Bernstein的分析指出，规模较小的企业无法负担自动化投资和中央加工设施，这意味着它们的成本劣势正在从“可管理”变成“致命”。这不是一个短期劳动力市场紧张的问题，而是日本人口结构、劳动力供给和最低工资政策共同作用下的长期趋势。

对于产业决策者而言，这个判断的蕴含是：整合的驱动力不是来自管理层的战略选择，而是来自成本结构的物理约束。那些能够通过规模效应实现自动化、中央化处理的企业，将获得不可逆的成本优势。那些不能的企业，无论管理层多么优秀，都将被结构性地挤出市场。

在整体零售利润微薄的背景下，Bernstein对PPIH（唐吉诃德母公司）给出了跑赢大盘的评级。这个判断背后的逻辑值得仔细拆解。

PPIH属于折扣零售中的特殊子类别。报告描述其运营模式为“高毛利率、高成本、高利润”——约7\%的经营利润率在行业中极为突出。这听起来有些反直觉：折扣店如何能实现高毛利？关键在于其独特的商品组合和供应链策略。免税销售和自有品牌的增长提供了额外的结构性增长空间，而竞争对手难以复制其复杂的采购和定价能力。

这个案例的价值在于，它展示了在整合浪潮中，真正的赢家可能不是最大规模的参与者，而是拥有独特运营模型的企业。对于投资者来说，这意味着需要超越传统的“规模效应”框架，去理解哪些

[... middle omitted ...]

这个案例提供了另一个维度的启示：在整合浪潮中，作为收购方，估值溢价是否合理，取决于整合执行的能力，而不仅仅是交易规模。

Seven \& i Holdings（7-Eleven母公司）的情况更为复杂。Bernstein给出了中性评级，理由是其核心便利店业务已过巅峰期，加盟商经济性正在恶化，治理压力也在上升。但同时，股票仍然保留了显著的“企业价值重估期权”——来自Couche-Tard的收购底价，以及伊藤忠与创始家族主导的便利店重组可能性。

这个判断的核心是：Seven \& i的价值不再完全取决于其运营表现，而是越来越多地取决于公司结构和所有权变化带来的价值释放。对于投资者来说，这意味着需要同时评估运营基本面和企业治理动态——两者之间的脱节可能正是机会所在。

对于产业决策者而言，这个案例提出了一个更广泛的问题：在日本零售业整合加速的背景下，哪些企业可能成为整合的目标，而不是整合的推动者？那些业务结构复杂、运营效率下降、但资产价值被低估的企业，可能成为资本运作的标的。

\subsection{[R020] GS：软件行业的真正分水岭不是AI本身，而是“好粘性”与“坏粘性”的重新定价}
{\small\color{refgray}归类：AI/算力：云巨头资本效率竞赛与内存成本冲击}

GS：软件行业的真正分水岭不是AI本身，而是“好粘性”与“坏粘性”的重新定价

Databricks年度大会结束后，GS软件团队发布了一份不同寻常的报告。它没有罗列产品更新清单，而是提出了一个判断：软件行业的估值分化将加速，而分化的核心变量不是AI能力的有无，而是“粘性”的质量。

这份报告的洞察起点是一个简单却尖锐的问题：当AI让数据迁移成本骤降，哪些软件公司的护城河会变浅？哪些反而会更深？GS分析师给出的框架是“好粘性”与“坏粘性”之分。前者靠持续创新和客户热爱来锁定用户，后者依赖历史遗留的迁移成本。AI正在系统性地瓦解后者。

这意味着，未来12-18个月，软件投资的核心命题不再是“谁有AI故事”，而是“谁的商业模式经得起重新定价”。

1. 迁移成本正在被AI工具系统性地压低，依赖“坏粘性”的公司面临估值重塑

GS的报告明确点出了一个被市场低估的结构性变化：AI不仅改变了软件的功能边界，更改变了竞争的基础设施。过去，一家SaaS公司只要把客户数据锁在自己的格式里，就能获得相当可观的定价权。客户即使不满意，也会因为迁移成本过高而选择留下。这种“坏粘性”曾经是很多软件公司利润率的基石。

但Databricks在大会上的产品发布，尤其是LTAP（Lake Transactional/Analytical Processing）和Open Sharing协议，清晰地展示了新范式：数据可以在开放格式下自由流动，分析引擎和事务引擎可以共享同一个数据源。当客户的数据不再被某个特定平台绑架，迁移成本就大幅下降。

GS分析师特别指出，Databricks的定价策略本身就反映了这种逻辑。公司内部明确表示，它可以选择定价高出两倍，但短期定价优化往往会牺牲长期竞争力和客户终身价值。这种“以开放换增长”的策略，实际上是在提高整个行业的竞争标准。

对于投资者而言，这意味着需要重新审视那些过去依赖高迁移成本来维持续费率的公司。当客户可以更容易地用AI工具完成数据迁移，那些产品创新速度放缓、客户满意度下降的公司，其估值中的“粘性溢价”需要被重新定价。

2. 软件行业正在从“打包应用”走向“定制代理”，中间层成为新的价值高地

GS报告中最具前瞻性的洞察之一，是对软件架构演进的判断。传统的SaaS应用是孤立的，每个应用管理自己的一套数据和流程。但真正有价值的AI代理应用，恰恰需要跨越这些孤岛，存在于SaaS系统之间的“空白地带”。

这意味着，未来软件行业的TAM（可寻址市场）将发生结构性变化。GS此前已经量化过，代理式AI将在2030年前为软件行业带来20\%的TAM提升，相当于行业增速从7\%提升到9\%。但这份报告进一步指出，增量中的大部分将流向“定制应用”而非“打包应用”。

支撑定制应用的核心是“操作系统层”——一个能够从各种记录系统和企业数据存储中拉取数据的智能中间层。GS识别了四种构建这一层的路径：Databricks的Ontology产品、Palantir的工程外包模式、云厂商的生态系统（如微软的“前沿智能”愿景），以及ServiceNow等AI点产品的控制塔。

这里的关键判断是：组织复杂度决定了选择。复杂的企业组织更可能选择定制路径，而简单的SMB可能继续使用现成的SaaS。这意味着，能够帮助客户构建和维护这个中间层的公司——无论是通过平台、工程服务还是治理工具——将获得更大的价值分配。

3. Databricks正在用“数据库革命”来支撑代理时代，而市场可能低估了LTAP的长期影响

GS报告用大量篇幅分析了Databricks的产品创新，但最核心的观点是：当整个行业都在追逐代理应用时，Databricks选择回头解决数据库基础设施层面的问题。这看起来是“向后看”，实际上是“向前看”的必要条件。

LTAP的

[... middle omitted ...]

定性高且可分支的环境。Lakebase引入了近乎实时的数据库分支能力，让代理可以安全地分叉数据集、实验、写入更改并在需要时回退。这不是一个渐进式改进，而是为代理工作负载重新设计数据库基础设施。

GS的数据显示，Databricks的核心收入增长（排除代币转售）已经连续五个季度加速。这并非偶然，而是产品市场匹配度提升的直接结果。

4. 定制代理应用的维护成本正在被系统性降低，但报告尚未完全回答“谁将获得最大份额”

GS报告提到了一个被很多人忽视的关键点：维护成本是定制应用最大的隐性障碍。ServiceNow此前曾量化，定制应用的总拥有成本是SaaS的5-10倍。这解释了为什么过去大多数企业选择打包SaaS而非自建。

但Databricks的多项创新直接瞄准了这一痛点。Unity Catalog作为通用治理层嵌入Ontology；Ontology利用单一数据源；LTAP消除了不同数据源之间脆弱的管道和ETL流程。这些创新的综合效果，是大幅降低定制应用的维护TCO。

\section{Disclaimer}
{\scriptsize\color{refgray}
This community is a paid learning and information-sharing community created for general educational purposes only. The membership fee is charged solely for access to community discussions, educational content, organization, and operational costs. It is not an advisory fee, management fee, performance fee, brokerage fee, or compensation for personalized financial, investment, legal, tax, or professional advice.\par
I participate and share information solely as an individual and for educational discussion among community members. I am not acting as a financial adviser, investment adviser, broker, dealer, portfolio manager, legal adviser, tax adviser, accountant, fiduciary, or any other licensed professional adviser. Nothing shared in this community constitutes, and should not be interpreted as, financial advice, investment advice, legal advice, tax advice, a recommendation, solicitation, offer, endorsement, instruction, or guarantee to buy, sell, hold, short, trade, or invest in any security, cryptocurrency, fund, derivative, strategy, or other financial product.\par
All content shared in this community, including messages, discussions, links, files, charts, examples, market commentary, watchlists, AI-generated or AI-polished summaries, and third-party materials, is general, impersonal, and educational in nature. It does not take into account any individual's financial situation, investment objectives, risk tolerance, investment horizon, liquidity needs, tax status, legal restrictions, or suitability requirements. No content should be relied upon as suitable for any specific person, account, portfolio, or financial decision.\par
Information shared in this community may be incomplete, inaccurate, outdated, speculative, AI-generated, AI-edited, or based on third-party internet sources that have not been independently verified. I make no representation or warranty as to the accuracy, completeness, timeliness, reliability, or suitability of any information shared.\par
Financial markets involve substantial risk, including the possible loss of principal. Past performance, historical data, hypothetical examples, simulated results, backtests, personal opinions, or educational examples do not guarantee future results. Each member is solely responsible for conducting their own research, exercising independent judgment, and making their own decisions. Before making any financial, investment, legal, or tax decision, members should consult qualified and licensed professionals.\par
Participation in this community, payment of any membership fee, or communication with me or other members does not create any adviser-client, fiduciary, professional, contractual, agency, partnership, employment, or other advisory relationship. By joining or participating in this community, you acknowledge and agree that you are solely responsible for your own decisions, actions, transactions, profits, losses, risks, and consequences, and that I shall not be liable for any direct, indirect, incidental, consequential, financial, legal, tax, or other loss or damage arising from or related to any content, discussion, information, or material shared in this community.\par
}
\end{document}
